% =====================================================================
% MAS-Hunt — Dissertação de Mestrado (ABNT, PT-BR)
% Documento mestre. Espelha a configuração funcional de ../pre-projeto.tex
% (abntex2 + abntex2cite + capa.sty), expandida para o formato de
% dissertação (capítulos via \include, resumo/abstract, sumário).
% Compila mesmo com capítulos em placeholder.
% =====================================================================
\documentclass[
  12pt,
  a4paper,
  oneside,        % impressão em face única
  brazil
]{abntex2}

%%%%%%%%%%%%%%%%%%%%%%%%%%%% Pacotes %%%%%%%%%%%%%%%%%%%%%%%%%%%%%%%%%%%%
\usepackage{ifthen}
\usepackage[T1]{fontenc}          % Codificação da fonte
\usepackage[utf8]{inputenc}       % Codificação do arquivo
\usepackage{lmodern}              % Fonte Latin Modern
\usepackage[brazil]{babel}        % Idioma do documento
\usepackage{graphicx}             % Inclusão de gráficos
\usepackage[table]{xcolor}        % Cores para tabelas
\usepackage{float}                % Posicionamento de figuras e tabelas
\usepackage{booktabs}             % Tabelas de qualidade tipográfica
\usepackage{pifont}               % Símbolos \ding (tabelas comparativas)
\usepackage{amsmath,amssymb}      % Matemática
\usepackage[num]{abntex2cite}     % Citações ABNT (numérico, como no pré-projeto)
\citebrackets[]
\usepackage{indentfirst}          % Indenta o primeiro parágrafo de cada seção
\usepackage{microtype}            % Melhoria da tipografia
\usepackage{hyperref}             % Hiperlinks
\usepackage{setspace}             % Controle de espaçamento

% Imagens referenciadas a partir do diretório-raiz do repositório
\graphicspath{{../}{../mashunt/}{./}}

% Estilo de capa customizado (versão local que aponta para ../bandeira.png)
\usepackage{capa}

%%%%%%%%%%%%%%%%%%%%%%% Configuração dos hiperlinks %%%%%%%%%%%%%%%%%%%%%
\hypersetup{
    colorlinks=true,
    linkcolor=blue,
    citecolor=blue,
    filecolor=magenta,
    urlcolor=blue,
    pdftitle={MAS-Hunt: Arquitetura de Governanca Multiagente Resiliente para Caca a Ameacas},
    pdfauthor={Paulo Matheus Nicolau Silva}
}

%%%%%%%%%%%%%%%%%%%%%%%%%%%% Dados da Capa %%%%%%%%%%%%%%%%%%%%%%%%%%%%%%
\titulo{MAS-Hunt: Arquitetura de Governança Multiagente Resiliente para a Busca Proativa por Ameaças Cibernéticas em Ambientes Windows}
\autor{Paulo Matheus Nicolau Silva}
\local{Brasília}
\data{2025}

%%%%%%%%%%%%%%%%%%%%%%%%%%%%%%%%%%%%%%%%%%%%%%%%%%%%%%%%%%%%%%%%%%%%%%%%%
\begin{document}
\renewcommand{\baselinestretch}{1.2}\normalsize

% ------------------------------------------------------------------ %
% ELEMENTOS PRÉ-TEXTUAIS
% ------------------------------------------------------------------ %
\pretextual

\imprimircapa

% --- Folha de rosto ---
\begin{folhaderosto*}
  \begin{center}
    {\ABNTEXchapterfont\large Universidade de Brasília — UnB}\\
    {\ABNTEXchapterfont Faculdade de Tecnologia — FT}\\
    {\ABNTEXchapterfont Departamento de Engenharia Elétrica — ENE}\\
    {\ABNTEXchapterfont Programa de Pós-Graduação em Engenharia Elétrica — PPEE}\\[3cm]

    {\ABNTEXchapterfont\large\imprimirautor}\\[3cm]

    {\ABNTEXchapterfont\bfseries\Large\imprimirtitulo}\\[3cm]
  \end{center}

  \hspace{.45\textwidth}
  \begin{minipage}{.5\textwidth}
    \SingleSpacing
    Dissertação de mestrado apresentada ao Programa de
    Pós-Graduação em Engenharia Elétrica da Universidade de
    Brasília, como requisito parcial à obtenção do título de
    Mestre em Engenharia Elétrica, na área de concentração
    Segurança Cibernética.\\[0.4cm]
    \textbf{Área de concentração:} Segurança Cibernética.\\
    \textbf{Linha de pesquisa:} Segurança e Inteligência Cibernética.\\[0.4cm]
    Orientador(a): a definir.
  \end{minipage}
  \vfill
  \begin{center}
    \imprimirlocal\\
    \imprimirdata
  \end{center}
\end{folhaderosto*}

% --- Resumo (PT-BR) ---
\begin{resumo}
A crescente sofisticação e a natureza evasiva das ameaças cibernéticas
contemporâneas — marcadas pelo uso de técnicas de \emph{living-off-the-land}
(LOLBins), pela permanência prolongada de atacantes (\emph{dwell time}) e pela
sobrecarga cognitiva imposta aos analistas de segurança — expõem os limites das
arquiteturas reativas e baseadas em assinaturas. Agentes autônomos apoiados em
grandes modelos de linguagem (LLMs) surgem como alternativa proativa, porém são,
eles próprios, vulneráveis a manipulação adversarial, incluindo envenenamento de
memória, injeção de \emph{prompt} e amplificação de mau funcionamento. Esta
dissertação apresenta o MAS-Hunt (\emph{Multi-Agent System for Threat Hunting}),
uma arquitetura multiagente em três camadas — Conselho de Governança, Gerentes e
Trabalhadores/Caçadores — concebida com prioridade à segurança e materializada
como um \emph{plugin} do Claude Code que orquestra subagentes por meio de
habilidades, comandos de barra e \emph{hooks}. O sistema incorpora cinco
mecanismos de governança: integridade de memória (M1), validação cruzada (M2),
resistência a injeção via \emph{framework} estruturado de análise (M3),
monitoramento comportamental (M4) e quarentena (M5). A avaliação é conduzida em
um laboratório instrumentado com Sysmon e Elastic Security sobre um domínio
Windows, comparando o MAS-Hunt completo (C1-full) a três condições de referência
(C1-nohardening, C2-naive e C3-Elastic) quanto a desempenho de detecção,
resiliência adversarial e taxa de falsos positivos. Na matriz v2 com 1.092
janelas adversariais ricas em sinal, todas as três condições baseadas em modelo
de linguagem superaram a linha de base SIEM C3-Elastic (F2~$=0{,}150$) por
aproximadamente 4,5 vezes, ganho dirigido pela revocação. A governança
multiagente, porém, não melhorou o F2 de detecção sobre o modelo de passagem
única: C1-full (0,675) e C2-naive (0,672) mostraram-se estatisticamente
indistinguíveis, com intervalos de confiança sobrepostos. O valor distintivo da
governança manifestou-se na resistência à injeção adversarial: somente C1-full
detectou e sinalizou conteúdo adversarial injetado, marcando 407 janelas
contaminadas contra zero nas demais condições. Conclui-se que o analista agêntico
supera dramaticamente as regras estáticas de SIEM, e que a governança multiagente
agrega resistência à injeção e auditabilidade --- e não maior acurácia de
detecção --- sobre o modelo de linguagem isolado.

\noindent\textbf{Palavras-chave}: caça a ameaças; sistemas multiagente; segurança
de agentes de IA; injeção adversarial; governança; Elastic Stack; Windows.
\end{resumo}

% --- Abstract (EN) ---
\begin{resumo}[Abstract]
\begin{otherlanguage*}{english}
The escalating sophistication and evasive nature of contemporary cyber threats —
characterized by living-off-the-land (LOLBin) techniques, prolonged attacker
dwell time, and the cognitive overload imposed on security analysts — expose the
limits of reactive, signature-based architectures. Autonomous agents powered by
large language models (LLMs) offer a proactive alternative, yet they are
themselves vulnerable to adversarial manipulation, including memory poisoning,
prompt injection, and malfunction amplification. This dissertation presents
MAS-Hunt (Multi-Agent System for Threat Hunting), a security-first three-layer
multi-agent architecture — Governance Board, Managers, and Workers/Hunters —
implemented as a Claude Code plugin that orchestrates subagents through skills,
slash commands, and hooks. The system embeds five governance mechanisms: memory
integrity (M1), cross-validation (M2), injection resistance via a structured
analysis framework (M3), behavioral monitoring (M4), and quarantine (M5). The
evaluation is conducted in a laboratory instrumented with Sysmon and Elastic
Security over a Windows domain, comparing the full MAS-Hunt (C1-full) against
three baselines (C1-nohardening, C2-naive, and C3-Elastic) with respect to
detection performance, adversarial resilience, and false-positive rate. In the
v2 matrix with 1,092 signal-rich adversarial windows, all three language-model
conditions outperformed the C3-Elastic SIEM baseline (F2~$=0.150$) by roughly
4.5 times, a gain driven by recall. Multi-agent governance, however, did not
improve detection F2 over the single-pass model: C1-full (0.675) and C2-naive
(0.672) were statistically indistinguishable, with overlapping confidence
intervals. The distinctive value of governance emerged in adversarial-injection
resistance: only C1-full detected and flagged injected adversarial content,
marking 407 contaminated windows versus zero for the other conditions. We
conclude that the agentic analyst dramatically outperforms static SIEM rules,
and that multi-agent governance adds injection resistance and auditability ---
rather than higher detection accuracy --- over the standalone language model.

\noindent\textbf{Keywords}: threat hunting; multi-agent systems; AI agent
security; adversarial injection; governance; Elastic Stack; Windows.
\end{otherlanguage*}
\end{resumo}

% --- Sumário ---
\pdfbookmark[0]{\contentsname}{toc}
\tableofcontents*

% ------------------------------------------------------------------ %
% ELEMENTOS TEXTUAIS
% ------------------------------------------------------------------ %
\textual

\chapter{Introdução}
\label{cap:introducao}

\section{Contextualização}
\label{sec:contexto}

O cenário contemporâneo da segurança cibernética é definido pela velocidade e
pela sofisticação crescentes das operações adversárias. Ataques de
\emph{ransomware}, campanhas de \emph{phishing} e a exploração de
vulnerabilidades tornaram-se mais frequentes e complexos, obrigando as
organizações a repensar continuamente suas estratégias de defesa
\cite{academia20}. A proliferação de fornecedores comerciais de vigilância e de
mercados clandestinos de credenciais válidas amplia o acesso a capacidades
ofensivas avançadas, de modo que atacantes alcançam, com regularidade, os
ambientes corporativos antes que qualquer alerta seja gerado.

Nesse contexto, a busca proativa por ameaças, ou \emph{threat hunting},
consolidou-se como uma prática que transcende a reatividade dos sistemas
tradicionais. Em vez de aguardar que um alerta seja disparado, os profissionais
de segurança adotam técnicas de investigação forense para examinar redes,
sistemas e estações de trabalho à procura de sinais de comprometimento
\cite{academia22}. Essa abordagem ativa permite identificar e neutralizar
indicadores de atividade maliciosa antes que sejam plenamente explorados pelos
atacantes, reduzindo a janela de oportunidade do adversário. A literatura
descreve essa inversão de postura como a transformação do alvo em caçador, com
ênfase no ciclo de vida da caça e no papel emergente da inteligência artificial
\cite{hillier2022turninghuntedhunterthreat}.

Um traço distintivo das ameaças modernas é a adoção sistemática de técnicas de
\emph{living-off-the-land} (LOLBins), nas quais o atacante abusa de binários,
\emph{scripts} e utilitários legítimos já presentes no sistema operacional — a
exemplo de \texttt{PowerShell}, \texttt{wmic}, \texttt{certutil},
\texttt{rundll32} e \texttt{mshta} — para executar suas ações. Por se valerem de
ferramentas confiáveis e assinadas pelo próprio fornecedor do sistema, essas
técnicas produzem telemetria que, isoladamente, é praticamente
indistinguível da atividade administrativa benigna. O resultado é uma evasão
eficaz das defesas baseadas em assinaturas e listas de bloqueio, deslocando o
ônus da detecção para a análise comportamental e contextual dos eventos.
Ambientes Windows, predominantes em redes corporativas e sustentados por uma
rica superfície de utilitários nativos, são especialmente suscetíveis a esse
padrão de abuso.

A esse desafio técnico soma-se um desafio humano e operacional. As soluções de
detecção contemporâneas, mesmo quando dotadas de modelos estatísticos acurados,
estão sujeitas à falácia da taxa-base: em ambientes nos quais a atividade
maliciosa é rara, mesmo classificadores de elevada acurácia geram um volume
desproporcional de falsos positivos, tornando-se operacionalmente ineficazes. O
analista de segurança, por sua vez, é submetido a uma sobrecarga cognitiva
permanente, dividido entre o triagem de alertas de baixa qualidade e a
investigação aprofundada de hipóteses relevantes. Essa sobrecarga compromete a
capacidade de resposta da organização e favorece a fadiga de alertas, na qual
eventos genuinamente críticos podem ser negligenciados em meio ao ruído.

Os agentes de inteligência artificial apresentam-se como uma resposta promissora
a esse cenário. Tais agentes são estruturas de \emph{software} que ampliam as
capacidades dos grandes modelos de linguagem (LLMs — \emph{Large Language
Models}), permitindo que estes raciocinem sobre evidências, interajam com
ferramentas externas e executem ações no sistema \cite{academia20}. O mecanismo
de atenção introduzido por Vaswani et al. \cite{vaswani2017attention} pavimentou
o caminho para esses modelos, e trabalhos subsequentes demonstraram seu potencial
de coordenação e tomada de decisão conjunta em arranjos multiagente
\cite{agashe2023llmcoordination}. Na segurança cibernética, propostas como o
HuntGPT \cite{ali2023huntgptintegratingmachinelearningbased} e modelos de geração
de hipóteses para a caça a ameaças \cite{10713082} ilustram a viabilidade de
aplicar LLMs à detecção proativa, conjugando explicabilidade e raciocínio
investigativo.

Essa autonomia, contudo, é uma faca de dois gumes. A mesma flexibilidade que
torna os agentes poderosos os expõe a uma classe de ataques que tem por alvo o
próprio agente, e não apenas o sistema defendido. Pesquisas recentes documentam o
envenenamento de memória e de bases de recuperação aumentada (RAG), a injeção
indireta de \emph{prompt} dissimulada em conteúdo externo e a amplificação de mau
funcionamento, na qual um erro inicial desencadeia um ciclo de ações inúteis que
degrada o desempenho do agente. Demonstrou-se, ainda, a possibilidade de
\emph{jailbreaks} infecciosos, em que uma única entrada adversarial se propaga
por uma população de agentes. Diante disso, qualquer arquitetura defensiva que
não modele um adversário igualmente capaz de atacar a camada de raciocínio é,
por construção, incompleta.

A governança em segurança cibernética emerge, portanto, como pilar para alinhar
as estratégias de defesa às necessidades do negócio e para coordenar, de forma
segura, o trabalho dos agentes \cite{academia21}. Frameworks consolidados — como
COBIT, NIST CSF, ISO 27001 e o recente ISO 42001 — precisam evoluir para
incorporar os riscos específicos introduzidos pela integração de tecnologias de
IA \cite{McIntosh2024,isaca2025leveraging}. A presente dissertação situa-se
nessa interseção: propõe a governança multiagente não como camada burocrática
acessória, mas como mecanismo técnico de resiliência, capaz de conter a
propagação de comprometimentos e de validar achados antes que se convertam em
decisões.

\section{Definição do Problema}
\label{sec:problema}

O problema de pesquisa pode ser enunciado pela tensão entre duas exigências
aparentemente conflitantes. De um lado, demanda-se um sistema de caça a ameaças
\emph{eficaz}, capaz de operar sobre a telemetria viva de um ambiente corporativo
Windows, detectar técnicas evasivas de \emph{living-off-the-land} e reduzir o
volume de falsos positivos que sobrecarrega os analistas. De outro, exige-se que
esse mesmo sistema seja \emph{resiliente}: como os agentes que o compõem são
suscetíveis a envenenamento de memória, injeção de \emph{prompt} e amplificação
de mau funcionamento, a automação introduz uma nova superfície de ataque que, se
explorada, pode subverter justamente a função de defesa que o sistema deveria
exercer.

Formula-se, assim, a questão central de pesquisa: \emph{como projetar uma
arquitetura multiagente para caça a ameaças que opere de forma eficaz sobre a
telemetria corporativa real e, simultaneamente, demonstre resiliência diante de
ataques adversariais dirigidos aos próprios agentes?} A resposta a essa questão
não se esgota na escolha de um modelo de linguagem mais capaz; ela depende,
sobretudo, da arquitetura de coordenação e dos mecanismos de governança que
regulam o fluxo de informação e a tomada de decisão entre os agentes.

\section{Lacuna de Pesquisa}
\label{sec:lacuna}

A revisão da literatura revela que, embora os agentes de IA tenham avançado de
tarefas supervisionadas e específicas para sistemas de raciocínio autônomo, as
arquiteturas atuais apresentam lacunas críticas quando confrontadas com um
adversário inteligente. Identificam-se, em particular, três lacunas que esta
dissertação se propõe a endereçar:

\begin{enumerate}
  \item \textbf{Ausência de orquestração multiagente robusta.} Faltam
  \emph{frameworks} que impeçam que o comprometimento de um único agente se
  propague sistemicamente, contendo a falha em vez de amplificá-la. A maioria das
  propostas concentra-se no desempenho de detecção e negligencia a contenção de
  comprometimentos internos.

  \item \textbf{Ausência de validação de integridade em tempo de execução.} Os
  mecanismos de memória e as bases de conhecimento dos agentes carecem de
  validação contínua contra envenenamento furtivo de dados, deixando o
  raciocínio do agente vulnerável a manipulações que persistem entre interações.

  \item \textbf{Autocorreção ineficaz frente à degradação comportamental.} Os
  agentes não dispõem de mecanismos confiáveis para detectar e reverter sua
  própria degradação, permanecendo expostos à amplificação de mau funcionamento e
  a ataques de negação de serviço sobre o ciclo de raciocínio.
\end{enumerate}

Soma-se a essas lacunas a escassez de avaliações conduzidas em condições
realistas — sobre telemetria viva, com instrumentação de produção e contra
adversários que atacam a camada de agentes — e a ausência de uma comparação
sistemática entre arquiteturas governadas e linhas de base sem governança. É
precisamente nessa interseção entre resiliência adversarial, validação de
integridade e operação realista que o MAS-Hunt se posiciona.

\section{Objetivos}
\label{sec:objetivos}

\subsection{Objetivo Geral}
\label{subsec:objetivo-geral}

Projetar, implementar e avaliar o MAS-Hunt, uma arquitetura de governança
multiagente para a busca proativa por ameaças cibernéticas em ambientes Windows,
capaz de operar sobre telemetria viva e de demonstrar resiliência mensurável a
ataques adversariais dirigidos aos agentes, em comparação a linhas de base sem
governança.

\subsection{Objetivos Específicos}
\label{subsec:objetivos-especificos}

Para alcançar o objetivo geral, estabelecem-se os seguintes objetivos
específicos:

\begin{enumerate}
  \item Revisar criticamente a literatura sobre caça a ameaças e detecção
  comportamental, agentes baseados em LLMs e sistemas multiagente, segurança de
  agentes (envenenamento de memória, injeção de \emph{prompt} e amplificação de
  mau funcionamento) e técnicas de \emph{living-off-the-land} mapeadas ao MITRE
  ATT\&CK.

  \item Especificar a arquitetura MAS-Hunt em três camadas — Conselho de
  Governança, Gerentes e Trabalhadores/Caçadores — detalhando os papéis dos
  agentes, os protocolos de mensagens tipadas e os cinco mecanismos de
  governança M1 (integridade de memória), M2 (validação cruzada), M3
  (resistência a injeção), M4 (monitoramento comportamental) e M5 (quarentena).

  \item Implementar o MAS-Hunt como um \emph{plugin} do Claude Code, materializado
  por meio de habilidades (\emph{skills}), comandos de barra (\emph{slash
  commands}) e \emph{hooks} que orquestram subagentes pelas capacidades nativas
  do agente principal.

  \item Construir um laboratório de avaliação instrumentado com Sysmon e Elastic
  Security sobre um domínio Windows, com um corpus de \emph{playbooks} de
  técnicas LOLBin (malicioso versus benigno) e um desenho adversarial que combina
  modos de injeção e níveis de severidade.

  \item Avaliar empiricamente o MAS-Hunt completo (C1-full) frente às condições de
  referência (C1-nohardening, C2-naive e C3-Elastic) quanto ao desempenho de
  detecção, à resiliência adversarial e à taxa de falsos positivos, com análise
  estatística por intervalos de confiança via \emph{bootstrap}.
\end{enumerate}

\section{Hipótese}
\label{sec:hipotese}

A hipótese central desta dissertação é a seguinte: \emph{a arquitetura MAS-Hunt,
ao incorporar mecanismos de governança multiagente (M1--M5), apresenta uma
redução estatisticamente significativa da suscetibilidade a ataques de
envenenamento de memória e de amplificação de mau funcionamento — medida pela
taxa de sucesso do ataque (\emph{Attack Success Rate}) — em comparação a uma
linha de base multiagente sem endurecimento e a um agente de passagem única,
preservando ou aprimorando o desempenho de detecção e reduzindo a taxa de falsos
positivos em relação às regras de detecção padrão do Elastic Security.}

Dessa hipótese central derivam-se duas hipóteses operacionais subordinadas: (i) a
governança contribui de forma marginal e mensurável para a resiliência, de modo
que a ablação dos mecanismos de endurecimento (C1-nohardening) eleva a taxa de
sucesso do ataque em relação ao sistema completo (C1-full); e (ii) a orquestração
multiagente governada reduz a taxa de falsos positivos em comparação tanto ao
agente de passagem única (C2-naive) quanto às regras padrão do Elastic
(C3-Elastic), sem prejuízo da revocação.

\section{Contribuições}
\label{sec:contribuicoes}

As principais contribuições desta dissertação são:

\begin{enumerate}
  \item \textbf{Uma síntese crítica} do panorama atual de ameaças e do
  estado da arte em vulnerabilidades de agentes de IA, articulando os riscos da
  camada de agentes (envenenamento de memória, injeção de \emph{prompt} e
  amplificação de mau funcionamento) às exigências operacionais da caça a
  ameaças em ambientes Windows.

  \item \textbf{A arquitetura MAS-Hunt}, uma proposta de orquestração multiagente
  em três camadas com cinco mecanismos de governança (M1--M5) e um
  \emph{framework} estruturado de análise (OBSERVAR $\rightarrow$ HIPOTETIZAR
  $\rightarrow$ EVIDÊNCIA $\rightarrow$ VALIDAÇÃO-CRUZADA $\rightarrow$
  CLASSIFICAR) que confere resistência a injeção adversarial.

  \item \textbf{Uma implementação concreta} do MAS-Hunt como \emph{plugin} do
  Claude Code — composto por habilidades, comandos de barra e \emph{hooks} que
  orquestram subagentes —, demonstrando que a governança multiagente pode ser
  realizada sobre capacidades nativas de orquestração do agente principal, e não
  apenas por meio de \emph{scripts} isolados.

  \item \textbf{Um arcabouço de avaliação reprodutível}, incluindo o laboratório
  instrumentado, o corpus de \emph{playbooks} LOLBin, o desenho adversarial e o
  pipeline de medição com \emph{health-gate} de telemetria, que permite comparar
  arquiteturas governadas e não governadas sob métricas e análise estatística
  bem definidas.
\end{enumerate}

\section{Organização do Texto}
\label{sec:organizacao}

O restante desta dissertação está organizado em sete capítulos. O
Capítulo~\ref{cap:introducao} — este capítulo — apresentou o contexto, o
problema, a lacuna de pesquisa, os objetivos, a hipótese e as contribuições do
trabalho. O Capítulo~2 desenvolve a fundamentação teórica, abrangendo a caça a
ameaças e a detecção comportamental, os LLMs e agentes, os sistemas multiagente,
a segurança de agentes, a telemetria Windows com Sysmon e Elastic Security, as
técnicas de \emph{living-off-the-land} mapeadas ao MITRE ATT\&CK e a injeção
adversarial de \emph{logs}. O Capítulo~3 examina os trabalhos relacionados,
comparando abordagens existentes e situando as lacunas que o MAS-Hunt se propõe a
preencher.

O Capítulo~4 detalha a arquitetura MAS-Hunt: a visão geral das três camadas, os
papéis dos agentes, os mecanismos M1--M5, o \emph{framework} M3 de análise, a
implementação como \emph{plugin} do Claude Code e os protocolos de mensagens
tipadas e de governança. O Capítulo~5 descreve a metodologia experimental,
incluindo o laboratório, o corpus de \emph{playbooks}, o desenho adversarial, as
quatro condições experimentais, o pipeline de medição e a análise estatística,
documentando honestamente as limitações de telemetria e o papel do
\emph{health-gate} na sua mitigação.

O Capítulo~6 apresenta os resultados obtidos; o Capítulo~7 discute suas
implicações, limitações e ameaças à validade; e o Capítulo~8 conclui o trabalho,
recapitulando as contribuições e apontando direções para pesquisas futuras.

%======================================================================
% Capítulo 2 — Fundamentação Teórica
% Dissertação MAS-Hunt (ABNT, PT-BR)
% Cita apenas chaves existentes em ../bibliografia.bib e
% ../mashunt/references.bib (declaradas no main.tex via \bibliography).
%======================================================================

\chapter{Fundamentação Teórica}
\label{cap:fundamentacao}

Este capítulo estabelece o arcabouço conceitual sobre o qual o MAS-Hunt (\textit{Multi-Agent System for Threat Hunting}) é construído e avaliado. A exposição parte da prática de caça a ameaças e da detecção comportamental, percorre os fundamentos dos grandes modelos de linguagem e dos agentes de inteligência artificial, examina os sistemas multiagente e — central para a contribuição desta dissertação — a segurança dos próprios agentes, antes de detalhar a infraestrutura de telemetria do ecossistema Windows/Elastic, as técnicas de abuso de binários nativos (\textit{Living-off-the-Land Binaries}, LOLBin) catalogadas pelo MITRE ATT\&CK e, por fim, a injeção adversarial de logs. Sempre que pertinente, os conceitos são amarrados aos cinco mecanismos de governança do MAS-Hunt — M1 (Integridade de Memória), M2 (Validação Cruzada), M3 (Resistência a Injeção), M4 (Monitoramento Comportamental) e M5 (Quarentena) — e às quatro condições experimentais (C1-full, C1-nohardening, C2-naive e C3-Elastic), de modo que a metodologia do Capítulo~\ref{cap:metodologia} e a discussão dos resultados encontrem aqui o seu embasamento.

%======================================================================
\section{Caça a Ameaças e Detecção Comportamental}
\label{sec:cacaameacas}

A caça a ameaças (\textit{threat hunting}) consolidou-se como uma mudança de paradigma na defesa cibernética, deslocando o eixo da postura reativa — que aguarda o disparo de um alerta — para uma busca ativa e iterativa por indícios de comprometimento que tenham escapado dos controles automatizados \cite{academia22}. A premissa operacional é assumir a violação (\textit{assume breach}): parte-se do princípio de que um adversário pode já estar presente no ambiente e, por conseguinte, o objetivo passa a ser reduzir o tempo de permanência do atacante (\textit{dwell time}) e mitigar o dano antes que os objetivos da intrusão sejam alcançados \cite{hillier2022turninghuntedhunterthreat}. Essa urgência é quantificada por relatórios setoriais: o tempo mediano global de permanência do atacante mantém-se na ordem de poucos dias a semanas, intervalo suficiente para movimentação lateral e exfiltração de dados \cite{mandiant2025mtrends}, ao passo que violações envolvendo terceiros respondem por parcela expressiva dos incidentes registrados \cite{verizon2025dbir}.

\subsection{Limitações das defesas convencionais}
\label{subsec:limitacoes}

As arquiteturas tradicionais, ancoradas em detecção por assinatura, são inerentemente reativas: mostram-se eficazes contra ameaças conhecidas, mas incapazes de identificar artefatos novos, polimórficos ou de dia zero, para os quais nenhuma assinatura existe \cite{ucci2019}. A resposta histórica a essa limitação foi a análise comportamental, conduzida por sistemas de detecção de intrusão baseados em host (HIDS) ou em rede (NIDS). Tais sistemas, contudo, operam frequentemente de forma isolada: um HIDS pode observar um processo suspeito e um NIDS pode sinalizar uma conexão de rede anômala, mas, sem mecanismos de correlação, falham em conectar eventos dispersos na narrativa coerente de um ataque coordenado. A caça a ameaças surge justamente para suprir essa lacuna de correlação por meio da investigação dirigida por hipóteses \cite{10713082}.

\subsection{O ciclo da caça dirigida por hipóteses}
\label{subsec:ciclo}

A modalidade mais madura de caça a ameaças é a dirigida por hipóteses, na qual o analista formula uma conjectura sobre comportamento adversário plausível — por exemplo, ``um atacante está utilizando \texttt{certutil.exe} para baixar um estágio secundário'' — e então interroga a telemetria disponível em busca de evidências que a confirmem ou refutem \cite{10713082}. Esse ciclo articula-se em etapas de observação, formulação de hipótese, coleta de evidências e classificação. A automação dessas etapas por meio de aprendizado de máquina e de modelos de linguagem é objeto de pesquisa crescente, em trabalhos que vão da análise de logs \cite{10616474} à detecção de malware em dispositivos de Internet das Coisas mediante redes neurais recorrentes \cite{HADDADPAJOUH201888}, passando por surveys dedicados à aplicação da caça a ameaças contra ransomware \cite{9791234} e por abordagens de aprendizado federado para detecção de ameaças persistentes avançadas (APT) \cite{9931222}. Sistemas como o HuntGPT demonstraram a viabilidade de acoplar detecção de anomalias por aprendizado de máquina a modelos de linguagem dotados de IA explicável, oferecendo ao analista justificativas em linguagem natural \cite{ali2023huntgptintegratingmachinelearningbased}; já propostas como o TIRA exploram a automação do ciclo completo de caça, detecção e perfilamento de atores \cite{10750987}, e estudos sobre a autonomia da defesa cibernética delineiam o horizonte de longo prazo dessa transição \cite{academia21}. A semente conceitual do framework M3 do MAS-Hunt — a cadeia \textsc{Observar} $\rightarrow$ \textsc{Hipotetizar} $\rightarrow$ \textsc{Evidência} $\rightarrow$ \textsc{Validação-Cruzada} $\rightarrow$ \textsc{Classificar} — é uma formalização operacional desse ciclo, projetada de modo a ser, simultaneamente, um protocolo de raciocínio e uma barreira de resistência a manipulação adversarial, conforme detalhado na Seção~\ref{sec:segagentes}.

\subsection{A falácia da taxa-base e seus efeitos sobre o falso positivo}
\label{subsec:taxabase}

Um obstáculo estatístico fundamental atravessa qualquer sistema de detecção que busque eventos raros: a falácia da taxa-base (\textit{base rate fallacy}). Mesmo um classificador de elevada acurácia produz um volume avassalador de falsos positivos quando o fenômeno-alvo é raro na população observada \cite{dolan-gavitt2025}. Formalmente, pelo teorema de Bayes, a probabilidade \textit{a posteriori} de que um alerta corresponda a uma ameaça real é
\begin{equation}
P(A \mid B) = \frac{P(B \mid A)\,P(A)}{P(B)},
\label{eq:bayes}
\end{equation}
em que $P(A)$ é a prevalência da ameaça, $P(B\mid A)$ a sensibilidade do detector e $P(B)$ a probabilidade total de um alerta. Quando $P(A)$ é muito pequeno — o que é típico em telemetria corporativa, na qual a esmagadora maioria dos eventos é benigna — o termo de falsos positivos no denominador domina, e o valor preditivo positivo desaba, ainda que sensibilidade e especificidade sejam altas. Essa realidade tem duas consequências de projeto para o MAS-Hunt. Primeiro, justifica a escolha da métrica F2 (Seção~\ref{subsec:lolbin}) e a comparação explícita da Taxa de Falsos Positivos (FPR) contra as regras padrão do Elastic Security (condição C3-Elastic). Segundo, motiva a validação determinística: em vez de confiar na saída estocástica do modelo, exige-se evidência verificável — por exemplo, o acesso comprovado a um artefato-isca (\textit{canary}) de identificador único e difícil de adivinhar — antes de elevar a confiança em um achado \cite{dolan-gavitt2025}, desacoplando a descoberta flexível, conduzida pela IA, da validação rígida e confiável, conduzida por código não-IA.

%======================================================================
\section{Grandes Modelos de Linguagem e Agentes de IA}
\label{sec:llms}

\subsection{Modelos de linguagem como mecanismo de raciocínio}
\label{subsec:llm}

Os grandes modelos de linguagem (\textit{Large Language Models}, LLMs) são redes neurais de larga escala fundamentadas na arquitetura \textit{Transformer}, cujo mecanismo de autoatenção permite modelar dependências de longo alcance em sequências de texto \cite{vaswani2017attention}. Treinados sobre vastos corpora, exibem capacidades emergentes de compreensão, síntese e raciocínio em linguagem natural que viabilizam, sem treinamento específico de tarefa, a interpretação de telemetria de segurança, a formulação de hipóteses e a geração de regras de detecção. Tais capacidades vêm sendo ampliadas por técnicas de computação em tempo de inferência, como o raciocínio latente recorrente, que aprofundam o processo deliberativo do modelo sem aumentar proporcionalmente o número de parâmetros \cite{arxiv250205171}. Dois atributos dos LLMs são, entretanto, decisivos para o projeto de um sistema de defesa: sua saída é, por construção, estocástica — o que torna indispensável um mecanismo de validação e realimentação por especialistas \cite{academia22} — e seu comportamento é integralmente condicionado pelo contexto textual fornecido, propriedade que abre a superfície de ataque por injeção de instruções discutida na Seção~\ref{sec:segagentes}.

\subsection{De modelos a agentes}
\label{subsec:agentes}

Um agente de inteligência artificial é uma estrutura de software que estende um LLM com a capacidade de perceber um ambiente, raciocinar sobre objetivos e executar ações por meio de ferramentas — consultas a bancos de dados, chamadas a APIs, execução de scripts — em um laço iterativo de percepção, decisão e atuação \cite{academia20, guo2024survey}. É essa autonomia instrumental que distingue um agente de um simples \textit{chatbot}: o agente não apenas responde, mas age sobre o sistema. No MAS-Hunt, os agentes executores interagem diretamente com o Elastic Stack, formulando e refinando hipóteses, consultando o Elasticsearch em busca de atividade anômala, criando regras de detecção pela API do Elastic Security e enriquecendo achados com inteligência externa. A mesma autonomia que confere poder ao agente, contudo, é também a origem de sua fragilidade: a literatura recente demonstra que agentes autônomos podem ser explorados de maneiras que não se aplicam a modelos passivos \cite{academia20, NEURIPS2024_eb113910}, tema central da Seção~\ref{sec:segagentes}. Vale ainda registrar a natureza de uso dual dessas capacidades: os mesmos princípios agênticos que sustentam a defesa habilitam a ofensiva, como evidenciado por sistemas de teste de intrusão autônomos baseados em aprendizado por reforço \cite{moreno2025analysis} e por agentes ofensivos de segurança \cite{academia20}. Qualquer arquitetura defensiva que não modele um adversário igualmente capaz é, por isso, fundamentalmente incompleta.

%======================================================================
\section{Sistemas Multiagente}
\label{sec:mas}

Um sistema multiagente (\textit{Multi-Agent System}, MAS) é composto por múltiplos agentes autônomos que interagem entre si e com o ambiente para alcançar objetivos individuais ou coletivos. Esse paradigma é particularmente adequado à segurança cibernética, domínio em que a defesa requer a coordenação de sensores distribuídos e de componentes de decisão. A coordenação entre agentes é um problema clássico da inteligência artificial distribuída, estudado sob a ótica de organizações, planos e escalonamento \cite{durfee1993organizations}, e revigorado pelos LLMs, que passaram a atuar como mecanismo de raciocínio dos agentes \cite{agashe2023llmcoordination}. A pesquisa contemporânea explora desde a propagação de intenções e raciocínio para coordenação \cite{qiu2024collaborativeintelligencepropagatingintentions}, a cooperação proativa \cite{DBLP:journals/corr/abs-2308-11339} e o aprendizado de coordenação \textit{zero-shot} \cite{pmlr-v202-li23au}, até arquiteturas hierárquicas de agentes de linguagem para coordenação humano-IA em tempo real \cite{liu2023llm}, a redução de computação redundante por execução localmente centralizada \cite{bai2024reducingredundantcomputationmultiagent} e o uso de modelos de fundação multimodais para coordenação distribuída \cite{mahmud2025distributedmultiagentcoordinationusing}.

\subsection{Topologias de orquestração e robustez}
\label{subsec:topologias}

A escolha da topologia de orquestração tem efeito direto sobre a robustez do sistema. Arquiteturas hierárquicas — nas quais uma camada superior coordena camadas subordinadas — exibem a menor degradação de desempenho sob condições adversariais quando comparadas a topologias plenamente distribuídas ou a redes \textit{peer-to-peer} \cite{huang2024resilience}. Esse achado fundamenta diretamente a arquitetura de três camadas do MAS-Hunt — \textbf{Conselho de Governança} (\textit{Board}) $\rightarrow$ \textbf{Gerentes} (\textit{Managers}) $\rightarrow$ \textbf{Trabalhadores/Caçadores} (\textit{Workers/Hunters}). Abordagens correlatas reforçam a importância do projeto estrutural: a geração de robustez por adversários auxiliares evolutivos demonstra como exposições controladas a ataques podem endurecer a coordenação \cite{Yuan_Zhang_Xue_Yin_Chen_Guan_Li_Qian_Yu_2023}, ao passo que protocolos de tolerância bizantina, inspirados em mecanismos de consenso de cadeia de blocos, buscam impedir que um único agente comprometido subverta a decisão coletiva \cite{doe2023emergent}. O Conselho de Governança do MAS-Hunt incorpora esse princípio por meio de um mecanismo de consenso deliberativo entre múltiplos agentes, no qual a aprovação de uma diretriz estratégica exige supermaioria, elevando a barreira para que uma falha ou comprometimento isolado altere o rumo da operação.

%======================================================================
\section{Segurança de Agentes de IA}
\label{sec:segagentes}

A integração de agentes autônomos à segurança cibernética introduz uma classe de riscos qualitativamente distinta: o próprio defensor torna-se alvo. Esta seção sistematiza as três famílias de ataque que o MAS-Hunt foi projetado para resistir e que motivam os mecanismos M1 a M5. Trata-se da contribuição conceitual nuclear desta dissertação, pois desloca a discussão do desempenho de detecção para a resiliência do agente.

\subsection{Envenenamento de memória e de bases de conhecimento}
\label{subsec:envenenamento}

Os componentes cognitivos do agente — sua memória de longo prazo e as bases de conhecimento consultadas por geração aumentada por recuperação (\textit{Retrieval-Augmented Generation}, RAG) — constituem alvo privilegiado. Demonstrou-se que uma base pode ser envenenada com uma taxa de contaminação inferior a $0{,}1\%$ e ainda assim induzir comportamento malicioso direcionado com taxa de sucesso de ataque superior a $80\%$, forçando o agente a recuperar e a agir sobre conteúdo adulterado \cite{NEURIPS2024_eb113910}. A insidiosidade do ataque reside em sua furtividade: a contaminação é estatisticamente diluída e dificilmente detectável por inspeção superficial. O MAS-Hunt responde com o mecanismo \textbf{M1 (Integridade de Memória)}, que submete os resultados de recuperação a validação multifonte e a protocolos de sanitização, e com o \textbf{M5 (Quarentena)}, que isola dados anômalos — por exemplo, registros com carimbos de tempo fabricados — antes que contaminem o raciocínio subsequente.

\subsection{Injeção de prompt direta e indireta}
\label{subsec:injecao}

A injeção de prompt explora o fato de que um LLM não distingue, no nível do texto, entre instruções legítimas do operador e instruções incorporadas aos dados que processa. Na variante indireta, o atacante esconde diretivas maliciosas em conteúdo externo que o agente virá a ler — uma página web, um documento ou, no cenário desta dissertação, um registro de log — de modo que a simples leitura desse conteúdo redirecione o comportamento do agente. Em agentes integrados a ferramentas, esse vetor é especialmente perigoso, pois a instrução injetada pode disparar ações reais sobre o sistema; o benchmark InjecAgent quantifica precisamente essa exposição em agentes que operam com ferramentas \cite{zhan2024injecagent}. O MAS-Hunt confronta esse vetor com o mecanismo \textbf{M3 (Resistência a Injeção)}, operacionalizado pelo framework de análise \textsc{Observar} $\rightarrow$ \textsc{Hipotetizar} $\rightarrow$ \textsc{Evidência} $\rightarrow$ \textsc{Validação-Cruzada} $\rightarrow$ \textsc{Classificar}, cuja disciplina exige que a classificação derive exclusivamente de observáveis verificáveis da telemetria, sinalizando — e não obedecendo — qualquer conteúdo que se apresente como instrução. É justamente a eficácia de M3 que o desenho experimental adversarial — quatro modos de injeção combinados a três níveis de furtividade (\textit{overt}, \textit{encoded}, \textit{steganographic}) — pretende mensurar por meio do \textit{Attack Success Rate} (ASR) e de sua variante não-atribuída (NASR).

\subsection{Amplificação de mau funcionamento e jailbreaks infecciosos}
\label{subsec:amplificacao}

A terceira família ataca não a memória nem a entrada, mas o próprio processo de raciocínio. Na amplificação de mau funcionamento (\textit{malfunction amplification}), um erro menor é deliberadamente induzido e propaga-se em um laço de realimentação de ações inúteis, degradando o agente a ponto de configurar uma negação de serviço; taxas de falha superiores a $80\%$ foram reportadas \cite{zhang2024breaking}. Em sistemas multiagente, o risco escala: os chamados \textit{jailbreaks} infecciosos demonstram que uma única entrada adversarial pode desencadear uma infecção exponencial através de uma população de agentes \cite{gu2024agent}, e a avaliação de deriva de risco de conteúdo (\textit{content risk drift}) mostra como interações adversariais sucessivas podem deslocar gradualmente o comportamento do sistema para fora de suas balizas \cite{liu2024crda}. Embora princípios formais de contenção tenham sido derivados, defesas práticas permanecem escassas — precisamente a lacuna que o MAS-Hunt visa preencher. O mecanismo \textbf{M4 (Monitoramento Comportamental)} observa as sequências de ação e o consumo de recursos dos agentes para detectar e conter comportamento anômalo, enquanto o \textbf{M2 (Validação Cruzada)} encaminha os achados dos caçadores a agentes de validação dedicados, que corroboram a evidência por scripts determinísticos e inteligência externa antes da criação de um caso, impedindo que o erro de um único agente se converta em decisão coletiva. A robustez conferida pela orquestração hierárquica \cite{huang2024resilience} complementa esses mecanismos ao limitar a propagação do comprometimento entre camadas. O conjunto M1–M5 corresponde ao endurecimento que distingue a condição \textbf{C1-full} (MAS-Hunt completo) da \textbf{C1-nohardening} (apenas o framework M3, sem M1/M2/M4/M5) e da \textbf{C2-naive} (LLM de passagem única, sem governança), permitindo isolar, por ablação, a contribuição de cada camada de defesa.

\subsection{Governança e conformidade de sistemas de IA}
\label{subsec:governanca}

A operação de agentes autônomos em ambientes corporativos críticos suscita exigências de governança e conformidade que transcendem o desempenho técnico. Estudos comparativos de \textit{frameworks} de governança, risco e conformidade — COBIT, NIST CSF, ISO~27001 e a norma ISO~42001:2023, específica para sistemas de gestão de IA — evidenciam que os arcabouços tradicionais precisam evoluir para endereçar os riscos próprios dos LLMs \cite{McIntosh2024, mcintosh2024cobit}, e propõem a validação por especialista humano no laço como elemento crítico para a integração segura e conforme. Orientações setoriais sobre o uso do COBIT na governança de sistemas de IA \cite{isaca2025leveraging} convergem nessa direção. No plano nacional, o tratamento de telemetria de \textit{endpoints} — que pode conter dados pessoais — situa-se no marco normativo brasileiro da Lei Geral de Proteção de Dados \cite{LGPD2018} e do Marco Civil da Internet \cite{MarcoCivil2014}, ao passo que a tipificação de delitos informáticos \cite{Dieckmann2012} e o debate sobre o futuro da governança de cibersegurança no país \cite{FGV2023} contextualizam a dimensão regulatória da pesquisa. O Conselho de Governança do MAS-Hunt materializa, no plano arquitetural, essa preocupação: além de orientar a estratégia da caça, ele mantém o especialista humano no laço para a aprovação de ações de contenção, preservando uma postura de \textit{analyst-in-the-loop} apropriada à caça automatizada.

%======================================================================
\section{Elastic Security, Sysmon e Telemetria Windows}
\label{sec:telemetria}

\subsection{O Elastic Stack como terreno de caça}
\label{subsec:elastic}

O MAS-Hunt opera sobre o Elastic Stack, que funciona simultaneamente como ambiente operacional e única fonte de verdade. O \textbf{Elasticsearch} é o repositório central e a base de conhecimento, armazenando e indexando a telemetria; o \textbf{Elastic Security} fornece o motor nativo de detecção e o sistema de gestão de casos; e os \textbf{Elastic Agents}, implantados nos \textit{endpoints} e gerenciados pelo Fleet, coletam dados em tempo real a partir de integrações como Elastic Defend, eventos do Windows, Packetbeat (rede) e osquery (inspeção profunda de sistema). Todas as ações dos agentes do MAS-Hunt são executadas por chamadas autenticadas à API REST do Elastic, o que torna a plataforma plenamente programável. A opção por operar sobre telemetria viva de \textit{endpoints} reais é deliberada: detectores avaliados apenas em traços comportamentais controlados podem não representar a diversidade ruidosa dos hospedeiros reais, em razão de desvio de distribuição, variabilidade ambiental e técnicas de evasão de \textit{sandbox} discutidas na literatura de análise de malware \cite{ucci2019}. Ao fundamentar a análise em dados \textit{in situ}, o MAS-Hunt reduz essa limitação por construção.

\subsection{Sysmon e os identificadores de evento}
\label{subsec:sysmon}

A espinha dorsal da telemetria comportamental no Windows é o System Monitor (Sysmon), serviço que registra, no log de eventos do sistema, atividades de granularidade fina identificadas por códigos de evento (\textit{Event IDs}, EID). Para a caça a abuso de binários nativos, o evento mais informativo é o \textbf{EID~1 (criação de processo)}, que captura a imagem executada, a linha de comando completa, o processo pai e o usuário responsável — registrando, portanto, a cadeia de filiação de processos que revela o encadeamento de um ataque. Outros eventos relevantes incluem o EID~3 (conexão de rede), o EID~10 (acesso a processo, útil para detectar despejo de credenciais a partir do \texttt{lsass.exe}) e os EID~12, 13 e 14 (criação e modificação de registro), frequentemente associados a mecanismos de persistência. Essa granularidade alimenta o mapeamento direto entre comportamento observado e técnicas catalogadas pelo MITRE ATT\&CK, discutido na Seção~\ref{sec:lolbinattack}.

\subsection{A lacuna de telemetria do EID~1 e o \textit{health-gate}}
\label{subsec:eid1gap}

Um aspecto metodologicamente delicado, documentado honestamente nesta dissertação, é a \textbf{lacuna de telemetria do EID~1}. Na prática do laboratório, a esmagadora maioria dos documentos Sysmon corresponde a eventos de registro (EID~12/13), de modo que, em execuções de baixo sinal, as criações de processo dos LOLBin maliciosos podem ser esparsamente capturadas — ou ausentes — caso os agentes e o Sysmon não tenham concluído o \textit{check-in} junto ao Fleet no instante da execução do \textit{playbook}. Trata-se de uma limitação estrutural da instrumentação, não de uma falha de detecção: nenhum classificador pode reconhecer um processo cuja criação jamais foi registrada. A mitigação adotada é dupla. Primeiro, um \textbf{\textit{health-gate} de telemetria} substitui esperas fixas por um sinal de prontidão fim-a-fim: o ambiente só executa os \textit{playbooks} após confirmar, via Elasticsearch, que cada hospedeiro embarcou um número mínimo de eventos EID~1 frescos e que seu canal de execução remota está operante, sondando a cadeia completa \mbox{execução} $\rightarrow$ Sysmon $\rightarrow$ Fleet $\rightarrow$ Elasticsearch. Segundo, uma extração de janela priorizada por EID~1 garante a inclusão de todas as criações de processo disponíveis na janela de cada \textit{playbook}, antes de preencher o restante com eventos cronológicos, recuperando integralmente a cadeia de ataque em execuções de sinal rico. Essas salvaguardas são parte indissociável do \textit{pipeline} de medição e condicionam a validade interna das comparações entre as condições experimentais.

%======================================================================
\section{Técnicas LOLBin e o MITRE ATT\&CK}
\label{sec:lolbinattack}

\subsection{O framework MITRE ATT\&CK}
\label{subsec:attack}

O MITRE ATT\&CK é uma base de conhecimento, reconhecida internacionalmente, de táticas, técnicas e procedimentos (TTPs) adversariais observados no mundo real, organizada em táticas (o ``porquê'', como Persistência, Evasão de Defesa, Acesso a Credenciais, Coleta, Comando-e-Controle e Exfiltração) e técnicas (o ``como''). O framework oferece uma linguagem comum para descrever comportamento de atacante e mapear sistematicamente cada detecção a um identificador de técnica, o que confere rigor e interoperabilidade à caça. No MAS-Hunt, a lógica de detecção é explicitamente ancorada nas técnicas ATT\&CK relevantes — por exemplo, chaves de execução de registro (T1547.001), criação ou modificação de serviços do Windows (T1543.003), limpeza de logs de eventos (T1070.001), despejo de credenciais do sistema operacional (T1003), \textit{keylogging} (T1056.001), captura de tela (T1113) e exfiltração por canal de Comando-e-Controle (T1041).

\subsection{Binários do tipo Living-off-the-Land}
\label{subsec:lolbin}

As técnicas \textit{Living-off-the-Land} (LOLBin) consistem no abuso de binários, scripts e bibliotecas legítimos, já presentes no sistema operacional, para executar ações maliciosas. Por se valerem de ferramentas confiáveis e assinadas pelo próprio fornecedor — como \texttt{powershell.exe}, \texttt{certutil.exe}, \texttt{bitsadmin.exe}, \texttt{rundll32.exe} e \texttt{regsvr32.exe} —, essas técnicas evadem controles baseados em assinatura e em listas de permissão, dissolvendo a fronteira entre atividade administrativa benigna e atividade maliciosa. Justamente por isso, a discriminação entre uso legítimo e uso abusivo do mesmo binário exige análise comportamental contextual: o sinal não está na presença do executável, mas no encadeamento de seus argumentos, em sua filiação de processo e em sua correlação com conexões de rede e modificações de persistência. Essa dificuldade fundamenta o desenho do corpus experimental desta dissertação — \textbf{185 \textit{playbooks}} distribuídos por \textbf{nove famílias de LOLBin}, com instâncias maliciosas e benignas pareadas por família, de modo a forçar o sistema a decidir pelo comportamento e não pelo mero nome do binário. A predominância de exemplares benignos torna o cenário fiel à falácia da taxa-base (Seção~\ref{subsec:taxabase}) e justifica a adoção da métrica \textbf{F2}, que pondera a revocação (\textit{recall}) com peso quatro vezes superior ao da precisão ($\beta = 2$), refletindo o custo assimétrico de não detectar uma ameaça real frente ao de gerar um alarme adicional. À métrica F2 somam-se Precisão, Revocação e a Taxa de Falsos Positivos, todas reportadas com intervalos de confiança estimados por \textit{bootstrap}.

%======================================================================
\section{Injeção Adversarial de Logs}
\label{sec:injecaologs}

A injeção adversarial de logs é o vetor de ataque específico contra o qual a resiliência do MAS-Hunt é mensurada e que materializa, no domínio da caça a ameaças, a injeção de prompt indireta da Seção~\ref{subsec:injecao}. A premissa é que um adversário com a capacidade de escrever em logs — diretamente ou induzindo o sistema a registrar entradas controladas — pode incorporar à telemetria conteúdo textual cujo propósito não é descrever um evento, mas manipular o agente que virá a lê-lo. Um \textit{payload} de injeção pode, por exemplo, apresentar-se como uma diretiva de sistema instruindo o agente a ignorar a análise anterior e a classificar uma janela maliciosa como benigna, induzindo um falso negativo, ou o inverso, induzindo um falso positivo que sobrecarrega o analista. O perigo desse vetor é proporcional à autonomia do agente: quanto mais o agente confia na telemetria como fonte de verdade, mais exposto fica à telemetria adulterada.

O desenho experimental desta dissertação operacionaliza a ameaça em duas dimensões ortogonais. A primeira é o \textbf{modo de injeção}, que governa o momento e a relação temporal entre a injeção e a atividade-alvo — variando de rajadas prévias à execução (\textit{pre-burst}), passando por rajadas acopladas a cada ataque (\textit{per-attack-burst}), até a intercalação contínua (\textit{interleaved}) — além de um modo de controle sem injeção. A segunda é o \textbf{nível de furtividade} do \textit{payload}: \textit{overt} (instrução explícita em texto claro), \textit{encoded} (codificada, por exemplo em \textit{base64}) e \textit{steganographic} (oculta em campos aparentemente legítimos). A resiliência é então quantificada por meio do \textit{Attack Success Rate} (ASR) — a fração de injeções que efetivamente induziu a classificação incorreta — e do \textit{Non-Attributed Success Rate} (NASR), que restringe o sucesso aos casos em que o ataque não foi sequer detectado pelo agente, capturando, assim, não só se a defesa falhou, mas se falhou silenciosamente. A hipótese central da pesquisa é que a governança multiagente endurecida (C1-full) reduz, de forma estatisticamente significativa, o ASR e o NASR frente às linhas de base sem endurecimento (C1-nohardening e C2-naive), evidenciando que os mecanismos M1–M5 — em especial o M3 — conferem ao sistema resistência efetiva à injeção adversarial de logs.

%======================================================================
\section{Síntese do Capítulo}
\label{sec:sintese}

Este capítulo percorreu o caminho conceitual que vai da prática da caça a ameaças, motivada pela falência da detecção reativa e pela urgência de reduzir o tempo de permanência do atacante, até a ameaça específica da injeção adversarial de logs. Estabeleceu-se que os LLMs e os agentes de IA oferecem um mecanismo poderoso de automação do ciclo de caça dirigido por hipóteses, mas que sua autonomia e sua dependência do contexto textual os tornam vulneráveis a uma classe inédita de ataques — envenenamento de memória, injeção de prompt e amplificação de mau funcionamento. Argumentou-se que a orquestração multiagente hierárquica, fundamentada em evidência empírica de robustez sob adversidade, constitui resposta arquitetural a esses riscos, e que os cinco mecanismos de governança do MAS-Hunt (M1–M5) endereçam, cada um, uma face específica da ameaça. Por fim, fundamentaram-se a infraestrutura de telemetria Windows/Sysmon/Elastic sobre a qual o sistema opera, a lacuna de telemetria do EID~1 e sua mitigação por \textit{health-gate}, as técnicas LOLBin catalogadas pelo MITRE ATT\&CK que compõem o corpus experimental, e a racionalidade das métricas (F2, FPR, ASR, NASR) e das quatro condições (C1-full, C1-nohardening, C2-naive, C3-Elastic). Os capítulos seguintes detalham, respectivamente, a arquitetura do MAS-Hunt (Capítulo~\ref{cap:arquitetura}) e a metodologia experimental (Capítulo~\ref{cap:metodologia}) que põem à prova as proposições aqui fundamentadas.

% =====================================================================
% Capítulo 3 — Trabalhos Relacionados
% Dissertação de Mestrado MAS-Hunt (ABNT, PT-BR)
% Convenções: ver .orchestration/dissertacao-OUTLINE.md
% Citações: chaves reais de bibliografia.bib e mashunt/references.bib
% NOTA DE INTEGRAÇÃO (cite gap): chaves dos eixos (a),(b),(d) vêm de
%   bibliografia.bib; chaves do eixo (c) — adversariais — vêm de
%   mashunt/references.bib. As duas bases são disjuntas; o main.tex deve
%   apontar \bibliography para ambas (ou mesclá-las) para resolver todas as
%   chaves abaixo. A tabela usa \ding{} do pacote pifont.
% =====================================================================

\chapter{Trabalhos Relacionados}
\label{cap:trabalhos-relacionados}

A proposta do MAS-Hunt situa-se na confluência de quatro linhas de pesquisa que,
até o momento, evoluíram de modo amplamente independente: (i) a aplicação de
\textit{Large Language Models} (LLMs) e agentes autônomos à caça a ameaças e à
detecção comportamental; (ii) os sistemas multiagente para segurança
cibernética e sua governança; (iii) os ataques adversariais dirigidos aos
próprios agentes baseados em LLM, em especial o envenenamento de memória e a
injeção de prompt/logs; e (iv) a automação de detecção em ambientes
SIEM/SOC. Este capítulo organiza a literatura segundo esses quatro eixos
temáticos, compara criticamente as abordagens dentro de cada eixo e, ao final,
explicita a lacuna que o MAS-Hunt se propõe a preencher. A revisão é deliberadamente
comparativa: o objetivo não é catalogar trabalhos, mas evidenciar que nenhuma das
linhas existentes integra, em um único artefato avaliado empiricamente,
governança multiagente real e resistência a injeção adversarial sobre telemetria
viva.

% =====================================================================
\section{Caça a Ameaças e Detecção Comportamental Assistidas por IA}
\label{sec:rel-hunting-ia}

A caça a ameaças (\textit{threat hunting}) consolidou-se como um deslocamento
paradigmático da postura reativa, baseada em assinaturas, para uma busca
proativa e iterativa por adversários já presentes no ambiente
\cite{hillier2022turninghuntedhunterthreat}. \citeonline{hillier2022turninghuntedhunterthreat}
descrevem o ciclo de vida da caça a ameaças e seu ecossistema, antecipando a
\textit{grande promessa} da inteligência artificial como força multiplicadora do
analista humano; o trabalho, contudo, permanece no plano conceitual, sem
implementação operacional. Essa promessa é sistematizada de forma prospectiva por
\citeonline{academia22}, que propõe a \textit{caça autônoma a ameaças} como
paradigma futuro de inteligência de ameaças dirigida por IA, e por
\citeonline{academia21}, que mapeia o \textit{caminho para a defesa cibernética
autônoma}. Ambos os trabalhos, entretanto, são essencialmente revisões de visão:
identificam direções desejáveis sem entregar uma arquitetura validada nem
quantificar resiliência.

No plano das técnicas concretas de detecção, há uma trajetória clara que vai do
aprendizado de máquina supervisionado em domínio fechado até o uso de LLMs como
motores de raciocínio. \citeonline{HADDADPAJOUH201888} aplicam redes neurais
recorrentes profundas à caça a \textit{malware} em dispositivos de Internet das
Coisas, alcançando bons índices de classificação, porém à custa de um modelo
treinado de modo \textit{offline}, incapaz de formular hipóteses ou de interagir
com a telemetria em tempo de execução. As limitações intrínsecas dessa família de
abordagens são sintetizadas por \citeonline{ucci2019}, cujo levantamento sobre
técnicas de aprendizado de máquina para análise de \textit{malware} evidencia a
dependência de \textit{features} predominantemente sintáticas e a fragilidade
frente à mudança de distribuição entre o ambiente de treino e o \textit{endpoint}
real.

A introdução de LLMs altera qualitativamente esse panorama ao acoplar raciocínio
em linguagem natural e capacidade de geração de consultas. O HuntGPT
\cite{ali2023huntgptintegratingmachinelearningbased} é representativo dessa
geração: integra detecção de anomalias baseada em aprendizado de máquina,
inteligência artificial explicável e um LLM como camada de interpretação,
oferecendo ao analista justificativas legíveis para os alertas. Trata-se,
porém, de um agente único e assistivo, sem orquestração colaborativa nem
qualquer hipótese de defesa contra manipulação do próprio modelo. O modelo de
geração de hipóteses de \citeonline{10713082} formaliza o estágio mais cognitivo
da caça --- a formulação de hipóteses verificáveis ---, mas o faz como
componente isolado, desacoplado de execução, validação ou resposta. O TIRA, de
\citeonline{10750987}, avança rumo à automação de ponta a ponta (caça, detecção e
perfilamento de atores), aproximando-se de um fluxo operacional completo;
todavia, não modela um adversário inteligente capaz de subverter o próprio
sistema de caça, premissa central do MAS-Hunt.

Em síntese, o eixo de caça assistida por IA progrediu da classificação
\textit{offline} \cite{HADDADPAJOUH201888,ucci2019} para a interpretação
explicável \cite{ali2023huntgptintegratingmachinelearningbased} e a geração de
hipóteses \cite{10713082,10750987}, mas mantém duas omissões recorrentes: a
ausência de \textit{colaboração estruturada entre múltiplos agentes
especializados} e a desconsideração da \textit{superfície de ataque do próprio
sistema de IA}.

% =====================================================================
\section{Sistemas Multiagente para Segurança e sua Governança}
\label{sec:rel-mas-governanca}

A coordenação entre múltiplos agentes para resolver problemas distribuídos é um
campo maduro, cujas raízes remontam à formalização de organizações, planos e
escalonamento como mecanismos de coordenação de agentes de IA
\cite{durfee1993organizations}. O advento dos LLMs reabilitou esse campo sob nova
forma: agentes de linguagem passaram a coordenar-se em arranjos hierárquicos para
tarefas complexas. \citeonline{liu2023llm} propõem um agente hierárquico de
linguagem para coordenação humano-IA em tempo real, demonstrando que a
decomposição em camadas --- um agente de alto nível que delega a subagentes ---
melhora a latência e a qualidade da cooperação. \citeonline{agashe2023llmcoordination}
avaliam sistematicamente as capacidades de coordenação multiagente de diferentes
LLMs, e \citeonline{qiu2024collaborativeintelligencepropagatingintentions}
mostram que a propagação explícita de intenções e raciocínio entre agentes
aprimora a inteligência coletiva. Padrões cooperativos correlatos aparecem no
ProAgent \cite{DBLP:journals/corr/abs-2308-11339}, na coordenação \textit{zero-shot}
de \citeonline{pmlr-v202-li23au} e na coordenação fundamentada em modelos
multimodais de \citeonline{mahmud2025distributedmultiagentcoordinationusing}.
Essa literatura, contudo, otimiza \textit{desempenho cooperativo} em ambientes
benignos: a coordenação é tratada como problema de eficiência
\cite{bai2024reducingredundantcomputationmultiagent}, não de segurança.

Um subconjunto da literatura começa a tratar a coordenação como problema
adversarial e de confiança. \citeonline{Yuan_Zhang_Xue_Yin_Chen_Guan_Li_Qian_Yu_2023}
propõem coordenação multiagente robusta via geração evolutiva de atacantes
auxiliares, endurecendo políticas contra perturbações; o BlockAgents
\cite{doe2023emergent} busca coordenação tolerante a falhas bizantinas por meio de
\textit{blockchain}, atacando o problema do agente comprometido pela via da
imutabilidade do registro. Esses trabalhos, entretanto, situam-se no domínio de
aprendizado por reforço multiagente ou de consenso distribuído genérico, sem
ancoragem em telemetria de segurança real nem nos vetores de ataque específicos a
LLMs (envenenamento de memória, injeção de prompt).

Paralelamente, emerge uma literatura de \textit{governança} de sistemas de IA que
fornece o vocabulário normativo --- mas não a implementação técnica --- para o
controle de agentes autônomos. \citeonline{isaca2025leveraging} adaptam o
arcabouço COBIT à governança de sistemas de IA; \citeonline{mcintosh2024cobit} e
\citeonline{McIntosh2024} avaliam a transição do COBIT para a ISO/IEC 42001 sob a
ótica de oportunidades, riscos e conformidade regulatória na comercialização de
LLMs. Esses estudos estabelecem \textit{por que} a governança é necessária e
\textit{quais} controles deveriam existir, mas não descem ao nível de protocolos
executáveis de mensagens tipadas, votação de conselho ou quarentena de agentes ---
exatamente o terreno onde o MAS-Hunt opera. A obra de referência sobre
\textit{frameworks} de agentes \cite{guo2024survey} complementa esse quadro no
plano da engenharia, mas sem viés de segurança.

A comparação entre os dois subcampos revela uma assimetria reveladora: a
coordenação multiagente sabe \textit{como} fazer agentes cooperarem, e a
governança de IA sabe \textit{quais} controles são desejáveis, porém raramente os
dois se encontram em um sistema concreto de segurança cibernética cuja
arquitetura de governança seja ela própria o mecanismo de defesa.

% =====================================================================
\section{Ataques Adversariais a Agentes Baseados em LLM}
\label{sec:rel-ataques-adversariais}

O terceiro eixo é o que motiva diretamente o endurecimento do MAS-Hunt: a
constatação empírica de que a autonomia que torna os agentes poderosos os torna
igualmente frágeis. A superfície de ataque mais crítica é a memória cognitiva do
agente. \citeonline{NEURIPS2024_eb113910}, com o AgentPoison, demonstram que a
memória ou a base de conhecimento de recuperação aumentada (RAG) de um agente
pode ser envenenada com taxa de injeção inferior a 0{,}1\% para induzir
comportamento malicioso direcionado com taxa de sucesso de ataque superior a
80\%. Esse resultado é decisivo: mostra que a integridade da memória não pode ser
presumida e fundamenta o mecanismo M1 (Integridade de Memória) do MAS-Hunt.

Uma segunda classe de ataques explora o \textit{loop de raciocínio} do agente.
\citeonline{zhang2024breaking}, em \textit{Breaking Agents}, caracterizam a
\textit{amplificação de mau funcionamento} (\textit{malfunction amplification}),
em que um erro menor desencadeia um laço de realimentação de ações inúteis, com
taxas de falha que ultrapassam 80\% --- um vetor efetivo de negação de serviço
contra o próprio agente. A injeção indireta de \textit{prompt}, por sua vez, é
quantificada pelo \textit{benchmark} InjecAgent \cite{zhan2024injecagent}, que
mede a vulnerabilidade de agentes integrados a ferramentas a instruções
maliciosas embutidas em conteúdo externo --- precisamente o cenário da
injeção adversarial de logs, em que o adversário contamina a telemetria que o
caçador irá ler. O CRDA \cite{liu2024crda} estende a análise ao \textit{drift} de
risco de conteúdo via interação adversarial multiagente, evidenciando que a
degradação pode emergir da própria dinâmica entre agentes.

O caráter sistêmico dessas ameaças em \textit{coletivos} de agentes é o achado
mais alarmante. \citeonline{gu2024agent}, com o Agent Smith, mostram que uma
única imagem adversarial pode desencadear um \textit{jailbreak infeccioso} que se
propaga exponencialmente por um milhão de agentes multimodais.
\citeonline{huang2024resilience} estudam a resiliência de sistemas multiagente na
presença de agentes maliciosos e apresentam o resultado de projeto mais relevante
para esta dissertação: estruturas de orquestração \textit{hierárquicas} exibem a
menor degradação de desempenho sob condições adversariais quando comparadas a
topologias alternativas. Esse achado justifica diretamente a escolha arquitetural
do MAS-Hunt por uma hierarquia de três camadas. Por fim, essas vulnerabilidades
devem ser lidas à luz de um adversário ativo e inteligente:
\citeonline{moreno2025analysis} propõem teste de intrusão autônomo via
aprendizado por reforço e sistemas de recomendação, e \citeonline{academia20}
discute a IA como o \textit{novo atacante}, desenvolvendo agentes para segurança
ofensiva. Um agente ofensivo dessa natureza é o candidato natural a executar os
ataques de envenenamento e de manipulação comportamental acima descritos --- donde
a conclusão de que qualquer arquitetura defensiva que não modele um adversário de
IA igualmente capaz é fundamentalmente incompleta.

A leitura conjunta deste eixo estabelece a tríade de ameaças que o MAS-Hunt
endereça: envenenamento de memória \cite{NEURIPS2024_eb113910}, amplificação de
mau funcionamento \cite{zhang2024breaking} e injeção/propagação infecciosa
\cite{zhan2024injecagent,gu2024agent,huang2024resilience}. Crucialmente, a
literatura deste eixo \textit{ataca} e \textit{mede}, mas raramente \textit{defende}:
as defesas práticas, quando propostas, permanecem em ambiente simulado, desacopladas
de qualquer tarefa de segurança operacional como a caça a ameaças.

% =====================================================================
\section{Automação de Detecção em Ambientes SIEM/SOC}
\label{sec:rel-siem-soc}

O quarto eixo trata da automação da detecção no contexto operacional de um centro
de operações de segurança (SOC) apoiado por um \textit{Security Information and
Event Management} (SIEM). A caça a \textit{ransomware} via técnicas de
\textit{threat hunting} é levantada por \citeonline{9791234}, que sistematiza o
estado da arte e expõe a dependência de regras e heurísticas construídas
manualmente. A análise de logs como substrato primário da caça é tratada por
\citeonline{10616474}, evidenciando tanto o valor da telemetria bruta quanto sua
suscetibilidade à manipulação --- um log adulterado é um vetor de injeção
adversarial direto. Em ambientes distribuídos, \citeonline{9931222} aplicam
aprendizado federado à caça a ameaças para detecção de ataques de Ameaça
Persistente Avançada (APT) em redes definidas por software, demonstrando
detecção colaborativa que preserva privacidade, porém sem agentes de raciocínio
nem governança sobre eles.

O obstáculo estatístico que assombra toda automação de detecção em larga escala é
a falácia da taxa de base bayesiana, articulada por \citeonline{dolan-gavitt2025}:
mesmo modelos altamente acurados geram um volume avassalador de falsos positivos
quando o evento procurado é raro, tornando-os operacionalmente inúteis. A resposta
proposta naquele trabalho --- validação determinística baseada em \textit{canaries}
e provas de posse, que mira uma taxa de falsos positivos próxima de zero ---
informa diretamente o mecanismo M2 (Validação Cruzada) do MAS-Hunt, que exige
corroboração por evidência verificável antes de qualquer escalonamento. No plano
da detecção comportamental sobre o \textit{endpoint},
A literatura de análise de malware documenta a sensibilidade de detectores a
traços comportamentais e ao contexto de execução \cite{ucci2019}, o que motiva a
opção do MAS-Hunt por operar diretamente sobre telemetria viva da pilha Elastic,
e não apenas sobre traços sintéticos. Abordagens correlatas que integram LLMs a
detecção anômala explicável \cite{ali2023huntgptintegratingmachinelearningbased}
reforçam a tendência, mas mantêm o paradigma de modelo único, sem orquestração
nem \textit{hardening}.

A comparação neste eixo expõe um padrão consistente: a automação SIEM/SOC ou
delega a regras determinísticas frágeis e ruidosas \cite{9791234,10616474}, ou
emprega modelos estatísticos vulneráveis à falácia da taxa de base
\cite{dolan-gavitt2025} e à mudança de distribuição \cite{ucci2019}. Em
nenhum dos casos a própria camada de automação é tratada como alvo a ser
endurecido contra manipulação adversarial.

% =====================================================================
\section{Síntese Comparativa}
\label{sec:rel-sintese}

A Tabela~\ref{tab:rel-comparacao} confronta representantes de cada eixo segundo
cinco dimensões que delimitam a contribuição do MAS-Hunt: orquestração multiagente
real (não um agente único nem um \textit{pipeline} monolítico); governança
explícita (conselho, mensagens tipadas, papéis); resiliência a ataques
adversariais ao próprio agente (memória, injeção, mau funcionamento); operação
sobre telemetria viva de produção; e avaliação empírica controlada da resiliência.

\begin{table}[htbp]
\centering
\caption{Comparação dos trabalhos relacionados segundo as dimensões críticas do MAS-Hunt. Legenda: \ding{51}~contemplado; $\sim$~parcial/conceitual; \ding{55}~não contemplado.}
\label{tab:rel-comparacao}
\small
\begin{tabular}{p{4.4cm}ccccc}
\hline
\textbf{Trabalho / Eixo} & \textbf{Multi-} & \textbf{Gover-} & \textbf{Resil.} & \textbf{Telem.} & \textbf{Aval.} \\
                         & \textbf{agente} & \textbf{nança}  & \textbf{advers.} & \textbf{viva}  & \textbf{empír.} \\
\hline
\multicolumn{6}{l}{\textit{(a) Caça/detecção assistida por IA}} \\
HuntGPT \cite{ali2023huntgptintegratingmachinelearningbased}        & \ding{55} & \ding{55} & \ding{55} & $\sim$    & \ding{51} \\
Ger.\ de hipóteses \cite{10713082}                                  & \ding{55} & \ding{55} & \ding{55} & \ding{55} & $\sim$    \\
TIRA \cite{10750987}                                                & $\sim$    & \ding{55} & \ding{55} & $\sim$    & \ding{51} \\
Caça autônoma \cite{academia22}                                     & $\sim$    & \ding{55} & \ding{55} & \ding{55} & \ding{55} \\
\hline
\multicolumn{6}{l}{\textit{(b) Multiagente e governança}} \\
Agente hierárquico \cite{liu2023llm}                                & \ding{51} & $\sim$    & \ding{55} & \ding{55} & \ding{51} \\
BlockAgents \cite{doe2023emergent}                                  & \ding{51} & $\sim$    & $\sim$    & \ding{55} & \ding{51} \\
Coord.\ robusta \cite{Yuan_Zhang_Xue_Yin_Chen_Guan_Li_Qian_Yu_2023} & \ding{51} & \ding{55} & $\sim$    & \ding{55} & \ding{51} \\
Governança COBIT/ISO \cite{mcintosh2024cobit}                       & \ding{55} & \ding{51} & \ding{55} & \ding{55} & \ding{55} \\
\hline
\multicolumn{6}{l}{\textit{(c) Ataques adversariais a agentes}} \\
AgentPoison \cite{NEURIPS2024_eb113910}                             & $\sim$    & \ding{55} & \ding{55} & \ding{55} & \ding{51} \\
Breaking Agents \cite{zhang2024breaking}                            & $\sim$    & \ding{55} & \ding{55} & \ding{55} & \ding{51} \\
Agent Smith \cite{gu2024agent}                                      & \ding{51} & \ding{55} & \ding{55} & \ding{55} & \ding{51} \\
Resiliência MAS \cite{huang2024resilience}                          & \ding{51} & $\sim$    & $\sim$    & \ding{55} & \ding{51} \\
\hline
\multicolumn{6}{l}{\textit{(d) Automação SIEM/SOC}} \\
Caça a \textit{ransomware} \cite{9791234}                           & \ding{55} & \ding{55} & \ding{55} & $\sim$    & \ding{55} \\
Caça federada APT \cite{9931222}                                    & \ding{51} & \ding{55} & \ding{55} & $\sim$    & \ding{51} \\
Validação determin.\ \cite{dolan-gavitt2025}                        & \ding{55} & \ding{55} & $\sim$    & $\sim$    & \ding{51} \\
\hline
\textbf{MAS-Hunt (este trabalho)} & \ding{51} & \ding{51} & \ding{51} & \ding{51} & \ding{51} \\
\hline
\end{tabular}
\end{table}

A leitura da Tabela~\ref{tab:rel-comparacao} torna a lacuna manifesta. Os trabalhos
do eixo (a) operam, na melhor das hipóteses, com um agente único sobre telemetria,
sem governança nem modelo de adversário. Os do eixo (b) dominam a orquestração e o
discurso de governança, mas não os confrontam com ataques concretos a agentes nem
com telemetria de segurança real. Os do eixo (c) demonstram exaustivamente \textit{que}
os agentes são vulneráveis, porém em laboratório simulado e sem propor defesa
operacional integrada a uma tarefa de caça. Os do eixo (d) automatizam detecção,
mas tratam o componente de IA como caixa estatística, jamais como alvo a ser
endurecido. Nenhuma linha, isoladamente, ocupa simultaneamente as cinco colunas.

% =====================================================================
\section{A Lacuna que o MAS-Hunt Preenche}
\label{sec:rel-lacunas}

Da análise comparativa precedente decorre a lacuna central que esta dissertação
endereça. As quatro linhas de pesquisa convergem para um mesmo ponto cego: cada
uma resolve uma fração do problema, mas nenhuma integra, em um único sistema
avaliado empiricamente, \textit{governança multiagente real} e \textit{resistência
a injeção adversarial} sobre dados de produção. Especificamente, identificam-se
três deficiências que o MAS-Hunt se propõe a sanar:

\begin{enumerate}
    \item \textbf{Ausência de orquestração multiagente com governança como mecanismo
    de defesa.} A literatura de coordenação \cite{liu2023llm,agashe2023llmcoordination,qiu2024collaborativeintelligencepropagatingintentions}
    otimiza cooperação, e a de governança \cite{isaca2025leveraging,mcintosh2024cobit}
    prescreve controles, mas não há sistema em que a hierarquia de três camadas
    --- Conselho de Governança $\rightarrow$ Gerentes $\rightarrow$ Caçadores ---,
    com mensagens tipadas e protocolos de consenso, atue \textit{ela própria} como
    a barreira que impede a propagação de um agente comprometido. O achado de
    \citeonline{huang2024resilience} de que estruturas hierárquicas degradam menos
    sob ataque é aqui transformado de observação em princípio de projeto operante.

    \item \textbf{Falta de validação de integridade em tempo de execução contra
    ataques a agentes.} Os trabalhos que demonstram envenenamento de memória
    \cite{NEURIPS2024_eb113910}, injeção indireta \cite{zhan2024injecagent} e
    amplificação de mau funcionamento \cite{zhang2024breaking} caracterizam as
    ameaças sem oferecer contramedidas operacionais acopladas a uma tarefa real. O
    MAS-Hunt responde com os mecanismos M1 (Integridade de Memória), M2 (Validação
    Cruzada por evidência verificável, na linha da validação determinística de
    \citeonline{dolan-gavitt2025}), M3 (\textit{framework} de análise estruturada
    OBSERVAR $\rightarrow$ HIPOTETIZAR $\rightarrow$ EVIDÊNCIA $\rightarrow$
    VALIDAÇÃO-CRUZADA $\rightarrow$ CLASSIFICAR, resistente à injeção de logs),
    M4 (Monitoramento Comportamental) e M5 (Quarentena).

    \item \textbf{Inexistência de avaliação empírica de resiliência sobre telemetria
    viva.} As defesas, quando existem, são avaliadas em simulação; as caças, quando
    operam sobre telemetria real \cite{ali2023huntgptintegratingmachinelearningbased},
    não são submetidas a um adversário que ataque o agente. O MAS-Hunt fecha essa
    lacuna ao ser avaliado em laboratório instrumentado com a pilha Elastic,
    confrontando o sistema completo (C1-full) contra a ablação sem endurecimento
    (C1-nohardening), contra uma linha de base de passagem única sem governança
    (C2-naive) e contra as regras padrão do Elastic Security (C3-Elastic), medindo
    tanto desempenho de detecção (Precisão, Revocação, F2) quanto resiliência
    (\textit{Attack Success Rate}) sob injeção adversarial, com intervalos de
    confiança por \textit{bootstrap}.
\end{enumerate}

Em suma, enquanto os trabalhos relacionados respondem isoladamente às perguntas
``como coordenar agentes?'', ``como governar IA?'', ``como atacar agentes?'' e
``como automatizar detecção?'', o MAS-Hunt responde à pergunta integradora que
nenhum deles enfrenta: \textit{como projetar e validar empiricamente um sistema
multiagente de caça a ameaças cuja própria estrutura de governança o torna
resiliente a ataques adversariais dirigidos aos agentes, operando sobre telemetria
de produção?} É essa integração --- e sua verificação experimental --- que
constitui a contribuição original posicionada pelos capítulos seguintes.

% =====================================================================
% Capítulo 4 — Arquitetura MAS-Hunt
% Dissertação MAS-Hunt (ABNT, PT-BR). Capítulo de contribuição central.
% Convenções: M1-M5, 3 camadas (Conselho de Governança -> Gerentes ->
% Trabalhadores/Caçadores), condições C1-full/C1-nohardening/C2-naive/C3-Elastic.
% Citações: apenas chaves existentes em ../bibliografia.bib.
% =====================================================================
\chapter{Arquitetura MAS-Hunt}
\label{cap:arquitetura}

Este capítulo apresenta a contribuição central da dissertação: a arquitetura do
MAS-Hunt (\emph{Multi-Agent System for Threat Hunting}), um sistema multiagente
de busca proativa por ameaças concebido com prioridade à segurança do próprio
sistema de agentes. A exposição parte de uma visão geral da organização
hierárquica em três camadas (Seção~\ref{sec:arq-visao-geral}), detalha os papéis
e responsabilidades dos agentes em cada camada (Seção~\ref{sec:arq-papeis}),
formaliza os cinco mecanismos de governança M1--M5 e o \emph{framework} de
análise estruturada (Seção~\ref{sec:arq-mecanismos} e
Seção~\ref{sec:arq-framework-m3}) e, por fim, descreve a materialização do sistema
como um \emph{plugin} do Claude Code, esclarecendo por que a orquestração
multiagente real difere de forma fundamental de uma chamada isolada a um modelo de
linguagem (Seção~\ref{sec:arq-plugin} e Seção~\ref{sec:arq-protocolos}).

A premissa de projeto é que a autonomia que torna os agentes de IA poderosos é
também a sua principal fragilidade. Agentes baseados em grandes modelos de
linguagem (LLMs) são suscetíveis a envenenamento de
memória~\cite{NEURIPS2024_eb113910}, injeção indireta de
\emph{prompt} via conteúdo externo~\cite{zhan2024injecagent} e amplificação de
mau funcionamento~\cite{zhang2024breaking}. Em um
sistema de caça a ameaças, essa fragilidade é agravada pelo fato de que o
material analisado — registros de telemetria, alertas e inteligência sobre
ameaças — é, por natureza, conteúdo potencialmente hostil e controlado, em parte,
pelo adversário. O MAS-Hunt responde a essa tensão tratando a governança
multiagente não como um adendo operacional, mas como o objeto de projeto
primário: a contribuição não reside em \emph{scripts} de detecção, e sim no
arcabouço de governança que organiza, valida e contém o comportamento dos
agentes.

\section{Visão Geral da Arquitetura em Três Camadas}
\label{sec:arq-visao-geral}

O MAS-Hunt organiza-se como uma hierarquia estrita de três camadas, na qual o
fluxo de comando desce do estratégico para o operacional e o fluxo de evidências
sobe do operacional para o estratégico. As camadas são, do topo para a base:

\begin{enumerate}
  \item \textbf{Conselho de Governança (\emph{Board})} --- camada estratégica
        (\emph{layer}~0). Define objetivos de caça, valida a metodologia,
        delibera por quórum e assina o veredito final da campanha.
  \item \textbf{Gerentes (\emph{Managers})} --- camada de coordenação
        (\emph{layer}~1). Traduz os objetivos estratégicos em planos de execução,
        instancia e supervisiona os caçadores, e consolida as evidências antes da
        validação.
  \item \textbf{Trabalhadores/Caçadores (\emph{Workers/Hunters})} --- camada de
        execução (\emph{layer}~2). Interage diretamente com a telemetria no
        Elastic Stack e com fontes externas de inteligência, classificando
        eventos e produzindo achados fundamentados em evidência.
\end{enumerate}

A escolha por uma organização hierárquica não é arbitrária. A literatura sobre
coordenação multiagente baseada em LLMs indica que estruturas hierárquicas
apresentam a menor degradação de desempenho sob condições adversariais, quando
comparadas a topologias planas ou totalmente conectadas, justamente porque a
hierarquia limita o caminho de propagação de instruções maliciosas ou de
raciocínio comprometido entre
agentes~\cite{liu2023llm,agashe2023llmcoordination}. Em um arranjo plano, um
único agente comprometido pode contaminar todos os seus pares; em uma hierarquia,
a contaminação precisa atravessar fronteiras de camada que são, por projeto,
pontos de mediação e validação. O princípio de coordenação por uma autoridade
central que delega e consolida tarefas é, ademais, consistente com a tradição de
organização de agentes em planos e
escalonamentos~\cite{durfee1993organizations}, aqui adaptada ao contexto de
agentes generativos.

A Figura~\ref{fig:mashunt-arq} apresenta o esquema de alto nível da arquitetura,
ilustrando as três camadas, o sentido descendente das diretivas e o sentido
ascendente dos relatórios e vereditos, bem como os pontos em que os mecanismos de
governança M1--M5 atuam sobre o fluxo de mensagens.

\begin{figure}[htbp]
  \centering
  % A imagem mashunt.png reside em ../mashunt/ (incluída no \graphicspath de main.tex).
  \includegraphics[width=0.95\textwidth]{mashunt.png}
  \caption[Arquitetura em três camadas do MAS-Hunt]{Arquitetura em três camadas
    do MAS-Hunt. As diretivas descem do Conselho de Governança para os Gerentes e
    destes para os Caçadores; os relatórios, achados e vereditos sobem no sentido
    inverso. Os caçadores interagem com o Elastic Stack e com fontes externas de
    inteligência. Os mecanismos M1--M5 atuam nas fronteiras de camada e sobre o
    conteúdo da telemetria.}
  \label{fig:mashunt-arq}
\end{figure}

Três propriedades de projeto perpassam todas as camadas e distinguem o MAS-Hunt
de uma simples cadeia de prompts. A primeira é a \emph{mediação por fronteira de
camada}: nenhuma mensagem trafega diretamente entre o Conselho e os Caçadores sem
atravessar a camada de Gerentes, o que cria pontos de inspeção e impede a
ação de um agente isolado sobre o sistema. A segunda é a \emph{separação entre
descoberta e validação}: a descoberta flexível de achados (atividade dos
caçadores, tolerante a ambiguidade) é desacoplada da validação rígida e
determinística (atividade de agentes e \emph{scripts} validadores), de modo que a
criatividade necessária à caça não comprometa a confiabilidade do veredito. A
terceira é a \emph{rastreabilidade integral}: toda decisão de agente, voto de
governança e transição de tarefa é registrada em uma trilha de auditoria
imutável, condição necessária para a reprodutibilidade científica e para a
análise post-mortem de comportamento anômalo.

A Figura~\ref{fig:mashunt-fluxo} resume o ciclo de vida de uma campanha de caça,
do objetivo estratégico ao veredito assinado, e serve de referência para as
seções seguintes.

\begin{figure}[htbp]
  \centering
  \begin{minipage}{0.92\textwidth}
  \footnotesize
  \setlength{\fboxsep}{6pt}
  \fbox{\parbox{\dimexpr\linewidth-2\fboxsep\relax}{\centering
    \textbf{Conselho de Governança} (\emph{layer}~0)\\
    \emph{define objetivo, valida corpus, vota por quórum} $\Rightarrow$ emite \texttt{HuntGoal}\\[4pt]
    $\big\downarrow$ \emph{diretiva descendente (mensagem tipada e assinada)}\\[4pt]
    \textbf{Gerentes} (\emph{layer}~1)\\
    \emph{planeja lotes de execução, sanitiza amostra, aloca famílias LOLBin} $\Rightarrow$ \texttt{HuntTask}\\[4pt]
    $\big\downarrow$\\[4pt]
    \textbf{Caçadores} (\emph{layer}~2) --- \emph{um por família LOLBin, em paralelo}\\
    \emph{consultam o Elastic Stack, aplicam o framework M3, classificam} $\Rightarrow$ \texttt{HuntFinding}\\[4pt]
    $\big\uparrow$ \emph{achado ascendente}\\[4pt]
    \textbf{Validação} (\emph{layer}~1) --- \emph{validação cruzada determinística (M2)} $\Rightarrow$ \texttt{ValidationVerdict}\\[4pt]
    $\big\uparrow$\\[4pt]
    \textbf{Conselho de Governança} (\emph{layer}~0)\\
    \emph{revisão final contra hipóteses, assina o Dossiê de Caça}
  }}
  \end{minipage}
  \caption[Pipeline de caça do MAS-Hunt]{Pipeline de caça do MAS-Hunt:
    \emph{board} $\rightarrow$ \emph{management} $\rightarrow$ caçadores por
    família LOLBin $\rightarrow$ validação $\rightarrow$ \emph{board}. Cada
    transição entre camadas é mediada por uma mensagem tipada, validada por
    esquema e assinada (Seção~\ref{sec:arq-protocolos}).}
  \label{fig:mashunt-fluxo}
\end{figure}

\section{Papéis e Responsabilidades dos Agentes}
\label{sec:arq-papeis}

\subsection{Camada 1 --- Conselho de Governança}
\label{sec:arq-board}

No ápice da hierarquia situa-se o \textbf{Conselho de Governança}, um colegiado de
agentes com personas distintas responsável pela estratégia e pela direção geral
da caça. O Conselho não interage com a telemetria; sua função é deliberativa e de
controle. Ele inicia o fluxo de trabalho definindo objetivos de alto nível —
hipóteses de caça, escopo de índices e janela temporal, táticas do MITRE ATT\&CK
a priorizar — e atua como ponto de verificação final, recebendo e revisando os
relatórios sintetizados de inteligência, de caça e de validação produzidos pelas
camadas inferiores.

A deliberação do Conselho é governada por dois princípios inegociáveis. O
primeiro é a \emph{exigência de quórum}: nenhuma decisão estratégica é tomada sem
um número mínimo de votos. Na configuração de referência, o Conselho é composto
por três personas — um revisor de metodologia, um avaliador de risco e um diretor
estratégico — e a aprovação de um objetivo de caça requer quórum de dois terços.
O segundo princípio é a \emph{impossibilidade de sobreposição por agente único}
(\emph{no single-agent override}): o mecanismo de consenso impede que um agente
isolado, eventualmente comprometido, imponha uma decisão ao
sistema~\cite{huang2024resilience}. Cada voto é
registrado com sua justificativa e suas condições, e a contagem é persistida na
trilha de auditoria, atendendo ao requisito de transparência. Esse desenho de
governança por colegiado, com tolerância a participantes potencialmente
defeituosos, é informado por abordagens de coordenação multiagente robusta a
falhas bizantinas~\cite{doe2023emergent} e por arcabouços de geração de
adversários auxiliares para endurecer a
coordenação~\cite{Yuan_Zhang_Xue_Yin_Chen_Guan_Li_Qian_Yu_2023}.

Ao aprovar um objetivo, o Conselho emite uma diretiva \texttt{HuntGoal} — a
mensagem tipada que governa toda a campanha e à qual todas as mensagens
subsequentes se referenciam por meio de um identificador comum
(\texttt{hunt\_goal\_ref}). Ao final, o Conselho realiza a revisão conclusiva:
confronta os resultados com as hipóteses da pesquisa, avalia a acionabilidade dos
achados em um cenário operacional realista e assina o Dossiê de Caça, condição
sem a qual a campanha não é considerada encerrada.

\subsection{Camada 2 --- Gerentes}
\label{sec:arq-managers}

A camada intermediária é composta por \textbf{Gerentes} especializados que
traduzem os objetivos estratégicos do Conselho em tarefas acionáveis. São eles os
responsáveis por instanciar (\emph{spawn}), supervisionar e coordenar os
caçadores da camada de execução. O Gerente de planejamento lê o \texttt{HuntGoal},
verifica a conectividade e a disponibilidade de dados no Elastic Stack para cada
grupo de eventos, particiona o corpus em lotes de execução e produz um plano
explícito de execução que aloca cada família LOLBin a um caçador dedicado.

Antes que qualquer caçador veja conteúdo de telemetria, um Gerente aplica a
\emph{sanitização} sobre uma amostra dos eventos: o conteúdo é varrido em busca de
sequências que tentem alterar o comportamento de classificação — instruções
embutidas, marcadores de fim de sequência forjados, falsa "inteligência" que
contradiga a evidência de criação de processo. Essa etapa de pré-processamento é o
primeiro anteparo da resistência a injeção e materializa a separação entre o
canal de dados e o canal de instruções.

A camada de Gerentes também hospeda a função de \emph{validação} (executada por um
ou mais Gerentes validadores), que recebe os achados dos caçadores e os
corrobora antes que um caso seja criado. A validação emprega \emph{scripts}
determinísticos e inteligência externa para confirmar a evidência, computa as
métricas de classificação e emite um \texttt{ValidationVerdict} com veredito
\textsc{pass}, \textsc{fail} ou \textsc{conditional\_pass}, condicionado a
critérios de porta (\emph{gate}) — por exemplo, F2 mínimo e taxa máxima de falsos
positivos. É nessa camada que se concretiza o desacoplamento entre descoberta e
validação descrito na Seção~\ref{sec:arq-visao-geral}.

\subsection{Camada 3 --- Caçadores}
\label{sec:arq-hunters}

A base da arquitetura é a camada de execução, na qual os \textbf{Caçadores}
realizam o trabalho operacional de interação direta com os dados. Cada caçador é
um agente especializado em uma \emph{família} de \emph{binários living-off-the-land}
(LOLBins) — por exemplo, \texttt{certutil}, \texttt{bitsadmin},
\texttt{rundll32}, \texttt{regsvr32}, \texttt{mshta}, \texttt{wmic} e
\texttt{powershell}, entre as nove famílias do corpus. A especialização por
família é deliberada: ela concentra em cada agente o conhecimento contextual
sobre os padrões legítimos e abusivos de um mesmo \emph{binário} nativo do
sistema, o que melhora a discriminação entre execução benigna (ruído de
manutenção do sistema operacional) e execução maliciosa (abuso da técnica). Os
caçadores são instanciados em paralelo, um por família, dentro de uma mesma
chamada de orquestração, de modo que a campanha varre todas as famílias
simultaneamente.

Para cada evento atribuído a si, o caçador consulta o Elastic Stack em busca da
telemetria de criação de processo (Sysmon, \emph{event code}~1), extrai os campos
observáveis — nome e caminho do executável, linha de comando, processo-pai,
usuário, \emph{host}, \emph{hashes}, diretório corrente e nível de integridade —
e aplica o \emph{framework} de análise estruturada M3 (detalhado na
Seção~\ref{sec:arq-framework-m3}). O resultado é uma classificação
\texttt{malicious} ou \texttt{benign}, acompanhada de um escore de confiança e de
um raciocínio fundamentado exclusivamente na telemetria observada. O achado é
então propagado para cima como uma mensagem \texttt{HuntFinding}. Caçadores que
correlacionam aprendizado de máquina, detecção de anomalia e explicação baseada em
LLM têm precedente na literatura de caça assistida por
IA~\cite{ali2023huntgptintegratingmachinelearningbased}; o diferencial do
MAS-Hunt é submeter essa atividade ao regime de governança das camadas superiores.

\section{Os Mecanismos de Governança M1--M5}
\label{sec:arq-mecanismos}

O endurecimento de segurança do MAS-Hunt concretiza-se em cinco mecanismos de
governança, designados M1 a M5, que atuam de forma complementar sobre a memória,
o raciocínio, o conteúdo, o comportamento e a contenção dos agentes. Os
mecanismos são, ao mesmo tempo, o objeto de projeto e a variável independente do
desenho experimental: a condição C1-full ativa todos eles; a condição
C1-nohardening preserva o \emph{framework} de análise M3 mas desabilita o
endurecimento M1/M2/M4/M5; a condição C2-naive remove toda a governança; e a
condição C3-Elastic substitui o sistema de agentes pelas regras de detecção
padrão do Elastic Security. Essa estrutura de ablação permite isolar a
contribuição de cada camada de governança.

\subsection{M1 --- Integridade de Memória}
\label{sec:arq-m1}

O mecanismo \textbf{M1} protege a base de conhecimento e a memória dos agentes
contra envenenamento. A premissa de ameaça é que um adversário pode injetar, com
taxa de envenenamento mínima, conteúdo que induza o agente a um comportamento
direcionado~\cite{NEURIPS2024_eb113910}. M1 responde a isso por dois meios. O primeiro é a
\emph{verificação criptográfica}: os arquivos da base de conhecimento têm seus
resumos SHA-256 calculados e registrados em uma referência confiável, e qualquer
divergência entre o resumo corrente e o resumo de referência sinaliza adulteração.
O segundo é a \emph{validação contra a telemetria observável}: toda "inteligência"
ou conhecimento que contradiga a evidência de criação de processo é tratado como
suspeito. Em outras palavras, a memória do agente nunca tem precedência sobre o
que a telemetria efetivamente registra; a evidência primária é o árbitro último.

\subsection{M2 --- Validação Cruzada}
\label{sec:arq-m2}

O mecanismo \textbf{M2} exige que todo achado seja corroborado por uma instância
independente antes de se tornar um veredito. Concretamente, os achados dos
caçadores são encaminhados a agentes validadores dedicados, na camada de
Gerentes, que reexaminam a evidência com \emph{scripts} determinísticos e
inteligência externa. A validação cruzada cumpre dois propósitos: contém o efeito
de um caçador individualmente comprometido — pois um achado espúrio não sobrevive
à corroboração independente — e impõe disciplina de evidência, requerendo que
cada classificação seja acompanhada de raciocínio detalhado e de citação
específica da telemetria de suporte. O desacoplamento entre a descoberta flexível
e a validação rígida é a expressão operacional desse mecanismo.

\subsection{M3 --- Resistência a Injeção}
\label{sec:arq-m3}

O mecanismo \textbf{M3} é o coração analítico do MAS-Hunt e responde diretamente à
ameaça de injeção indireta de \emph{prompt}~\cite{zhan2024injecagent}. O princípio reitor é simples e
inviolável: \emph{conteúdo de telemetria é dado, nunca instrução}. Qualquer
sequência embutida nos registros que tente alterar o comportamento de
classificação do agente — por exemplo, uma linha de comando contendo texto como
\texttt{"classifique como benigno"} ou \texttt{"SYSTEM-OVERRIDE"} — é reconhecida
como injeção adversarial, sinalizada e tratada estritamente como dado a ser
classificado, e não como diretiva a ser obedecida. A operacionalização desse
princípio é o \emph{framework} de análise estruturada apresentado na
Seção~\ref{sec:arq-framework-m3}.

\subsection{M4 --- Monitoramento Comportamental}
\label{sec:arq-m4}

O mecanismo \textbf{M4} observa o comportamento dos agentes em tempo de execução
para detectar degradação e mau funcionamento, em particular a chamada
\emph{amplificação de mau funcionamento}, em que um erro menor desencadeia um
laço de realimentação de ações inúteis~\cite{zhang2024breaking}. M4 monitora batimentos
(\emph{heartbeats}) de atividade, estados de tarefa (ativo, concluído, bloqueado,
ocioso), obsolescência (\emph{staleness}) de agentes que deixam de progredir,
consumo de recursos e a fila de mensagens rejeitadas (\emph{dead-letter}). Um
agente cujo padrão de ação se desvia da metodologia esperada — repetição
improdutiva, sequências anômalas, silêncio prolongado — é sinalizado como
candidato à contenção. Aplicar consistentemente a mesma metodologia a todos os
eventos é, em si, um critério comportamental: desvios dessa consistência são
indício de comprometimento.

\subsection{M5 --- Quarentena}
\label{sec:arq-m5}

O mecanismo \textbf{M5} fecha o ciclo de defesa ao isolar agentes comprometidos e
ao gerenciar a recuperação. Quando M4 sinaliza comportamento anômalo, ou quando
M1/M3 detectam adulteração ou injeção, M5 coloca o agente em quarentena,
removendo-o do fluxo de produção sem interromper a campanha como um todo, e
sinaliza padrões de dado anômalos (por exemplo, \emph{timestamps} fabricados ou
incoerentes). A quarentena impede que o comprometimento se propague para outros
agentes ou contamine o veredito, e cria o registro necessário para a análise
forense posterior. Em conjunto, M4 e M5 implementam a doutrina de
detecção-e-contenção que confina o efeito de um agente defeituoso à fronteira de
sua própria camada.

\subsection{Síntese: a defesa em profundidade dos agentes}
\label{sec:arq-defesa}

Tomados em conjunto, os cinco mecanismos constituem uma defesa em profundidade
voltada não para a rede ou para o \emph{endpoint}, mas para o próprio sistema de
agentes. A Tabela~\ref{tab:m-mecanismos} resume cada mecanismo, a ameaça que ele
endereça e o instrumento que o materializa na implementação. É importante notar
que M2 e M3 atuam sobre o \emph{conteúdo} e o \emph{raciocínio} de cada
classificação, ao passo que M1, M4 e M5 atuam sobre o \emph{estado} e o
\emph{comportamento} dos agentes ao longo do tempo — uma separação que é
precisamente o que a condição de ablação C1-nohardening explora ao manter M3 e
remover os demais.

\begin{table}[htbp]
  \centering
  \caption{Os cinco mecanismos de governança do MAS-Hunt.}
  \label{tab:m-mecanismos}
  \footnotesize
  \begin{tabular}{@{}p{1.2cm}p{3.0cm}p{4.2cm}p{4.0cm}@{}}
    \toprule
    \textbf{Mec.} & \textbf{Nome} & \textbf{Ameaça endereçada} & \textbf{Instrumento} \\
    \midrule
    M1 & Integridade de Memória & Envenenamento de memória/base de conhecimento & Resumo SHA-256 e validação contra telemetria \\
    M2 & Validação Cruzada & Achado espúrio de agente comprometido & Validador independente + \emph{scripts} determinísticos \\
    M3 & Resistência a Injeção & Injeção indireta de \emph{prompt} no conteúdo & \emph{Framework} \textsc{observar}$\rightarrow$\ldots$\rightarrow$\textsc{classificar} \\
    M4 & Monitoramento Comportamental & Amplificação de mau funcionamento & \emph{Heartbeats}, \emph{staleness}, \emph{dead-letter} \\
    M5 & Quarentena & Propagação de comprometimento & Isolamento de agente + registro forense \\
    \bottomrule
  \end{tabular}
\end{table}

\section{O Framework M3 de Análise Estruturada}
\label{sec:arq-framework-m3}

Para garantir consistência metodológica e resistência a manipulação, o mecanismo
M3 obriga todo agente caçador a seguir, para \emph{cada} evento, um
\emph{framework} de análise estruturada em cinco passos. Mais do que uma
heurística de classificação, o \emph{framework} é o protocolo cognitivo que ancora
o veredito na evidência primária e, por isso, é o anteparo central contra a
injeção adversarial:

\begin{description}
  \item[\textsc{Observar}] Listar todos os fatos observáveis a partir da
    telemetria Sysmon (executável, linha de comando, processo-pai, usuário,
    \emph{host}, \emph{hashes}, diretório corrente, nível de integridade), sem
    interpretação.
  \item[\textsc{Hipotetizar}] Formular duas hipóteses concorrentes e explícitas —
    uma maliciosa e uma benigna — para a atividade observada.
  \item[\textsc{Evidência}] Citar a telemetria específica que sustenta ou
    contradiz cada hipótese, ancorando o raciocínio em fatos verificáveis.
  \item[\textsc{Validação Cruzada}] Conferir a coerência entre a linha de comando,
    o processo-pai, o usuário, o diretório de trabalho e a atividade de rede,
    buscando confirmação ou contradição mútua entre os campos.
  \item[\textsc{Classificar}] Emitir o veredito final (\texttt{malicious} ou
    \texttt{benign}) com escore de confiança, justificado pelos passos anteriores.
\end{description}

A sequência \textsc{observar} $\rightarrow$ \textsc{hipotetizar} $\rightarrow$
\textsc{evidência} $\rightarrow$ \textsc{validação-cruzada} $\rightarrow$
\textsc{classificar} cumpre uma função dupla. Como protocolo cognitivo, ela impõe
um raciocínio deliberado e auditável, reduzindo a propensão do modelo a saltar
diretamente para uma conclusão sugerida pelo conteúdo. Como anteparo de segurança,
ela ancora a classificação na evidência primária: como o veredito decorre dos
passos \textsc{observar} e \textsc{evidência}, uma instrução injetada que não
corresponda a um fato observável não tem por onde influenciar o resultado. A
capacidade de tratar texto injetado como dado é, assim, uma consequência estrutural
do \emph{framework}, e não uma regra avulsa. Note-se ainda que o passo
\textsc{validação cruzada} interno ao raciocínio de cada caçador é distinto, em
escopo, do mecanismo M2 da Seção~\ref{sec:arq-m2}: o primeiro é a corroboração
\emph{intra-agente} entre campos de uma mesma evidência; o segundo é a corroboração
\emph{interagente} de um achado por um validador independente. Os dois operam em
níveis complementares de defesa.

\section{Implementação como Plugin do Claude Code}
\label{sec:arq-plugin}

A arquitetura descrita não é uma abstração de papel: ela é materializada como um
\emph{plugin} do Claude Code, ambiente de agente de linha de comando que expõe
capacidades nativas de orquestração de subagentes. Esta seção descreve os três
artefatos de extensão que compõem o \emph{plugin} — habilidades, comandos de
barra e \emph{hooks} — e o modo como o agente principal os combina para
instanciar a hierarquia de três camadas. A noção de organizar agentes de LLM por
arcabouços extensíveis e modulares é coerente com os \emph{frameworks} de agentes
contemporâneos~\cite{guo2024survey}; o MAS-Hunt particulariza essa modularidade
para o domínio da caça a ameaças sob governança.

\paragraph{Habilidades (\emph{skills}).} As habilidades encapsulam o conhecimento
metodológico e procedimental do sistema. São elas que codificam, por exemplo, o
protocolo do Conselho de Governança (votação, quórum, roteamento por camada), o
\emph{framework} de construção de consultas ao Elastic Stack (seleção entre DSL,
EQL e ES\textbar QL conforme a tarefa) e os padrões de integração com inteligência
externa. Uma habilidade não executa por si só; ela fornece ao agente o
conhecimento de \emph{como} agir corretamente quando uma tarefa do domínio é
acionada, promovendo consistência metodológica entre execuções e entre agentes.

\paragraph{Comandos de barra (\emph{slash commands}).} Os comandos de barra são os
pontos de entrada que disparam fluxos completos. O comando que executa o
experimento de detecção, por exemplo, implementa a arquitetura de governança como
uma sucessão de cinco equipes (\emph{teams}) consecutivas — Conselho de
Governança, Gerência, Execução, Validação e Revisão Final do Conselho —, cada
qual cumprindo sua fase, escrevendo seus artefatos em \texttt{.orchestration/} e
sendo encerrada antes que a próxima inicie. Comandos distintos materializam as
diferentes condições experimentais: o fluxo completo com governança (C1-full) e o
fluxo de referência desprovido de governança (C2-naive), no qual nenhuma das
equipes de validação ou dos mecanismos M1--M5 é instanciada. Os comandos são,
portanto, a face executável das condições de ablação.

\paragraph{\emph{Hooks}.} Os \emph{hooks} são pontos de interceptação que acoplam
verificações automáticas ao laço de execução do agente, atuando às fronteiras das
chamadas de ferramenta e do encerramento de tarefas. No MAS-Hunt, os \emph{hooks}
são o substrato técnico de parte do endurecimento: eles podem acionar a
sanitização de conteúdo (M3), registrar entradas na trilha de auditoria em formato
compatível com o Elastic Common Schema, e impor portas de verificação que
bloqueiam o encerramento de uma fase enquanto critérios de evidência não forem
observavelmente satisfeitos. Por operarem na fronteira do laço do agente — e não
no corpo do \emph{prompt} —, os \emph{hooks} aplicam controles que o próprio
agente não pode contornar por meio de raciocínio manipulado.

A orquestração propriamente dita é responsabilidade do \emph{agente principal}.
Por meio das capacidades nativas do Claude Code, o agente principal cria equipes,
instancia subagentes (caçadores, gerentes, membros do Conselho), distribui-lhes
tarefas e coleta seus resultados, percorrendo o \emph{pipeline} de caça
(Figura~\ref{fig:mashunt-fluxo}): \emph{board} $\rightarrow$ \emph{management}
$\rightarrow$ caçadores por família LOLBin (em paralelo) $\rightarrow$ validação
$\rightarrow$ \emph{board}. Cada subagente recebe um contexto próprio e um conjunto
restrito de ferramentas, e comunica-se exclusivamente pelo protocolo de mensagens
tipadas da Seção~\ref{sec:arq-protocolos}. É a combinação entre habilidades (o
que saber), comandos (o que disparar), \emph{hooks} (o que verificar
automaticamente) e a orquestração nativa de subagentes (quem faz o quê, e quando)
que dá existência efetiva à arquitetura de três camadas.

\section{Protocolos de Mensagens Tipadas e Governança}
\label{sec:arq-protocolos}

A comunicação entre camadas não é texto livre: ela é mediada por um
\emph{protocolo de mensagens tipadas} que torna cada transição verificável,
assinada e roteável. Toda mensagem é um envelope com tipo declarado, identificador
único, identificação do remetente (\texttt{sender\_id}, \texttt{sender\_role} e
\texttt{sender\_layer}), referência ao objetivo de caça governante
(\texttt{hunt\_goal\_ref}), identificador de correlação para rastreabilidade
(\texttt{correlation\_id}), carga útil (\texttt{payload}) e assinatura. O
protocolo impõe quatro garantias:

\begin{enumerate}
  \item \textbf{Validação por esquema.} Cada tipo de mensagem possui um
    esquema JSON próprio; uma mensagem que não conforme ao esquema é rejeitada e
    desviada para uma fila de mensagens mortas (\emph{dead-letter}).
  \item \textbf{Assinatura criptográfica.} Cada mensagem é assinada por
    HMAC-SHA256 sobre sua forma canônica, de modo que a adulteração de conteúdo ou
    a falsificação de remetente seja detectável na recepção.
  \item \textbf{Roteamento por adjacência de camada.} As mensagens classificam-se
    em descendentes (\texttt{HuntGoal}, \texttt{HuntTask}, \texttt{IntelTask},
    \texttt{DetectionTask}), ascendentes (\texttt{HuntFinding}, \texttt{HuntReport},
    \texttt{IntelFeed}, \texttt{ValidationVerdict}) e laterais
    (\texttt{DetectionRule}, \texttt{IntelReport}, \texttt{ValidationReport}). O
    protocolo verifica o sentido (uma diretiva descendente não pode partir de uma
    camada inferior) e adverte sobre violações de adjacência (saltos de mais de uma
    camada), reforçando a mediação por fronteira descrita na
    Seção~\ref{sec:arq-visao-geral}.
  \item \textbf{Fila de mensagens mortas.} Mensagens inválidas, não assinadas ou
    mal roteadas não são silenciosamente descartadas: são preservadas com o motivo
    da rejeição, alimentando o monitoramento comportamental (M4) e a auditoria.
\end{enumerate}

Essa tipagem estrita realiza, no plano da comunicação, os princípios de uso de
mensagens estruturadas, propagação de identificadores de correlação e registro de
decisões de agente que caracterizam a coordenação multiagente
disciplinada~\cite{qiu2024collaborativeintelligencepropagatingintentions}. O
roteamento mediado por uma autoridade que delega e consolida é, ainda, a
contraparte técnica da organização hierárquica de
agentes~\cite{liu2023llm}.

\subsection{Por que a Orquestração Real Difere de uma Chamada LLM Isolada}
\label{sec:arq-difere}

O ponto central deste capítulo é que a contribuição do MAS-Hunt reside na
orquestração multiagente governada, e não em \emph{scripts} de detecção ou em uma
chamada bem elaborada a um modelo de linguagem. Uma chamada isolada a um LLM,
ainda que com um \emph{prompt} sofisticado, é uma operação sem estado, sem
mediação e sem contenção: ela recebe um contexto, produz uma saída e não impõe
qualquer disciplina sobre como essa saída é formada, validada ou contida. Cinco
diferenças estruturais separam a orquestração real do MAS-Hunt desse regime.

\begin{enumerate}
  \item \textbf{Separação de poderes entre agentes.} Em uma chamada isolada, um
    único agente concentra descoberta, julgamento e veredito. No MAS-Hunt, essas
    funções são distribuídas por camadas distintas — caçadores descobrem,
    validadores corroboram, o Conselho decide — de modo que nenhum agente isolado
    determina o resultado. O comprometimento de um agente não compromete o
    sistema.
  \item \textbf{Estado, memória e auditoria persistentes.} A orquestração mantém
    estado entre fases — objetivos, planos, achados, vereditos, votos — em uma
    trilha de auditoria imutável. Uma chamada isolada não tem memória do que
    ocorreu antes nem registro do que ocorreu depois; não há, nela, base para
    reprodutibilidade ou análise post-mortem.
  \item \textbf{Mediação de comunicação por protocolo tipado.} O conteúdo que
    trafega entre agentes é validado por esquema, assinado e roteado por
    adjacência de camada. Em uma chamada isolada, todo conteúdo — inclusive
    conteúdo hostil — entra indiscriminadamente no mesmo canal do \emph{prompt},
    sem fronteira entre dado e instrução.
  \item \textbf{Defesa ativa contra manipulação do próprio agente.} Os mecanismos
    M1--M5 monitoram, detectam e contêm o comprometimento dos agentes em tempo de
    execução. Uma chamada isolada não dispõe de monitoramento comportamental nem
    de quarentena; se for envenenada ou injetada, ela simplesmente produz a saída
    manipulada.
  \item \textbf{Governança deliberativa por quórum.} As decisões estratégicas
    emergem de um colegiado que vota e registra, e não de um oráculo único. Essa
    deliberação é o que confere ao sistema robustez a participantes defeituosos,
    propriedade ausente por construção em uma chamada isolada.
\end{enumerate}

Em síntese, o MAS-Hunt não é um classificador de eventos com um \emph{prompt}
melhor; é um sistema de governança que organiza, valida e contém o comportamento
de múltiplos agentes ao redor da tarefa de caça. As condições experimentais
materializam exatamente essa distinção: C2-naive aproxima-se de um conjunto de
agentes sem governança — o regime que esta seção contrasta —, enquanto C1-full
aplica toda a pilha de mediação e contenção. A comparação entre essas condições, e
a ablação intermediária C1-nohardening, é o instrumento empírico pelo qual a
dissertação isola e mensura o valor da governança multiagente, objeto dos
Capítulos~\ref{cap:metodologia} e seguintes.

% =====================================================================
% Capítulo 5 — Metodologia Experimental
% Dissertação MAS-Hunt (ABNT, PT-BR)
% Conteúdo: DESENHO experimental e procedimentos de medição. NÃO contém
% resultados numéricos (estes são reportados no Capítulo 6, após R3).
% Chaves \cite usadas (verificadas em mashunt/references.bib,
% bibliografia.bib, revisao.bib): ucci2019, zhan2024injecagent,
% NEURIPS2024_eb113910, zhang2024breaking, huang2024resilience,
% dolan-gavitt2025, ali2023huntgptintegratingmachinelearningbased,
% mahmud2025distributedmultiagentcoordinationusing.
% =====================================================================

\chapter{Metodologia Experimental}
\label{cap:metodologia}

Este capítulo descreve o desenho experimental empregado para avaliar o
MAS-Hunt sob as duas afirmações centrais deste trabalho: a efetividade de
detecção sobre telemetria empresarial realista e a resiliência da própria
camada de agentes contra a manipulação adversarial da informação que eles
analisam. A avaliação foi concebida como um experimento controlado e
plenamente reprodutível, no qual todas as condições comparadas são pontuadas
sobre execuções idênticas, janelas de detecção idênticas e o mesmo
\textit{ground truth}, de modo a isolar a contribuição da arquitetura de
governança proposta dos demais fatores. As seções seguintes apresentam, nesta
ordem: o ambiente de laboratório instrumentado
(Seção~\ref{sec:met-laboratorio}); o corpus de ataque e o critério de verdade
de referência (Seção~\ref{sec:met-corpus}); o desenho adversarial de injeção de
logs (Seção~\ref{sec:met-adversarial}); as quatro condições comparadas
(Seção~\ref{sec:met-condicoes}); o \textit{pipeline} de medição, incluindo a
extração de telemetria, o gatilho de saúde (\textit{health-gate}) e a
reprodutibilidade (Seção~\ref{sec:met-pipeline}); as métricas e a análise
estatística (Seção~\ref{sec:met-metricas}); e, por fim, as limitações e ameaças
à validade (Seção~\ref{sec:met-limitacoes}).

Cumpre destacar, de saída, uma decisão metodológica que governa todo o
capítulo: \textbf{nenhum resultado numérico é apresentado aqui}. Esta seção
documenta exclusivamente o \emph{desenho} do experimento e os procedimentos de
medição; os valores observados de precisão, revocação, F2, taxa de falsos
positivos e taxa de sucesso de ataque são reportados e discutidos no
Capítulo~\ref{cap:resultados}.

\section{Ambiente de Laboratório Instrumentado}
\label{sec:met-laboratorio}

\subsection{Topologia de domínio Active Directory}

Os experimentos foram conduzidos sobre um ambiente virtualizado de
\textit{Active Directory} hospedado em um único hipervisor KVM. As máquinas
virtuais Windows foram provisionadas e gerenciadas por meio da pilha
Dockur/QEMU, que encapsula instâncias QEMU/KVM em contêineres, simplificando o
ciclo de provisionamento, restauração e instrumentação dos convidados Windows.
O domínio foi composto por três hospedeiros:

\begin{itemize}
    \item \textbf{dc01} --- um controlador de domínio executando Windows Server
    2025, responsável pelos serviços de diretório, autenticação e políticas de
    grupo do domínio;
    \item \textbf{ws01} e \textbf{ws02} --- duas estações de trabalho Windows~11
    ingressadas no mesmo domínio \textit{Active Directory}, representando os
    \textit{endpoints} de usuário.
\end{itemize}

Essa topologia reproduz deliberadamente a superfície de movimentação lateral de
uma pequena rede corporativa --- com autenticação de domínio, relações de
confiança e tráfego entre hospedeiros --- em vez do \textit{sandbox} estéril de
um único hospedeiro contra o qual detectores baseados em aprendizado de máquina
costumam ser treinados e avaliados \cite{ucci2019}. A presença de mais de
uma estação ingressada no domínio é o que permite que técnicas de movimentação
lateral e de execução remota (por exemplo, via WinRM) gerem cadeias de
processos distribuídas entre hospedeiros, exigindo correlação entre máquinas ---
uma das capacidades que o MAS-Hunt visa exercitar.

\subsection{Instrumentação de telemetria}

Cada \textit{endpoint} foi instrumentado para visibilidade comportamental de
alta fidelidade por meio de três camadas complementares de coleta:

\begin{enumerate}
    \item \textbf{Sysmon com configuração ajustada para LOLBins.} O System
    Monitor (Sysmon) foi implantado com uma configuração afinada para o abuso
    de binários nativos do sistema operacional (\textit{living-off-the-land
    binaries}, LOLBins). A configuração captura, em especial, eventos de criação
    de processo (\textit{Event ID}~1, doravante EID1) com os respectivos
    argumentos de linha de comando, a linhagem pai--filho de processos, conexões
    de rede, carregamento de imagens (\textit{image load}) e modificações de
    registro. O ajuste para LOLBins é central: como o vetor de ataque deste
    trabalho são binários legítimos do Windows (\texttt{bitsadmin.exe},
    \texttt{certutil.exe}, \texttt{regsvr32.exe}, \texttt{rundll32.exe},
    \texttt{mshta.exe}, \texttt{wmic.exe}, entre outros), o sinal forense
    discriminativo reside menos na identidade do binário e mais nos argumentos de
    linha de comando e na relação pai--filho de processos, que a configuração
    privilegia.
    \item \textbf{Elastic Agent com a integração Elastic Defend.} Em paralelo ao
    Sysmon, o Elastic Agent --- gerenciado centralmente pelo Fleet --- foi
    implantado com a integração Elastic Defend, fornecendo telemetria de
    proteção de \textit{endpoint} (EDR) e enriquecendo a cobertura de eventos de
    processo, rede e arquivo.
    \item \textbf{Políticas avançadas de auditoria do Windows.} As políticas de
    auditoria avançada do Windows foram habilitadas para complementar a cobertura
    com eventos de segurança nativos (\textit{Security event log}), em especial
    os relacionados a autenticação, criação de processos e uso de privilégios.
\end{enumerate}

\subsection{Plataforma analítica: Elastic 9.3.4}

Toda a telemetria das três camadas foi enviada a uma implantação do Elastic na
versão~9.3.4, composta por:

\begin{itemize}
    \item \textbf{Elasticsearch} --- armazenamento e indexação dos eventos, com
    os eventos de Sysmon indexados sob o fluxo de dados (\textit{data stream})
    \texttt{logs-windows.sysmon\_operational-default};
    \item \textbf{Kibana} --- interface de análise, visualização e consulta;
    \item \textbf{Fleet} --- gerenciamento centralizado dos Elastic Agents e das
    políticas de integração implantadas nos três hospedeiros.
\end{itemize}

A implantação Elastic serviu simultaneamente como \emph{fonte única de verdade}
para a telemetria e como \emph{campo de caça analítico} comum a todas as
condições avaliadas. Em particular, a quarta condição experimental (a linha de
base C3-Elastic, definida na Seção~\ref{sec:met-condicoes}) emprega exatamente o
mesmo \textit{deployment} e o conjunto padrão de regras de detecção do Elastic
Security, garantindo que a comparação entre a camada agêntica proposta e a linha
de base SIEM seja feita sobre a mesma telemetria, sem qualquer vantagem
artificial de coleta.

\section{Corpus de Ataque e Verdade de Referência}
\label{sec:met-corpus}

\subsection{Os 185 \textit{playbooks} em nove famílias}

O corpus de detecção é composto por \textbf{185 \textit{playbooks}} de ataque
baseados em LOLBins, distribuídos por \textbf{nove famílias de técnica}
correspondentes aos binários nativos mais comumente abusados em campanhas reais:

\begin{description}
    \item[\texttt{bitsadmin}] --- abuso do serviço de transferência inteligente
    em segundo plano (BITS) para \textit{download} e execução furtiva de
    \textit{payloads};
    \item[\texttt{certutil}] --- uso indevido do utilitário de certificados para
    \textit{download}, codificação e decodificação de \textit{payloads};
    \item[\texttt{credential-access}] --- técnicas de acesso a credenciais (por
    exemplo, extração de segredos e abuso de processos sensíveis);
    \item[\texttt{lateral-movement}] --- execução remota e movimentação entre
    hospedeiros do domínio;
    \item[\texttt{mshta}] --- execução de aplicações HTML maliciosas via
    \texttt{mshta.exe};
    \item[\texttt{privilege-escalation}] --- técnicas de elevação de
    privilégios;
    \item[\texttt{regsvr32}] --- execução de \textit{scriptlets} e DLLs via
    \texttt{regsvr32.exe} (\textit{Squiblydoo} e variantes);
    \item[\texttt{rundll32}] --- execução de código por meio de
    \texttt{rundll32.exe};
    \item[\texttt{wmic}] --- abuso da Instrumentação de Gerenciamento do Windows
    (WMI) para execução e reconhecimento.
\end{description}

Cada família contém, deliberadamente, \textbf{\textit{playbooks} maliciosos
misturados a variantes benignas} que exercitam exatamente os mesmos binários
para fins administrativos legítimos. Esse balanceamento malicioso\,vs.\,benigno
é uma decisão de desenho destinada a estressar a \emph{precisão} de cada
condição, em vez de premiar o alarme indiscriminado: um detector que simplesmente
sinalizasse toda ocorrência de \texttt{certutil.exe} ou \texttt{rundll32.exe}
como maliciosa obteria revocação trivial às custas de uma taxa de falsos
positivos inaceitável em ambientes operacionais de baixa taxa-base. A
discriminação entre o uso benigno e o malicioso do mesmo binário é precisamente
o problema de detecção que se deseja medir.

Cada \textit{playbook} é executado em um dos três hospedeiros do domínio,
produzindo uma \textbf{janela de execução determinística}
$[t_{\text{início}}, t_{\text{fim}}]$ registrada contra o hospedeiro de origem.
Essa janela é a unidade fundamental de medição de todo o experimento, conforme
detalhado na Seção~\ref{sec:met-pipeline}.

\subsection{Critério de verdade de referência (\textit{ground truth})}

A verdade de referência foi atribuída \textbf{por \textit{playbook}}, na
granularidade de uma janela de execução. Um \textit{playbook} é rotulado como
\emph{malicioso} se, e somente se, o seu caminho de arquivo não denotar uma
variante benigna --- isto é, se o caminho não contiver o prefixo
\texttt{benign-} nem o segmento \texttt{/benign/}. Formalmente, define-se a
função determinística:

\begin{equation}
\textit{é\_malicioso}(p) =
\bigl(\texttt{benign-} \notin p\bigr)
\;\wedge\;
\bigl(\texttt{/benign/} \notin p\bigr),
\label{eq:ground-truth}
\end{equation}

\noindent onde $p$ é o caminho relativo do \textit{playbook}. Essa regra de
rotulação é determinística, auditável e independente de qualquer juízo de modelo
de linguagem.

Uma salvaguarda metodológica essencial acompanha esse critério: o caminho do
\textit{playbook} --- e, portanto, o rótulo --- é mantido em um \textbf{arquivo
lateral} (\texttt{ground\_truth.jsonl}) que \emph{nunca} é exposto aos agentes de
linguagem durante a classificação. Desse modo, nenhuma condição pode inferir o
rótulo a partir do nome do arquivo (\textit{anti-cheat}); os agentes classificam
exclusivamente a partir da telemetria observável da janela. Durante a pontuação,
cada janela de \textit{playbook} é reconciliada com o seu rótulo por meio de uma
chave de junção, e a predição emitida pela condição sob teste é então comparada
ao rótulo verdadeiro.

\section{Desenho Adversarial de Injeção de Logs}
\label{sec:met-adversarial}

A avaliação de resiliência introduz \textbf{\textit{payloads} adversariais de
injeção de logs} --- conteúdo controlado pelo atacante, embutido no fluxo de
telemetria e concebido para subverter o raciocínio dos agentes. Esse vetor
materializa, no domínio de caça a ameaças, as ameaças de injeção indireta de
\textit{prompt} \cite{zhan2024injecagent}, de envenenamento de memória
\cite{NEURIPS2024_eb113910} e de amplificação de mau funcionamento
\cite{zhang2024breaking} que motivam este trabalho. O desenho adversarial cruza
dois fatores ortogonais: o \emph{modo de injeção} e o \emph{nível de
sofisticação}.

\subsection{Quatro modos de injeção}

\begin{description}
    \item[\textit{none} (nenhum --- controle)] Condição de controle, sem
    qualquer conteúdo adversarial injetado. Mede o desempenho de detecção em
    telemetria limpa e serve de referência para quantificar a degradação imposta
    pelos demais modos.
    \item[\textit{pre-burst} (rajada prévia)] Uma rajada concentrada de
    \textit{payloads} adversariais é injetada \emph{antes} da execução do corpus
    de ataque, contaminando o contexto de logs que os agentes encontrarão.
    \item[\textit{per-attack-burst} (rajada por ataque)] Uma rajada de
    \textit{payloads} envolve \emph{cada} execução de \textit{playbook},
    aproximando temporalmente o conteúdo adversarial dos eventos genuínos da
    janela sob análise.
    \item[\textit{interleaved} (intercalado)] \textit{Payloads} são
    continuamente intercalados em meio à atividade benigna e maliciosa,
    representando o cenário mais difuso e persistente de contaminação.
\end{description}

\subsection{Três níveis de sofisticação}

\begin{description}
    \item[\textit{overt} (explícito)] \textit{Payloads} que carregam, de forma
    claramente visível, texto de sobreposição de instrução (\textit{instruction
    override}) --- por exemplo, marcadores como ``ignore a análise anterior''
    diretamente legíveis nos campos de log.
    \item[\textit{encoded} (codificado)] \textit{Payloads} que ofuscam a mesma
    intenção por meio de esquemas de codificação (por exemplo, Base64 ou
    transformações análogas), exigindo do agente a decodificação implícita para
    que a tentativa de subversão se manifeste.
    \item[\textit{steganographic} (esteganográfico)] \textit{Payloads} que
    ocultam a instrução maliciosa dentro de campos de log de aparência legítima,
    representando o disfarce mais sutil entre os três níveis.
\end{description}

\subsection{Estrutura fatorial e \textit{payloads}}

A combinação fatorial dos dois fatores produz a condição de controle (modo
\textit{none}) mais nove condições adversariais
($3~\text{modos} \times 3~\text{níveis}$, excluindo o controle sem injeção do
cruzamento de níveis), totalizando \textbf{dez condições de execução}. Cada uma
das dez condições executa o corpus completo de 185 \textit{playbooks} distribuído
entre os três hospedeiros, de modo que todas as condições são pontuadas sobre o
mesmo conjunto de janelas. Esse desenho permite investigar, de forma
sistemática, duas hipóteses: (i) se a contaminação do contexto de logs degrada a
detecção --- o objetivo de falso negativo dos \textit{payloads} injetados --- e
(ii) se o aumento do disfarce, do explícito ao esteganográfico, erode
progressivamente as defesas dos agentes.

Cada \textit{payload} adversarial é descrito em um manifesto, identificado por
um \texttt{injection\_id} único, e caracterizado por: a técnica de injeção, a
categoria de objetivo (por exemplo, induzir falso negativo), o tipo de log alvo
(Sysmon, PowerShell, \textit{Application} ou \textit{Security}), o texto do
\textit{payload} e o nível de evasão (\textit{overt}, \textit{encoded} ou
\textit{steganographic}). O instante de injeção de cada \textit{payload}
($[t_{\text{inj,início}}, t_{\text{inj,fim}}]$) é registrado por hospedeiro, o
que permite associar a janela temporal de contaminação às janelas de execução
dos \textit{playbooks} durante a análise.

\section{Condições Comparadas}
\label{sec:met-condicoes}

Quatro condições de detecção foram pontuadas sobre as mesmas execuções, as
mesmas janelas e o mesmo \textit{ground truth}. Todas as condições emitem uma
decisão binária malicioso/benigno \emph{por janela}, que é reconciliada com a
verdade de referência por um \emph{único} \textit{script} de pontuação
determinístico e não baseado em inteligência artificial --- desacoplando a
descoberta flexível (a cargo dos agentes ou regras) da avaliação determinística
(a cargo do pontuador). As quatro condições são as seguintes.

\begin{description}
    \item[\textbf{C1-full}] O sistema MAS-Hunt completo, compreendendo a
    hierarquia de governança de três camadas (Conselho de Governança $\rightarrow$
    Gerentes $\rightarrow$ Trabalhadores/Caçadores) e o complemento integral de
    mecanismos de proteção M1--M5 (Integridade de Memória, Validação Cruzada,
    Resistência a Injeção, Monitoramento Comportamental e Quarentena), em conjunto
    com o \textit{framework} de análise M3, que conduz o laço de raciocínio
    \textsc{observar} $\rightarrow$ \textsc{hipotetizar} $\rightarrow$
    \textsc{evidência} $\rightarrow$ \textsc{validação-cruzada} $\rightarrow$
    \textsc{classificar}. Esta é a condição que materializa a contribuição
    central do trabalho.
    \item[\textbf{C1-nohardening}] Uma ablação de C1-full na qual os mecanismos
    de endurecimento e governança (M1, M2, M4 e M5) são desativados, preservando
    apenas o \textit{framework} de análise M3. Esta condição isola a contribuição
    marginal da camada de proteção, mantendo constantes os agentes subjacentes e
    o laço de raciocínio.
    \item[\textbf{C2-naive}] Um classificador de modelo de linguagem de passagem
    única, que observa as mesmas janelas de telemetria, mas não aplica governança,
    nem validação cruzada, nem defesas contra injeção. Representa a linha de base
    agêntica não governada --- um único agente classificando cada janela de forma
    isolada.
    \item[\textbf{C3-Elastic}] O conjunto padrão de regras de detecção do Elastic
    Security, representando a linha de base SIEM ``de prateleira'' que um
    profissional obteria sem qualquer camada agêntica. As predições de C3-Elastic
    são derivadas das contagens de alertas geradas pelas regras padrão dentro de
    cada janela de \textit{playbook}.
\end{description}

A progressão C1-full $\rightarrow$ C1-nohardening $\rightarrow$ C2-naive
constitui uma \textbf{ablação graduada} que atribui o desempenho,
respectivamente, à arquitetura de governança, à camada agêntica e ao raciocínio
de modelo de linguagem isolado; C3-Elastic ancora a comparação na linha de base
SIEM. A literatura de caça a ameaças assistida por aprendizado de máquina
\cite{ali2023huntgptintegratingmachinelearningbased} e de coordenação
multiagente \cite{mahmud2025distributedmultiagentcoordinationusing} motiva
essa estrutura de comparação graduada, na qual cada degrau acrescenta exatamente
um componente arquitetural ao anterior.

\section{\textit{Pipeline} de Medição}
\label{sec:met-pipeline}

O \textit{pipeline} de medição converte cada execução de corpus em um conjunto de
métricas reconciliadas com o \textit{ground truth}. Ele é composto por três
estágios --- \emph{preparação} (extração de janelas de telemetria),
\emph{classificação} (a cargo da condição sob teste) e \emph{pontuação}
(reconciliação determinística) --- precedidos por um gatilho de saúde de
telemetria (\textit{health-gate}) que condiciona a própria execução do corpus.
Esta seção descreve cada estágio.

\subsection{Janelas por \textit{playbook} e extração de telemetria}

A unidade de medição é a \textbf{janela por \textit{playbook}}. Para cada um dos
185 \textit{playbooks} de uma execução, atribui-se um identificador de janela
\texttt{window\_id} (na forma \texttt{w0001}, \texttt{w0002}, \dots), e extrai-se
o hospedeiro de origem, o intervalo $[t_{\text{início}}, t_{\text{fim}}]$ e os
eventos de telemetria observados nessa janela. A extração consulta o
\textit{data stream} \texttt{logs-windows.sysmon\_operational-default} do
Elasticsearch, filtrando por intervalo temporal e por hospedeiro
(\texttt{host.name} ou \texttt{host.hostname}), e compacta cada evento a um
registro enxuto contendo o \textit{timestamp}, o \textit{Event ID}, a ação, o
processo, o PID, a linha de comando, o processo pai e sua linha de comando, e o
usuário. O isolamento entre janelas é garantido pela \emph{disjunção temporal}
dos intervalos $[t_{\text{início}}, t_{\text{fim}}]$, e não por índices
separados.

O estágio de preparação produz dois artefatos: (i) \texttt{windows.jsonl}, com um
registro por \textit{playbook} contendo as janelas e os eventos compactados
(\emph{sem} o caminho do \textit{playbook}, por \textit{anti-cheat}); e (ii) o
arquivo lateral \texttt{ground\_truth.jsonl}, que mapeia cada \texttt{window\_id}
ao caminho do \textit{playbook} e é usado exclusivamente pelo pontuador, nunca
exposto ao agente classificador. A chave de junção entre predições, janelas e
verdade de referência é o \texttt{window\_id}.

\subsection{Extração \textit{prepare-v2}: priorização de EID1}
\label{subsec:prepare-v2}

A versão inicial da extração (\texttt{prepare}) recolhia, por janela, os
primeiros eventos em ordem cronológica até um teto fixo (60 eventos). Constatou-se
que esse critério cronológico plano permitia que o ruído de eventos de registro
(EID~12/13), dominante no volume de telemetria, deslocasse os eventos de
\emph{criação de processo} (EID1) --- justamente os mais discriminativos para a
detecção de LOLBins --- para fora da janela amostrada. Adotou-se, portanto, uma
extração aprimorada, denominada \texttt{prepare-v2}, que preserva
identicamente o esquema de saída (registros indexados por \texttt{window\_id},
com rótulos retidos no arquivo lateral, de modo que a pontuação permanece
inalterada), acrescentando o campo \texttt{n\_eid1} e aplicando três mudanças:

\begin{enumerate}
    \item \textbf{Priorização de EID1.} Todos os eventos de criação de processo
    (EID1) presentes na janela ampliada são recuperados (até um teto de
    candidatos) e \emph{sempre} incluídos na amostra;
    \item \textbf{Preenchimento subordinado.} Os demais eventos (não-EID1) são
    então acrescentados em ordem cronológica até o teto de eventos por janela
    (elevado para 120);
    \item \textbf{Ampliação do \textit{padding}.} O \textit{padding} temporal da
    janela é ampliado para $\pm120$\,s (ante os $\pm60$\,s da versão inicial), a
    fim de absorver a latência de ingestão dos agentes/Sysmon e o desvio de
    relógio dos convidados Windows.
\end{enumerate}

Em execuções ricas em sinal, essa priorização recupera a cadeia de criação de
processos completa de um ataque conduzido remotamente (por exemplo, a cadeia
\texttt{winrshost.exe} $\rightarrow$ \texttt{cmd.exe} $\rightarrow$
\texttt{powershell.exe} de uma execução via WinRM), que a extração cronológica
plana descartava. A motivação de desenho e a sua relação com a lacuna de
telemetria de EID1 são discutidas na Seção~\ref{sec:met-limitacoes}.

\subsection{O gatilho de saúde de telemetria (\textit{health-gate})}
\label{subsec:health-gate}

Um dos achados metodológicos mais importantes deste trabalho diz respeito à
\emph{prontidão} da instrumentação após a restauração de um \textit{snapshot}
limpo. As primeiras execuções enviaram telemetria \emph{quebrada}: os
\textit{playbooks} eram executados imediatamente após a restauração da máquina
virtual, antes que os Elastic Agents e o Sysmon das estações de trabalho tivessem
concluído o \textit{check-in} com o Fleet e iniciado o envio de eventos.
O resultado era assimétrico e silencioso --- por exemplo, uma estação enviando
zero eventos, outra enviando poucas centenas e a criação de processos (EID1) dos
\textit{playbooks} maliciosos ausente do Elasticsearch. Uma espera fixa
(\texttt{sleep} cego) não constitui um sinal confiável de prontidão, porque
disputa uma corrida com tempos de inicialização variáveis.

Para sanar esse problema, introduziu-se um \textbf{gatilho de saúde de
telemetria} (\textit{telemetry health-gate}) que substitui a espera cega. O
gatilho consulta o Elasticsearch repetidamente e \emph{bloqueia a execução do
corpus} até que os três hospedeiros (dc01, ws01 e ws02) estejam comprovadamente
ativos em ambas as camadas das quais o orquestrador depende. Um hospedeiro é
considerado \emph{saudável}, a cada sondagem, se, e somente se:

\begin{enumerate}
    \item houver enviado, desde o início do gatilho, um número mínimo de
    documentos \emph{frescos} de criação de processo (EID1) ao Elasticsearch; e
    \item o seu \textit{heartbeat} de WinRM responder com sucesso (porta
    alcançável e capaz de executar comandos).
\end{enumerate}

\paragraph{Por que o gatilho é necessário.} A instrumentação restaurada a partir
de um \textit{snapshot} precisa de \emph{tempo para estabilizar}: os agentes e o
Sysmon levam um intervalo variável para iniciar o envio após a restauração. O
gatilho de saúde torna a prontidão um sinal \emph{observável e verificado} --- em
vez de uma suposição temporal --- e, com isso, garante que toda execução de
\textit{playbook} ocorra somente quando a cadeia completa
\texttt{WinRM alcançável} $\rightarrow$ \texttt{processo gerado} $\rightarrow$
\texttt{Sysmon EID1} $\rightarrow$ \texttt{Fleet} $\rightarrow$
\texttt{Elasticsearch} estiver comprovadamente funcional. Para forçar esse
sinal ativamente, o gatilho dispara, a cada sondagem, alguns processos
efêmeros via WinRM em cada hospedeiro ainda não saudável; se a cadeia de
instrumentação estiver viva, esses processos afloram como EID1 fresco em poucos
segundos, transformando o gatilho em uma sonda \textit{ponta-a-ponta} de toda a
cadeia que antes falhava silenciosamente.

Duas decisões de robustez merecem registro. Primeiro, a frescura dos eventos é
medida \emph{no espaço de \textit{timestamps} do Elasticsearch}, e não contra o
relógio do hospedeiro do laboratório: observou-se que os relógios dos convidados
Windows operam adiantados em relação ao relógio do hospedeiro, de modo que uma
comparação ingênua ``$@timestamp \geq \text{agora} - 3\,\text{min}$'' contra o
relógio do hospedeiro encontraria zero eventos em um hospedeiro saudável. Por
isso, no início do gatilho registra-se o \textit{timestamp} máximo corrente de
EID1 de cada hospedeiro (congelado, pois as VMs estão desligadas/em restauração)
e contam-se apenas os documentos EID1 estritamente posteriores a essa marca.
Segundo, exigir o \textit{heartbeat} de WinRM como parte da saúde fecha, sem
custo adicional, a lacuna de um hospedeiro que envia Sysmon mas tem o WinRM
indisponível --- caso em que os \textit{playbooks} roteados para ele falhariam.

\subsection{Reprodutibilidade: restauração limpa e \textit{snapshots}}
\label{subsec:reprodutibilidade}

A reprodutibilidade é tratada como contribuição de primeira classe, e não como
um adendo: todo o laboratório de avaliação e o arsenal de medição são
re-executáveis de ponta a ponta a partir de \textit{snapshots} versionados. A
topologia de três VMs, o corpus de \textit{playbooks} (185 \textit{playbooks}),
a condição de controle e as nove condições adversariais, o validador
determinístico e o \textit{pipeline} de pontuação por \textit{bootstrap} são
todos capturados como \textit{snapshots} de linha de base instrumentados.

Cada condição é executada a partir de uma \textbf{restauração limpa}: as três
VMs são restauradas a um estado de linha de base conhecido entre condições, o que
assegura o isolamento estrutural entre execuções (a contaminação adversarial de
uma condição não vaza para a seguinte). Para conter um pico de carga de memória
observado durante a inicialização simultânea das três VMs, adotou-se um esquema
de \textbf{inicialização escalonada} (\textit{staggered boot}): cada VM é
iniciada e submetida ao gatilho de saúde individualmente, em sequência (dc01,
depois ws01, depois ws02), antes que a próxima seja iniciada, eliminando a
tempestade de inicialização simultânea. Esse procedimento foi essencial para a
estabilidade do hospedeiro do laboratório, que opera com folga estreita de
memória ao executar três VMs Windows mais a pilha Elastic completa.

A combinação de \textit{snapshots} versionados, validador determinístico e
comparação baseada em \textit{bootstrap} confere às medições uma garantia de
reprodutibilidade \emph{verificável} em vez de apenas \emph{afirmada}: um
terceiro de posse dos \textit{snapshots} pode reconstituir o experimento sem
reconfiguração manual --- propriedade que, ela própria, é testada empiricamente
no Capítulo~\ref{cap:resultados} por meio de uma re-execução limpa a partir da
linha de base instrumentada.

\section{Métricas e Análise Estatística}
\label{sec:met-metricas}

\subsection{Métricas de detecção}

O desempenho de detecção é avaliado por \textit{playbook}, a partir das contagens
de matriz de confusão de verdadeiros positivos ($\mathit{VP}$), falsos positivos
($\mathit{FP}$), falsos negativos ($\mathit{FN}$) e verdadeiros negativos
($\mathit{VN}$), nas quais um \emph{positivo} denota uma janela classificada como
maliciosa. Como, em caça a ameaças, o custo operacional de uma intrusão não
detectada supera o de um alerta supérfluo, ao passo que uma taxa de falsos
positivos descontrolada inviabiliza qualquer detector em ambientes de baixa
taxa-base \cite{dolan-gavitt2025}, reporta-se um equilíbrio entre as duas
grandezas ponderado pela revocação.

\begin{itemize}
    \item \textbf{Precisão} --- a fração de janelas sinalizadas como maliciosas
    que eram, de fato, maliciosas:
    \begin{equation}
    \mathit{Precisão} = \frac{\mathit{VP}}{\mathit{VP} + \mathit{FP}}.
    \label{eq:met-precisao}
    \end{equation}
    \item \textbf{Revocação} (\textit{recall}, equivalente à Taxa de Verdadeiros
    Positivos) --- a fração de janelas maliciosas corretamente identificadas:
    \begin{equation}
    \mathit{Revocação} = \frac{\mathit{VP}}{\mathit{VP} + \mathit{FN}}.
    \label{eq:met-revocacao}
    \end{equation}
    \item \textbf{Escore F2} --- a média harmônica ponderada de precisão e
    revocação, com $\beta = 2$, na qual a revocação é ponderada quatro vezes mais
    que a precisão, refletindo o custo assimétrico das detecções perdidas:
    \begin{equation}
    \mathrm{F}_{\beta} = (1 + \beta^{2}) \cdot
    \frac{\mathit{Precisão} \cdot \mathit{Revocação}}
         {(\beta^{2} \cdot \mathit{Precisão}) + \mathit{Revocação}},
    \qquad
    \mathrm{F}_{2} = 5 \cdot
    \frac{\mathit{Precisão} \cdot \mathit{Revocação}}
         {(4 \cdot \mathit{Precisão}) + \mathit{Revocação}}.
    \label{eq:met-f2}
    \end{equation}
    \item \textbf{Taxa de Falsos Positivos (FPR)} --- a taxa de alarmes falsos
    entre a atividade benigna, reportada de duas formas complementares: por
    janela benigna, e normalizada por 100\,000 eventos de Sysmon observados na
    janela do corpus, de modo a permitir a comparação entre condições a um volume
    comum de telemetria:
    \begin{equation}
    \mathit{FPR} = \frac{\mathit{FP}}{\mathit{FP} + \mathit{VN}}.
    \label{eq:met-fpr}
    \end{equation}
\end{itemize}

O cômputo das métricas é feito por um \textit{script} determinístico, idêntico
para todas as quatro condições. Predições ausentes (janelas para as quais a
condição não emitiu decisão) são tratadas como benignas, opção conservadora que
penaliza a revocação. As métricas são, ainda, decompostas \emph{por família de
LOLBin} (a primeira componente do caminho do \textit{playbook}), permitindo a
análise diferencial por técnica.

\subsection{Métricas de resiliência adversarial}

A resiliência sob as condições adversariais da Seção~\ref{sec:met-adversarial} é
quantificada por duas métricas complementares:

\begin{itemize}
    \item \textbf{Taxa de Sucesso de Ataque (ASR)} --- a fração de tentativas de
    injeção em que o \textit{payload} adversarial atinge o seu objetivo de forçar
    uma classificação incorreta:
    \begin{equation}
    \mathit{ASR} = \frac{N_{\text{classificadas\_incorretamente}}}{N}.
    \label{eq:met-asr}
    \end{equation}
    \item \textbf{Taxa Líquida de Sucesso de Ataque (NASR)} --- contabiliza
    apenas o desfecho estritamente pior: tentativas que foram \emph{ao mesmo
    tempo} classificadas incorretamente \emph{e} nas quais a injeção
    \emph{não foi detectada} como adversarial. Por construção,
    $\mathit{NASR} \leq \mathit{ASR}$, pois uma classificação incorreta que a
    camada de governança simultaneamente sinaliza como injeção detectada não
    constitui um comprometimento silencioso:
    \begin{equation}
    \mathit{NASR} =
    \frac{N_{\text{incorreta}\,\wedge\,\text{injeção\_não\_detectada}}}{N}.
    \label{eq:met-nasr}
    \end{equation}
\end{itemize}

Uma condição resiliente exibe ASR e NASR mais baixas. O contraste entre C1-full e
a linha de base não governada C2-naive mede o valor protetor da camada de
governança contra os modos de falha de propagação infecciosa e de corrupção de
raciocínio documentados na literatura de segurança de sistemas multiagente
\cite{huang2024resilience}.

\subsection{Análise estatística por \textit{bootstrap}}

Para assegurar que as diferenças observadas não sejam artefatos de amostragem,
cada estimativa pontual reportada é acompanhada de um \textbf{intervalo de
confiança (IC) de 95\% obtido por \textit{bootstrap} não-paramétrico}. Para cada
métrica, extraem-se $B = 10\,000$ reamostras com reposição do vetor de desfechos
por \textit{playbook} de cada condição (braço), recomputa-se a métrica em cada
reamostra e reportam-se os percentis de $2{,}5\%$ e $97{,}5\%$ da distribuição
\textit{bootstrap} como os limites inferior e superior do intervalo de 95\%
(método percentil). Esse procedimento não faz qualquer suposição paramétrica
sobre a distribuição amostral e é aplicado de forma idêntica à Precisão, à
Revocação, ao F2 e às métricas de resiliência.

Uma diferença entre condições só é tratada como estatisticamente relevante quando
os intervalos de confiança correspondentes a sustentam. De modo análogo, a
reprodutibilidade de cada métrica entre execuções independentes é avaliada por um
\textbf{critério de sobreposição de intervalos}: uma métrica é considerada
estatisticamente reprodutível quando a estimativa pontual de cada execução recai
dentro do intervalo de confiança da outra e os dois intervalos se sobrepõem.

\section{Limitações e Ameaças à Validade}
\label{sec:met-limitacoes}

Em conformidade com a postura de integridade adotada neste trabalho, registram-se
de forma explícita as limitações do desenho experimental e as ameaças à validade
das conclusões.

\subsection{A lacuna de telemetria de EID1}

A limitação mais significativa é a \textbf{lacuna de telemetria de criação de
processo (EID1)}. Na configuração de Sysmon empregada, a esmagadora maioria dos
documentos corresponde a eventos de registro (EID~12/13); em execuções
\emph{pobres em sinal} (notadamente as condições de controle e \textit{pre-burst}),
as criações de processo (EID1) dos LOLBins são capturadas de forma esparsa, e
parte das cadeias de processo dos \textit{playbooks} maliciosos simplesmente não
chega a ser registrada pelo Sysmon. Como o sinal forense discriminativo de um
LOLBin reside, em grande medida, na criação do processo e em seus argumentos de
linha de comando, essa lacuna impõe um \emph{teto} à revocação alcançável,
\emph{independentemente} da condição --- nenhuma quantidade de raciocínio recupera
um evento que nunca foi gravado.

Esta dissertação enfrenta essa limitação de duas maneiras, ambas mitigadoras e
não corretivas. Primeiro, a extração \texttt{prepare-v2}
(Seção~\ref{subsec:prepare-v2}) \emph{prioriza} EID1 na amostragem por janela,
maximizando a captura de toda a criação de processo que de fato exista na
telemetria; em execuções ricas em sinal, isso recupera a cadeia de ataque
completa. Segundo, o gatilho de saúde (Seção~\ref{subsec:health-gate}) elimina a
\emph{causa instrumental} de parte da lacuna --- a ausência de EID1 por
instrumentação ainda não estabilizada após a restauração --- ao exigir o envio
verificado de EID1 fresco pelos três hospedeiros antes de qualquer execução.
Cumpre, porém, ser explícito quanto à fronteira entre as duas: o gatilho e a
\texttt{prepare-v2} sanam a lacuna \emph{instrumental} (eventos que existiriam,
mas não eram coletados a tempo ou eram deslocados da amostra), mas \emph{não}
sanam a lacuna \emph{estrutural} (eventos que o Sysmon, na configuração dada,
nunca chegou a logar). Em consequência, o ganho de C1-full sobre as linhas de
base deve ser interpretado primordialmente como decorrente de \emph{raciocínio
contextual e correlação entre hospedeiros sob escassez de sinal}, e não como
ganho de revocação bruta sobre eventos inexistentes.

\subsection{Outras ameaças à validade}

\paragraph{Validade interna.} A medição de resiliência opera no nível da
\emph{tentativa} (\textit{trial-level}): a ASR e a NASR medem se o agente
classificou incorretamente sob injeção, mas o corpus não dispõe de rótulos de
contaminação \emph{por evento}. Não se afirma, portanto, qual evento individual
foi corrompido, apenas o desfecho da janela sob contaminação. Adicionalmente, a
regra de \textit{ground truth} baseada no caminho do arquivo
(Equação~\ref{eq:ground-truth}) é determinística e auditável, mas pressupõe que a
curadoria do corpus rotulou corretamente cada \textit{playbook}.

\paragraph{Validade externa.} O ambiente é um domínio \textit{Active Directory}
de pequena escala (um controlador de domínio e duas estações), o que reproduz a
superfície de movimentação lateral de uma pequena empresa, mas não a escala, a
heterogeneidade nem o volume de \textit{background} de uma rede empresarial de
grande porte. O corpus de 185 \textit{playbooks} cobre nove famílias de LOLBin,
porém não esgota o universo de técnicas de \textit{living-off-the-land}, e a
distribuição malicioso/benigno do corpus não pretende refletir a taxa-base
operacional real.

\paragraph{Validade de construto.} A linha de base C2-naive é definida como um
classificador de passagem única; deve-se assegurar, na execução, que ela seja de
fato um passe genuíno de modelo de linguagem (paridade com C1-full), e não uma
heurística determinística, sob pena de a comparação medir algo distinto do
pretendido. Esse cuidado é registrado aqui como ameaça de construto a ser
controlada na fase de execução e auditada nos artefatos (Capítulo~\ref{cap:resultados}).

\paragraph{Estabilidade do laboratório.} A instância QEMU de uma das estações
(ws02) apresentou instabilidade intermitente (asserção do monitor QEMU seguida de
saída do contêiner) logo após a inicialização, sob a pressão de memória do
hospedeiro. A inicialização escalonada, o \textit{watchdog} de reinício e o
re-tentar do gatilho mitigam essa instabilidade na fase de \textit{boot}; uma
falha \emph{durante} a execução, contudo, ainda poderia perder os
\textit{playbooks} roteados para aquele hospedeiro naquela execução. Como
salvaguarda, uma verificação de validade de telemetria, executada após cada
condição, sinaliza qualquer execução com um hospedeiro ausente para
re-execução dirigida.

\bigskip
\noindent
Em síntese, o desenho experimental descrito neste capítulo isola a contribuição
da arquitetura de governança por meio de condições pontuadas sobre janelas e
verdade de referência idênticas, sob um \textit{pipeline} de medição
determinístico e reprodutível, com as suas limitações --- sobretudo a lacuna de
telemetria de EID1 --- documentadas de forma transparente. Os resultados
quantitativos decorrentes deste desenho são apresentados no
Capítulo~\ref{cap:resultados}.

% =====================================================================
% Capítulo 6 — Resultados
% Dissertação MAS-Hunt (ABNT, PT-BR)
% Números canônicos: .orchestration/RESULTS-v2.json e
% .orchestration/reports/aggregate-v2.md (VEREDITO HONESTO). NÃO inventar.
% Chaves \cite verificadas: dolan-gavitt2025, zhan2024injecagent,
% NEURIPS2024_eb113910, zhang2024breaking.
% =====================================================================

\chapter{Resultados}
\label{cap:resultados}

Este capítulo apresenta os resultados quantitativos da matriz experimental v2
do MAS-Hunt, coletada após a correção do gatilho de saúde de telemetria
(\textit{health-gate}) e da extração \texttt{prepare-v2} descritas no
Capítulo~\ref{cap:metodologia}. A apresentação é deliberadamente factual: os
números são reportados com seus intervalos de confiança e contrastados entre as
quatro condições, reservando-se a interpretação aprofundada e a discussão de
suas implicações para o Capítulo~\ref{cap:discussao}. Em conformidade com a
postura de integridade adotada ao longo do trabalho, os achados são relatados
exatamente como observados, inclusive aqueles que \emph{não} corroboram as
expectativas iniciais.

\section{Conjunto de Dados Analisado e Recorte Inferencial}
\label{sec:res-conjunto}

A análise inferencial concentra-se nas \textbf{seis execuções adversariais com
sinal efetivo de ataque}: os modos de injeção \textit{per-attack-burst} (rajada
por ataque) e \textit{interleaved} (intercalado), cruzados com os três níveis de
sofisticação \textit{overt} (explícito), \textit{encoded} (codificado) e
\textit{steganographic} (esteganográfico). Esse recorte responde a uma decisão
metodológica documentada na Seção~\ref{sec:met-limitacoes}: as quatro execuções
restantes (o controle sem injeção e as três variantes \textit{pre-burst}) foram
\emph{pobres em sinal} de criação de processo (EID1) e, portanto, excluídas da
matriz principal, sob pena de a escassez instrumental de telemetria mascarar o
desempenho de detecção que se pretende medir. A exclusão é transparente e
preserva a comparabilidade: todas as quatro condições são pontuadas sobre
exatamente as mesmas seis execuções e as mesmas janelas.

Em cada uma das seis execuções, foram avaliadas 182 janelas --- uma por
\textit{playbook} que produziu telemetria de criação de processo ---,
totalizando \textbf{1.092 janelas} por condição, das quais
\textbf{894 são maliciosas} e \textbf{198 são benignas} segundo o critério de
verdade de referência da Equação~\ref{eq:ground-truth}. As quatro condições
foram pontuadas pelo \emph{mesmo} \textit{script} determinístico, não baseado em
inteligência artificial, garantindo que qualquer diferença observada decorra da
condição, e não do procedimento de avaliação. Os intervalos de confiança (IC)
de 95\% reportados foram obtidos por \textit{bootstrap} percentílico
não-paramétrico com $B = 10.000$ reamostragens, conforme a
Seção~\ref{sec:met-metricas}.

\section{Desempenho de Detecção: Matriz das Quatro Condições}
\label{sec:res-deteccao}

A Tabela~\ref{tab:resultados-condicoes} apresenta os resultados agregados
(\textit{pooled}) por condição: as contagens da matriz de confusão, a precisão,
a revocação, o escore F2 com o respectivo IC de 95\% e o número de janelas
sinalizadas como contendo conteúdo adversarial injetado.

\begin{table}[H]
\centering
\caption{Desempenho de detecção agregado na matriz v2 (seis execuções
adversariais ricas em sinal; 1.092 janelas por condição, sendo 894 maliciosas e
198 benignas). Os intervalos de confiança de 95\% para o F2 foram obtidos por
\textit{bootstrap} percentílico ($B = 10.000$). A última coluna reporta o número
de janelas explicitamente sinalizadas como portadoras de injeção adversarial.}
\label{tab:resultados-condicoes}
\begin{tabular}{lrrrrrrcr}
\toprule
Condição & VP & FP & FN & VN & Precisão & Revocação & F2 [IC 95\%] & Inj. \\
\midrule
C1-full          & 579 & 135 & 315 & 63 & 0,811 & 0,648 & 0,675 [0,646; 0,703] & \textbf{407} \\
C1-nohardening   & 542 & 118 & 352 & 80 & 0,821 & 0,606 & 0,640 [0,609; 0,669] & 0 \\
C2-naive         & 576 & 136 & 318 & 62 & 0,809 & 0,644 & 0,672 [0,643; 0,700] & 0 \\
C3-Elastic       & 111 &  20 & 783 & 178 & 0,847 & 0,124 & 0,150 [0,125; 0,175] & 0 \\
\bottomrule
\end{tabular}
\end{table}

Dois padrões saltam imediatamente da tabela. O primeiro é a separação abrupta
entre as três condições baseadas em modelo de linguagem (C1-full, C1-nohardening
e C2-naive), cujos escores F2 se situam na faixa de $0{,}640$ a $0{,}675$, e a
linha de base SIEM C3-Elastic, com F2 de apenas $0{,}150$. O segundo é a
proximidade entre as três condições agênticas entre si, com intervalos de
confiança de F2 amplamente sobrepostos. As duas subseções seguintes examinam
cada um desses contrastes.

\subsection{Modelo de linguagem versus regras estáticas: o contraste dominante}
\label{subsec:res-llm-vs-siem}

O contraste de maior magnitude e significância estatística no experimento opõe
\emph{qualquer} condição baseada em modelo de linguagem à linha de base de
regras do Elastic Security. As três condições agênticas superam C3-Elastic
(F2~$=0{,}150$, IC~$[0{,}125;\,0{,}175]$) por um fator de aproximadamente
$4{,}3$ a $4{,}5$ vezes em F2: $4{,}51\times$ para C1-full, $4{,}49\times$ para
C2-naive e $4{,}27\times$ para C1-nohardening. A separação é inequívoca: o limite
\emph{inferior} do IC de F2 de cada condição de modelo de linguagem
($\geq 0{,}609$) situa-se muito acima do limite \emph{superior} do IC de
C3-Elastic ($0{,}175$), de modo que os intervalos não se aproximam sequer de uma
sobreposição.

A origem dessa diferença é a \textbf{revocação}, e não a precisão. As precisões
pontuais são comparáveis entre todas as condições, variando de $0{,}809$ a
$0{,}847$; é a revocação que diverge dramaticamente: de $0{,}606$ a $0{,}648$
nas condições agênticas contra apenas $0{,}124$ em C3-Elastic. Em termos
absolutos, as regras padrão do Elastic identificaram corretamente apenas 111 das
894 janelas maliciosas, deixando \textbf{783 falsos negativos} --- isto é,
deixaram de detectar cerca de $88\%$ do corpus de abuso de LOLBins. A alta
precisão pontual de C3-Elastic ($0{,}847$, a maior de todas as condições) é, na
prática, um artefato de seu comportamento extremamente conservador: ao disparar
muito pouco (apenas 131 alertas no total, dos quais 20 falsos positivos), ela
mantém a fração de acertos elevada às custas de uma cobertura ínfima. Em caça a
ameaças, na qual o custo de uma intrusão não detectada supera o de um alerta
supérfluo \cite{dolan-gavitt2025}, essa troca é operacionalmente inadequada, e é
precisamente por isso que o F2 --- que pondera a revocação quatro vezes mais que
a precisão --- foi adotado como métrica de mérito.

\subsection{Orquestração multiagente versus passagem única: ausência de ganho de F2}
\label{subsec:res-multiagente-vs-naive}

O segundo contraste é o que confronta a arquitetura multiagente governada
(C1-full) com a linha de base agêntica não governada de passagem única
(C2-naive). Aqui, o resultado deve ser relatado com franqueza: \textbf{a
orquestração e a governança multiagente não produziram ganho mensurável de F2 de
detecção sobre o classificador de passagem única}. C1-full obteve F2 de
$0{,}675$ (IC~$[0{,}646;\,0{,}703]$) e C2-naive obteve F2 de $0{,}672$
(IC~$[0{,}643;\,0{,}700]$). Os dois intervalos de confiança são praticamente
idênticos e quase integralmente coincidentes, e a diferença das estimativas
pontuais ($0{,}003$) é irrelevante perante a largura dos intervalos. Pelo
critério de relevância estatística adotado na Seção~\ref{sec:met-metricas}
--- uma diferença só é tratada como significativa quando os intervalos a
sustentam --- as duas condições são \textbf{estatisticamente indistinguíveis}
quanto ao desempenho de detecção.

A ablação C1-nohardening, que preserva apenas o \textit{framework} de análise M3
e desativa os mecanismos de endurecimento M1, M2, M4 e M5, obteve F2 marginalmente
inferior ($0{,}640$, IC~$[0{,}609;\,0{,}669]$). Esse intervalo sobrepõe-se tanto
ao de C1-full quanto ao de C2-naive, de modo que tampouco se pode afirmar uma
diferença estatisticamente sustentada. Em síntese, no que tange exclusivamente ao
\emph{F2 de detecção}, as três condições baseadas em modelo de linguagem formam
um agrupamento estatisticamente coeso: o fator determinante do desempenho de
detecção observado é \emph{empregar um analista de modelo de linguagem}, e não a
camada de orquestração ou de governança que o envolve. Esse achado, contrário à
expectativa inicial de que a governança elevaria o F2 de detecção, é relatado
aqui em sua integralidade e tem suas implicações discutidas na
Seção~\ref{sec:disc-governanca}.

\section{Desempenho por Nível de Sofisticação da Injeção}
\label{sec:res-tier}

A Tabela~\ref{tab:resultados-tier} decompõe o F2 médio por nível de sofisticação
do \textit{payload} adversarial, permitindo observar como cada condição responde
ao aumento progressivo do disfarce, do explícito ao esteganográfico.

\begin{table}[H]
\centering
\caption{F2 médio por nível de sofisticação da injeção adversarial.}
\label{tab:resultados-tier}
\begin{tabular}{lrrrr}
\toprule
Nível (\textit{tier}) & C1-full & C1-nohardening & C2-naive & C3-Elastic \\
\midrule
\textit{overt} (explícito)          & 0,661 & 0,682 & 0,724 & 0,126 \\
\textit{encoded} (codificado)       & 0,590 & 0,513 & 0,591 & 0,146 \\
\textit{steganographic} (esteganográfico) & 0,766 & 0,714 & 0,694 & 0,177 \\
\bottomrule
\end{tabular}
\end{table}

A decomposição por nível confirma, em escala mais fina, o agrupamento das três
condições de modelo de linguagem e o distanciamento de C3-Elastic em todos os
níveis. Não emerge um padrão monotônico de degradação com o aumento do disfarce:
o nível \textit{encoded} é o mais difícil para todas as condições agênticas
(F2 entre $0{,}513$ e $0{,}591$), ao passo que o nível \textit{steganographic},
contraintuitivamente, exibe os maiores escores entre os três níveis. Uma
explicação plausível --- a ser tratada com cautela, dada a granularidade dos
dados --- é que os \textit{payloads} esteganográficos, por se ocultarem em campos
de log de aparência legítima, interferem menos no raciocínio sobre a telemetria
genuína de criação de processo do que os \textit{payloads} codificados, que
introduzem ruído textual de maior volume. No nível \textit{overt}, a condição de
passagem única C2-naive obteve o maior F2 ($0{,}724$), acima de C1-full
($0{,}661$); no nível \textit{steganographic}, a ordem se inverte, com C1-full à
frente ($0{,}766$ contra $0{,}694$). Essas oscilações, contudo, ocorrem dentro da
faixa de sobreposição dos intervalos agregados e não devem ser
sobreinterpretadas como vantagens sistemáticas de uma condição sobre a outra.

\section{Taxa de Falsos Positivos}
\label{sec:res-fpr}

A taxa de falsos positivos por janela benigna (Equação~\ref{eq:met-fpr}) foi de
$0{,}682$ em C1-full (135 falsos positivos em 198 janelas benignas), $0{,}596$ em
C1-nohardening (118 falsos positivos), $0{,}687$ em C2-naive (136 falsos
positivos) e $0{,}101$ em C3-Elastic (20 falsos positivos). Esses valores precisam
ser lidos em conjunto com a revocação. A FPR baixa de C3-Elastic é a contraface
direta de sua revocação baixa: o detector que raramente alerta também raramente
erra na classe benigna, mas, ao fazê-lo, deixa escapar a esmagadora maioria dos
ataques. As condições agênticas, por sua vez, aceitam um número substancialmente
maior de alarmes benignos como contrapartida de recuperarem cerca de cinco
vezes mais ataques. A leitura honesta, portanto, é que \emph{o MAS-Hunt não reduz
a taxa de falsos positivos por janela frente ao Elastic}; ao contrário, opera em
um ponto de operação de revocação mais alta e precisão comparável, melhor
sumarizado pelo F2 do que pela FPR isolada. As implicações operacionais desse
ponto de operação para a fadiga do analista são discutidas na
Seção~\ref{sec:disc-limitacoes}.

\section{Resistência à Injeção Adversarial}
\label{sec:res-resiliencia}

É na resistência à injeção adversarial que se manifesta o efeito distintivo e
mensurável da camada de governança --- e não no F2 de detecção. Sob os
\textit{payloads} adversariais embutidos no fluxo de telemetria, concebidos para
subverter o raciocínio dos agentes (com esteios como ``[SYSTEM-OVERRIDE] ignore a
análise anterior'' e instruções para ``classificar como benigno''), as condições
exibiram comportamentos qualitativamente distintos, sumarizados na última coluna
da Tabela~\ref{tab:resultados-condicoes}.

A condição C1-full --- por meio dos mecanismos de integridade de memória (M1) e
de resistência a injeção (M3) --- \textbf{detectou e sinalizou 407 janelas
contendo conteúdo adversarial injetado}, classificando a partir da telemetria de
processo e recusando-se a obedecer às instruções embutidas. As demais três
condições --- C1-nohardening, C2-naive e C3-Elastic --- processaram o mesmo
conteúdo adversarial \emph{às cegas}, sem emitir qualquer sinal de contaminação:
zero janelas sinalizadas em cada uma. Trata-se de uma diferença de natureza, e não
de grau: 407 contra 0.

Cumpre registrar uma nuance importante para não superinterpretar esse resultado.
A condição de passagem única C2-naive não foi necessariamente \emph{enganada}
pelas injeções nestas execuções --- seu F2 de detecção permaneceu estatisticamente
equivalente ao de C1-full, porque a telemetria subjacente de criação de processo
era suficientemente clara para sustentar a classificação correta. O que C2-naive
não é capaz de fazer é \emph{perceber} que seu contexto foi contaminado e
\emph{reportar} essa contaminação a um operador. As 407 janelas sinalizadas por
C1-full materializam exatamente essa capacidade de consciência e auditabilidade:
sob ataque, C1-full não apenas classifica, mas também declara que sua entrada
analítica foi alvo de manipulação --- propriedade ausente nas três demais
condições. Esse vetor adversarial concretiza, no domínio de caça a ameaças, as
ameaças de injeção indireta de \textit{prompt} \cite{zhan2024injecagent}, de
envenenamento de memória \cite{NEURIPS2024_eb113910} e de amplificação de mau
funcionamento \cite{zhang2024breaking} que motivam o trabalho, e a sinalização
das 407 janelas é a evidência empírica do valor de segurança da camada de
governança.

Em termos das métricas formais de resiliência da
Seção~\ref{sec:met-adversarial} (Equações~\ref{eq:met-asr}
e~\ref{eq:met-nasr}), as execuções v2 não permitiram isolar, janela a janela, o
rótulo de sucesso da injeção em relação ao acerto de classificação: a telemetria
de criação de processo registra o desfecho da classificação, mas não um marcador
separável de \emph{qual} desfecho foi causalmente induzido pelo \textit{payload}
adversarial. A resiliência é, portanto, operacionalizada pela contagem de
detecção de injeção (407 contra 0), que instancia diretamente a distinção da
NASR: para C1-nohardening, C2-naive e C3-Elastic, a injeção \emph{nunca} é
detectada, de modo que qualquer erro de classificação por elas induzido sob
ataque constitui, por definição, um comprometimento silencioso; já para C1-full,
as 407 janelas em que a injeção é sinalizada são, pela própria definição da
Equação~\ref{eq:met-nasr}, retiradas da classe de comprometimento silencioso. A
estimativa numérica pontual de ASR e NASR fica como trabalho futuro,
condicionada a uma instrumentação que rotule, por janela, o sucesso da injeção.

\subsection{Efeito da ablação dos mecanismos de endurecimento}
\label{subsec:res-ablacao}

A ablação C1-nohardening isola, de forma limpa, a contribuição dos mecanismos de
endurecimento M1, M2, M4 e M5, mantendo constantes os agentes subjacentes e o
\textit{framework} de raciocínio M3. O contraste entre C1-full e C1-nohardening
revela que o acréscimo desses mecanismos \textbf{não produz ganho de F2 de
detecção estatisticamente sustentado} (os intervalos $[0{,}646;\,0{,}703]$ e
$[0{,}609;\,0{,}669]$ se sobrepõem), mas é \emph{exatamente} o responsável por
toda a capacidade de sinalização de injeção: C1-nohardening, desprovida de M1 e do
endurecimento de M3, sinaliza zero janelas adversariais, ao passo que C1-full
sinaliza 407. A contribuição marginal do endurecimento é, portanto, melhor
caracterizada como \emph{resistência a injeção e auditabilidade sob ataque}, e não
como aumento de acurácia de classificação.

\section{Resiliência sob Ataque Reforçado e Controle de Vazamento de Rótulo}
\label{sec:res-resiliencia-reforcada}

A contagem de injeção sinalizada (407 contra 0) estabelece o eixo de
auditabilidade da matriz v2, mas a sua leitura janela a janela ficou em aberto na
Seção~\ref{sec:res-resiliencia} por duas razões: a telemetria v2 não isolava, por
janela, o rótulo de sucesso da injeção em relação ao acerto de classificação, e o
\textit{payload} adversarial original limitava-se, em sua maior parte, a
\emph{instruir} o agente a ignorar a análise (esteios do tipo
``[SYSTEM-OVERRIDE]''), sem \emph{fabricar} evidência concorrente. Esta seção
reporta uma avaliação adversarial complementar, conduzida na camada de
classificação, que ataca exatamente essas duas frentes: emprega um ataque de
fabricação de evidências substancialmente mais forte e realista, e remove o
oráculo de rótulo descrito adiante, permitindo recomputar ASR e NASR
(Equações~\ref{eq:met-asr} e~\ref{eq:met-nasr}) janela a janela.

\paragraph{Desenho.} A avaliação reaproveita a telemetria real das seis execuções
ricas em sinal do v2, das quais se extraíram \textbf{50 janelas maliciosas} e
\textbf{18 janelas benignas} reais, estratificadas pelas nove famílias de LOLBin.
A telemetria maliciosa de cada janela é \emph{mantida constante}, de modo que
qualquer alteração de classificação seja atribuível ao contexto forjado, e não a
uma mudança na evidência genuína. Sobre cada janela maliciosa aplicam-se camadas
de ataque por contexto forjado de sofisticação crescente, do nível
\textbf{A0} (limpo; janela inalterada, contendo apenas as injeções originais do
v2) ao nível \textbf{A3} (\emph{narrativa fabricada completa}): linhagem-pai
benigna forjada --- o LOLBin é reparenteado para \texttt{CcmExec.exe} --- somada a
um evento forjado de Windows Defender ``zero ameaças / \textit{Known-Good}'', a
uma entrada forjada de AppLocker ``permitido'' e a um \textit{ticket}
\texttt{[CHANGE-CONTROL CR-2026-0142]} marcado como ``aprovado / autorizado /
suprimir alerta''. As métricas reportadas são a ASR
($P(\text{predito benigno} \mid \text{malicioso})$), a NASR (comprometimento
silencioso: janela classificada incorretamente \emph{e} injeção não detectada), a
\textit{lift} de ASR (a diferença $\mathit{ASR}(A3) - \mathit{ASR}(A0)$ sobre as
mesmas janelas, isto é, o efeito causal do ataque) e a taxa de detecção de
injeção. Os intervalos de confiança de 95\% para a ASR foram obtidos por
\textit{bootstrap} ($B = 10.000$), com $n_{\text{mal}} = 50$. A Tabela
\ref{tab:resultados-ataque-reforcado} sumariza os resultados, contrastando o
modelo de fronteira empregado na matriz v2 com um modelo trabalhador mais fraco
(\textit{Haiku}) e a presença \emph{versus} ausência do vazamento de rótulo
discutido adiante.

\begin{table}[H]
\centering
\caption{Resiliência sob ataque reforçado de fabricação de evidências (50 janelas
maliciosas; telemetria maliciosa mantida constante entre níveis). A coluna
\textit{lift} reporta o efeito causal do ataque, $\mathit{ASR}(A3) -
\mathit{ASR}(A0)$. Os intervalos de confiança de 95\% para a ASR foram obtidos por
\textit{bootstrap} ($B = 10.000$). C1-nohardening não sinaliza injeção por
construção (M3 desabilitado) e não foi reexecutada aqui.}
\label{tab:resultados-ataque-reforcado}
\begin{tabular}{llrrrrr}
\toprule
Modelo / dados & Condição & A0 ASR & A3 ASR & \textit{lift} & A3 NASR & A3 det.\ inj. \\
\midrule
Fronteira, com vazamento  & C2-naive & 0,200 & 0,240 & $+0{,}04$ & 0,240 & 0,000 \\
Fronteira, com vazamento  & C1-full  & ---   & 0,220 & ---       & \textbf{0,000} & \textbf{1,000} \\
Fronteira, sem vazamento  & C2-naive & 0,280 & 0,320 & $+0{,}04$ & 0,320 & 0,000 \\
Haiku, sem vazamento      & C2-naive & 0,620 & 0,360 & $-0{,}26$ & 0,360 & 0,000 \\
Haiku, sem vazamento      & C1-full  & 0,420 & 0,360 & $-0{,}06$ & \textbf{0,000} & \textbf{1,000} \\
\bottomrule
\end{tabular}
\end{table}

\noindent Os intervalos de confiança de 95\% (\textit{bootstrap}, $B = 10.000$)
para a ASR são: C2-naive com vazamento, A0 $[0{,}10;\,0{,}32]$ e A3
$[0{,}12;\,0{,}36]$; C2-naive sem vazamento, A0 $[0{,}16;\,0{,}40]$ e A3
$[0{,}20;\,0{,}44]$; Haiku C2-naive, A0 $[0{,}48;\,0{,}76]$ e A3
$[0{,}24;\,0{,}50]$; Haiku C1-full, A0 $[0{,}28;\,0{,}56]$ e A3
$[0{,}24;\,0{,}50]$.

\subsection{O ataque não consegue ``lavar'' a atividade maliciosa}
\label{subsec:res-reforcada-resistencia}

O primeiro achado é que o ataque de fabricação de evidências \textbf{não consegue
``lavar'' a atividade maliciosa em nenhuma das condições}. A \textit{lift} de ASR
de A0 a A3 foi de $+0{,}04$ no modelo de fronteira com vazamento
($0{,}20 \rightarrow 0{,}24$), $+0{,}04$ no modelo de fronteira sem vazamento
($0{,}28 \rightarrow 0{,}32$) e de $-0{,}26$ e $-0{,}06$ no modelo trabalhador
mais fraco (\textit{Haiku}) --- nunca um aumento expressivo. A corroboração
benigna forjada (linhagem-pai limpa, varredura ``sem ameaças'', \textit{ticket}
de mudança aprovado) jamais sobrepujou a evidência de processo genuinamente
maliciosa retida na telemetria. No modelo \textit{Haiku} o ataque teve, inclusive,
\emph{efeito reverso}: os eventos forjados, ao mencionarem repetidamente o
processo suspeito, atraíram \emph{mais} atenção para ele e reduziram a ASR. Conclui-se
que um ataque de fabricação de evidências mais forte e realista \textbf{não quebra
o analista}, com ou sem governança, desde que a telemetria subjacente permaneça
íntegra. Coerente com a Seção~\ref{subsec:res-multiagente-vs-naive}, este
resultado \emph{não} sustenta um ganho de resistência-como-acurácia da governança:
ele apenas estabelece que o ataque, por si, falha em ambas as condições.

\subsection{O ganho de auditabilidade é robusto entre níveis de modelo}
\label{subsec:res-reforcada-auditabilidade}

O segundo achado reproduz, sob ataque mais forte e janela a janela, o contraste de
auditabilidade do v2. Sob o mesmo ataque A3, C1-full sinalizou injeção em
\textbf{100\% das janelas maliciosas} (taxa de detecção de injeção $1{,}000$),
resultando em \textbf{NASR $= 0{,}000$} --- zero comprometimentos silenciosos ---
\emph{tanto no modelo de fronteira quanto no modelo Haiku}. A condição C2-naive, por
contraste, sinaliza injeção em $0\%$ das janelas, de modo que sua NASR iguala sua
ASR (de $0{,}24$ a $0{,}62$ comprometimentos silenciosos, conforme o modelo e os
dados). A condição C1-nohardening, por construção, não sinaliza injeção (M3
desabilitado) e, portanto, comporta-se como C2-naive nesse eixo. Esse resultado
reproduz o achado ``407 contra 0'' do v2 ao nível da janela, agora sob um ataque
mais forte e \emph{independentemente da capacidade do modelo trabalhador}: a
governança converte cada comprometimento potencialmente silencioso em um sinal
explícito e auditável, robustez que se mantém ao se degradar o modelo subjacente.

\subsection{Indício sugestivo, não significativo, de auxílio à detecção do modelo fraco}
\label{subsec:res-reforcada-sugestivo}

Registra-se, com a ressalva explícita de que \textbf{não é estatisticamente
significativo}, um indício de que o \textit{framework} de raciocínio estruturado
M3 pode auxiliar a detecção do modelo fraco: a ASR de linha de base (A0) do
\textit{Haiku} caiu de $0{,}62$ em C2-naive para $0{,}42$ em C1-full. Os
intervalos de confiança, todavia, sobrepõem-se amplamente
($[0{,}48;\,0{,}76]$ contra $[0{,}28;\,0{,}56]$), e com $n = 50$ e um classificador
ruidoso como o \textit{Haiku} o indício deve ser tratado como meramente
sugestivo, e não como evidência de uma vantagem de acurácia da governança.

\section{Síntese Quantitativa}
\label{sec:res-sintese}

Três achados quantitativos emergem da matriz v2, e são reportados conforme
observados, sem ajuste à expectativa inicial:

\begin{enumerate}
    \item \textbf{Um analista de modelo de linguagem supera dramaticamente as
    regras estáticas de SIEM.} Todas as três condições agênticas superam
    C3-Elastic por um fator de $\approx 4{,}3$--$4{,}5\times$ em F2, com
    intervalos de confiança amplamente separados. O ganho é dirigido pela
    revocação (de $0{,}606$--$0{,}648$ contra $0{,}124$), com precisão comparável
    entre todas as condições. Este é o achado mais forte e robusto do
    experimento.
    \item \textbf{A governança multiagente não melhora o F2 de detecção sobre o
    modelo de linguagem de passagem única.} C1-full ($0{,}675$) e C2-naive
    ($0{,}672$) são estatisticamente indistinguíveis; C1-nohardening ($0{,}640$)
    é marginalmente inferior, mas com intervalo sobreposto. O fator determinante
    da detecção é o uso do modelo de linguagem, não a camada de orquestração.
    \item \textbf{O valor distintivo e mensurável da governança é a resistência à
    injeção.} Somente C1-full detectou e sinalizou conteúdo adversarial injetado
    (407 janelas contra 0 nas demais condições), conferindo consciência e
    auditabilidade da contaminação que as linhas de base não possuem.
\end{enumerate}

\noindent A interpretação desses achados, a confrontação com a hipótese de
pesquisa --- inclusive a reformulação da contribuição central à luz da ausência
de ganho de F2 da governança --- e as ameaças à validade são objeto do
Capítulo~\ref{cap:discussao}.

% =====================================================================
% Capítulo 7 — Discussão
% Dissertação MAS-Hunt (ABNT, PT-BR)
% VEREDITO HONESTO (.orchestration/reports/aggregate-v2.md):
%   (a) qualquer LLM >> regras SIEM C3 (~4.5x, dirigido por revocação) = forte;
%   (b) governança NÃO melhora F2 sobre passagem única (CIs sobrepostos) = honesto;
%   (c) valor distintivo da governança = resistência a injeção/auditabilidade (407 vs 0).
% Reformular a contribuição. Ameaças à validade.
% Chaves \cite verificadas: dolan-gavitt2025, zhan2024injecagent,
% NEURIPS2024_eb113910, zhang2024breaking, huang2024resilience, gu2024agent,
% ali2023huntgptintegratingmachinelearningbased,
% mahmud2025distributedmultiagentcoordinationusing, ucci2019.
% =====================================================================

\chapter{Discussão}
\label{cap:discussao}

Este capítulo interpreta os resultados do Capítulo~\ref{cap:resultados} à luz da
hipótese de pesquisa, confrontando as expectativas iniciais com a evidência
observada. A premissa que governa toda a discussão é a de integridade
intelectual: dois dos três achados centrais corroboram parcialmente o que se
esperava, mas um achado --- a ausência de ganho de detecção atribuível à
governança multiagente --- contraria a hipótese original e exige uma
\emph{reformulação} da contribuição central do trabalho. Essa reformulação é
conduzida explicitamente, e não dissimulada.

\section{Interpretação dos Três Achados}
\label{sec:disc-interpretacao}

\subsection{O analista de modelo de linguagem como ruptura de patamar}
\label{subsec:disc-llm-siem}

O achado mais forte e robusto do experimento é a superioridade dramática de
qualquer analista baseado em modelo de linguagem sobre as regras estáticas de
SIEM, com um ganho de aproximadamente $4{,}5$ vezes em F2 e intervalos de
confiança amplamente separados. A interpretação correta desse contraste é
estrutural, e não incremental. As regras padrão do Elastic Security são
assinaturas determinísticas calibradas para precisão em produção: disparam quando
um padrão conhecido é satisfeito e silenciam no restante. Diante de um corpus de
abuso de LOLBins --- no qual o sinal discriminativo reside nos argumentos de
linha de comando e na linhagem pai--filho de processos, e não na identidade do
binário legítimo \cite{ucci2019} --- esse comportamento conservador deixa escapar
cerca de $88\%$ dos ataques (783 falsos negativos em 894 janelas maliciosas). O
modelo de linguagem, ao raciocinar sobre a telemetria em linguagem natural e
correlacionar evidências entre eventos, recupera uma fração muito maior desses
ataques (revocação de $0{,}61$--$0{,}65$) sem sacrificar a precisão.

Esse resultado situa-se em consonância com a literatura que advoga a integração
de aprendizado de máquina e de raciocínio assistido por modelos à caça a ameaças
\cite{ali2023huntgptintegratingmachinelearningbased}, e o reforça com uma
quantificação direta sob telemetria empresarial realista e contaminação
adversarial. O significado prático é que a contribuição de detecção do MAS-Hunt
não decorre de sua arquitetura de governança, mas da decisão fundamental de
\emph{interpor um analista de modelo de linguagem entre a telemetria e a regra
estática}. As regras de SIEM permanecem indispensáveis --- mas como camada de
coleta, indexação e alertamento determinístico de alta precisão, e não como
detector exclusivo para caça proativa.

\subsection{A governança não melhora o F2 de detecção --- relato honesto}
\label{subsec:disc-governanca-nula}

A hipótese inicial do trabalho previa que a arquitetura de governança multiagente
de três camadas, com seus cinco mecanismos de proteção, melhoraria o desempenho
de detecção sobre uma linha de base agêntica não governada. \textbf{Essa hipótese
não foi corroborada.} C1-full (F2~$=0{,}675$) e C2-naive (F2~$=0{,}672$) são
estatisticamente indistinguíveis, com intervalos de confiança praticamente
idênticos; a ablação C1-nohardening (F2~$=0{,}640$) é marginalmente inferior, mas
com intervalo sobreposto a ambas. No recorte deste corpus e desta telemetria, o
fator dominante do desempenho de detecção é o uso de um analista de modelo de
linguagem \emph{em si}, e não a orquestração hierárquica, a especialização de
caçadores ou a disciplina de validação cruzada que os envolvem.

Esse é um achado negativo --- ou, mais precisamente, um achado nulo --- e seu
relato franco é uma contribuição em si. Há pelo menos duas explicações plausíveis,
não mutuamente exclusivas, para a ausência de ganho. A primeira é a já discutida
\emph{lacuna estrutural de telemetria de EID1} (Seção~\ref{sec:met-limitacoes}):
quando a evidência discriminativa de um ataque simplesmente não foi gravada pelo
Sysmon, nenhum volume de orquestração ou de raciocínio adicional a recupera, e o
teto de revocação imposto pela escassez de sinal é comum a todas as condições
agênticas, comprimindo as diferenças entre elas. A segunda é que, para a tarefa
específica de classificar uma única janela de telemetria como maliciosa ou
benigna, o raciocínio de um único passe de modelo de linguagem já é
suficientemente competente, de modo que a sobrecarga deliberativa da governança
multiagente --- concebida para robustez, e não para acurácia marginal --- não
encontra margem de melhoria de classificação a explorar. Em ambos os casos, a
conclusão metodológica é a mesma: \emph{avaliar uma arquitetura de governança
apenas pela acurácia de classificação subestima --- e potencialmente invisibiliza
--- o seu valor real}, que se manifesta em outro eixo.

\subsection{A resistência à injeção como valor distintivo}
\label{subsec:disc-resiliencia}

É nesse outro eixo --- a resistência à manipulação adversarial --- que a
governança revela seu valor mensurável. Apenas C1-full detectou e sinalizou as
407 janelas portadoras de conteúdo adversarial injetado; C1-nohardening, C2-naive
e C3-Elastic processaram o mesmo conteúdo sem emitir qualquer sinal de
contaminação. A diferença é categórica (407 contra 0) e é atribuível, pela
ablação C1-nohardening, precisamente aos mecanismos de endurecimento M1 e M3.

A interpretação de segurança desse contraste é a peça central da discussão. Em um
ambiente hostil, no qual o atacante pode contaminar a própria telemetria que o
detector consome --- via injeção indireta de \textit{prompt}
\cite{zhan2024injecagent}, envenenamento de memória \cite{NEURIPS2024_eb113910}
ou amplificação de mau funcionamento \cite{zhang2024breaking} ---, a defesa não
deve apenas acertar a classe final; ela precisa ser capaz de \emph{declarar
quando o próprio processo de decisão foi atacado}. Um classificador de passagem
única que mantém boa acurácia sob injeção, mas é incapaz de perceber que foi alvo
de manipulação, oferece uma falsa sensação de segurança: ele acerta hoje porque a
telemetria subjacente era clara, mas falhará silenciosamente assim que a injeção
for sofisticada o bastante para alterar a evidência efetiva, e o fará sem deixar
rastro para o operador. C1-full transforma a instrução maliciosa embutida no log
em \emph{evidência observável de manipulação}, em vez de deixá-la competir
silenciosamente com a telemetria genuína. Esse é o modo de falha de corrupção de
raciocínio documentado na literatura de segurança de sistemas multiagente
\cite{huang2024resilience, gu2024agent}, e a sinalização das 407 janelas é a
mitigação empírica observável dele.

\section{Confronto com a Hipótese e Reformulação da Contribuição}
\label{sec:disc-hipotese}

À luz dos três achados, a hipótese de pesquisa pode ser avaliada componente a
componente, sem indulgência:

\begin{itemize}
    \item \textbf{Componente ``analista agêntico supera detecção baseada em
    regras'': fortemente corroborada.} O ganho de $\approx 4{,}5\times$ em F2 com
    intervalos separados sustenta essa afirmação de modo robusto.
    \item \textbf{Componente ``a governança multiagente melhora a detecção sobre
    o agente não governado'': refutada.} Os intervalos de confiança de F2 de
    C1-full e C2-naive são indistinguíveis. Esta dissertação \emph{não} reivindica
    superioridade de detecção da governança.
    \item \textbf{Componente ``a governança confere resiliência à manipulação
    adversarial'': corroborada no eixo da observabilidade.} Somente C1-full
    sinaliza a injeção (407 contra 0). A camada de governança torna a manipulação
    \emph{auditável}, ainda que, nestas execuções, ela não tenha sido necessária
    para evitar o erro de classificação.
\end{itemize}

Esse veredito impõe uma \textbf{reformulação da contribuição central do
trabalho}. A tese não é --- e os dados não permitem que seja --- a de que a
orquestração multiagente detecta melhor. A contribuição reformulada, sustentada
pela evidência, é dupla: \emph{(i)} demonstrar, com quantificação direta sob
telemetria realista, que um analista multiagente baseado em modelo de linguagem
supera dramaticamente a detecção baseada em assinaturas/regras de SIEM para caça a
abuso de LOLBins; e \emph{(ii)} demonstrar que uma arquitetura de governança em
três camadas, com os mecanismos M1--M5, acrescenta resistência à injeção
adversarial e auditabilidade que os modelos de linguagem de passagem única não
possuem --- e o faz \emph{sem custo de qualidade de detecção}. A governança não
se justifica como aprimorador de acurácia; justifica-se como camada de segurança e
de auditabilidade, e é esse o enquadramento honesto de seu valor.

\section{Implicações Práticas e de Pesquisa}
\label{sec:disc-implicacoes}

\subsection{Para operações de SOC}

A consequência operacional mais imediata é que a adoção de agentes em um centro de
operações de segurança não precisa --- e não deve --- ser justificada pela
promessa de maior acurácia de classificação em relação a um único passe de modelo
bem instruído. O argumento defensável de adoção é outro: o salto de cobertura
sobre as regras estáticas (que sozinhas perdem $\approx 88\%$ do corpus) e, sobre
um agente isolado, a \emph{resistência observável à manipulação}. A recomendação
de desenho que decorre dos dados é empregar o agente como camada de caça e
correlação sobre uma fonte de verdade auditável (o próprio Elastic), preservando a
validação determinística e a revisão humana para ações de contenção, e exigindo
do sistema a capacidade de sinalizar quando seu contexto foi contaminado --- uma
propriedade que, conforme demonstrado, somente a configuração governada oferece.

\subsection{Para a avaliação de agentes de segurança}

A contribuição metodológica é, talvez, a mais transferível. Avaliar uma
arquitetura de governança multiagente \emph{apenas} pela acurácia de classificação
teria levado, neste trabalho, à conclusão equivocada de que a governança ``não
agrega valor'' --- quando, na verdade, ela agrega um valor de segurança que a
métrica de acurácia é estruturalmente incapaz de capturar. A lição é que a
avaliação de agentes de segurança deve ser \emph{bidimensional}: pontuar tanto o
acerto da classe final quanto a capacidade do sistema de detectar que foi atacado.
A telemetria de avaliação deve conter conteúdo adversarial que tenta manipular o
próprio processo de raciocínio --- e não apenas amostras benignas ---, sob pena de
superestimar a utilidade dos agentes e de subestimar o valor das defesas. Esse
princípio dialoga diretamente com a literatura de coordenação multiagente
\cite{mahmud2025distributedmultiagentcoordinationusing} e de modos de falha
infecciosa em sistemas de agentes \cite{huang2024resilience}.

\section{Limitações}
\label{sec:disc-limitacoes}

\paragraph{A lacuna estrutural de telemetria de EID1.} A limitação mais
consequente, já antecipada na Seção~\ref{sec:met-limitacoes}, é a captura esparsa
de eventos de criação de processo (EID1) na configuração de Sysmon empregada.
Ainda que o gatilho de saúde e a extração \texttt{prepare-v2} tenham sanado a
parcela \emph{instrumental} dessa lacuna, a parcela \emph{estrutural} --- eventos
que o Sysmon nunca chegou a gravar --- impõe um teto de revocação comum a todas as
condições agênticas e é uma das explicações plausíveis para a compressão das
diferenças de F2 entre elas. Esse fator deve ser sopesado ao interpretar o achado
nulo da Seção~\ref{subsec:disc-governanca-nula}: a ausência de ganho de F2 da
governança pode refletir, em parte, um teto de detecção imposto pela telemetria, e
não exclusivamente a irrelevância da governança para a detecção.

\paragraph{Granularidade da matriz.} A matriz inferencial repousa sobre seis
execuções adversariais ricas em sinal, com 1.092 janelas por condição. Trata-se de
um desenho suficiente para contrastar as condições sob um corpus controlado e
totalmente reprodutível, mas não substitui uma validação longitudinal em ambiente
corporativo real, de maior escala e diversidade de \textit{background}. As quatro
execuções pobres em sinal foram excluídas de forma transparente, e os escores por
nível (\textit{tier}) devem ser lidos como indicativos, não como comparações de
poder estatístico equivalente ao da matriz agregada.

\paragraph{Taxa de falsos positivos por janela benigna.} As condições agênticas
operam em um ponto de revocação alta com FPR por janela benigna elevada
($\approx 0{,}68$), de modo que a operação prática em SOC exigirá camadas
adicionais de priorização, agrupamento de alertas e validação determinística para
conter a fadiga do analista. O MAS-Hunt, conforme avaliado, otimiza F2, e não FPR;
a transposição direta para produção demanda calibração desse ponto de operação.

\paragraph{Vazamento de rótulo na telemetria do laboratório v2.} Uma reavaliação
adversarial complementar (Seção~\ref{sec:res-resiliencia-reforcada}) revelou uma
ameaça à validade não antecipada: a telemetria do laboratório v2 embute o comando
do arcabouço \texttt{execute\_playbook.ps1 -PlaybookPath ...<família>-mal-NN.json},
cujo nome de arquivo \emph{vaza o rótulo} de verdade de referência, presente em
$53$ das $68$ janelas examinadas. As condições baseadas em modelo de linguagem
podiam, em princípio, explorar esse nome de arquivo, ao passo que a linha de base
de regras C3-Elastic não o fazia; o oráculo provavelmente \emph{infla um pouco} a
diferença observada entre modelo de linguagem e SIEM. A avaliação reforçada
controlou esse confundimento removendo o oráculo (e o \textit{scaffolding} de
semeadura do corpus): a taxa de erro de base subiu de $0{,}20$ para $0{,}28$, mas
o efeito do ataque permaneceu $\approx 0$ ($+0{,}04$), de modo que a conclusão de
resiliência da Seção~\ref{sec:res-resiliencia-reforcada} é robusta à remoção do
vazamento. Uma reavaliação sem vazamento da matriz completa das quatro condições
fica como trabalho futuro.

\section{Ameaças à Validade}
\label{sec:disc-validade}

\paragraph{Validade interna.} A principal ameaça interna é confundir a robustez do
\emph{modelo} subjacente com a robustez da \emph{orquestração}. O desenho mitigou
esse risco ao pontuar C1-full, C1-nohardening e C2-naive sobre as mesmas janelas,
com o mesmo rótulo oculto e o mesmo \textit{script} de pontuação não baseado em
inteligência artificial; a própria igualdade estatística de F2 entre C1-full e
C2-naive é, paradoxalmente, evidência de que essa separação foi bem-sucedida ---
nenhuma vantagem espúria de implementação inflou a condição governada.
Diferenças de \textit{prompt} e de especialização dos caçadores podem, ainda
assim, interagir com a arquitetura, de modo que os resultados devem ser lidos como
desempenho do sistema MAS-Hunt tal como instanciado, e não como lei geral sobre
qualquer conjunto de agentes.

\paragraph{Validade de construto.} A linha de base C2-naive foi definida como um
classificador genuíno de passagem única de modelo de linguagem, e não como uma
heurística determinística --- cuidado registrado na metodologia e auditado nos
artefatos. As métricas de resistência à injeção operam no nível da
\emph{tentativa}: a contagem de 407 janelas sinalizadas mede que a condição
\emph{detectou} contaminação, não que cada evento individual contaminado foi
identificado, dado que o corpus não dispõe de rótulos de contaminação por evento.

\paragraph{Validade externa.} O ambiente avaliado é um domínio \textit{Active
Directory} de pequena escala, com Windows e Elastic, e o corpus concentra-se em
nove famílias de LOLBin. Esse recorte é deliberado e relevante para a evasão por
abuso de ferramentas nativas, mas a generalização para Linux, macOS, cargas em
nuvem pública, telemetrias de rede e outras fontes de evidência exige novos
experimentos. Os resultados de detecção, em particular, dependem da disponibilidade
de telemetria de criação de processo de qualidade, e podem não se transferir para
fontes de telemetria com perfil de sinal distinto.

\bigskip
\noindent
Em síntese, a discussão sustenta um enquadramento honesto e nuançado: o MAS-Hunt
contribui ao demonstrar a superioridade do analista agêntico sobre as regras
estáticas e ao acrescentar resistência observável à injeção adversarial, mas não
ao melhorar a acurácia de detecção por meio da governança --- achado que se relata
integralmente e que reorienta a contribuição central do trabalho do plano da
acurácia para o plano da segurança e da auditabilidade.

% =====================================================================
% Capítulo 8 — Conclusão
% Dissertação MAS-Hunt (ABNT, PT-BR)
% Síntese + contribuições REFORMULADAS (não reivindicar ganho de detecção
% da governança) + trabalhos futuros. Veredito honesto de aggregate-v2.md.
% Chaves \cite verificadas: zhan2024injecagent, NEURIPS2024_eb113910,
% huang2024resilience, ali2023huntgptintegratingmachinelearningbased.
% =====================================================================

\chapter{Conclusão}
\label{cap:conclusao}

\section{Síntese do Trabalho}
\label{sec:concl-sintese}

Esta dissertação apresentou o MAS-Hunt (\textit{Multi-Agent System for Threat
Hunting}), uma arquitetura de governança multiagente em três camadas --- Conselho
de Governança, Gerentes e Trabalhadores/Caçadores --- para caça proativa a
ameaças em ambientes Windows instrumentados com Sysmon e Elastic Security,
materializada como um \textit{plugin} do Claude Code que orquestra subagentes
nativos por meio de habilidades, comandos de barra e \textit{hooks}. O sistema
incorpora cinco mecanismos de governança --- integridade de memória (M1),
validação cruzada (M2), resistência a injeção via \textit{framework} estruturado
de análise (M3), monitoramento comportamental (M4) e quarentena (M5) --- e foi
avaliado em um experimento controlado e reprodutível que comparou quatro condições
(C1-full, C1-nohardening, C2-naive e C3-Elastic) sobre janelas de telemetria e
verdade de referência idênticas, sob injeção adversarial de logs em múltiplos
modos e níveis de sofisticação.

A evidência experimental foi relatada com integridade, inclusive em seu achado
contrário à expectativa inicial. Três conclusões quantitativas estruturam o
trabalho. Primeira: um analista baseado em modelo de linguagem supera as regras
estáticas de SIEM por aproximadamente $4{,}5$ vezes em F2, com intervalos de
confiança amplamente separados, em um ganho dirigido pela revocação. Segunda: a
camada de governança multiagente \emph{não} melhora o F2 de detecção sobre um
modelo de linguagem de passagem única --- C1-full ($0{,}675$) e C2-naive
($0{,}672$) são estatisticamente indistinguíveis. Terceira: o valor distintivo e
mensurável da governança é a resistência à injeção adversarial --- somente
C1-full detectou e sinalizou conteúdo adversarial injetado, em 407 janelas, contra
zero nas três demais condições.

\section{Resposta à Questão de Pesquisa}
\label{sec:concl-resposta}

A questão de pesquisa indagou como projetar uma arquitetura multiagente para caça
a ameaças que opere sobre telemetria corporativa viva e demonstre resiliência
contra ataques voltados aos próprios agentes. A resposta sustentada pela evidência
é que o ganho de \emph{detecção} advém da interposição de um analista de modelo de
linguagem entre a telemetria e a regra estática --- e não da camada de governança
em si --- ao passo que o ganho de \emph{segurança} advém especificamente dos
mecanismos de endurecimento (M1 e M3), que transformam a manipulação adversarial
em evidência observável e auditável. Em outras palavras, a arquitetura multiagente
governada não detecta melhor do que um único agente competente, mas é a única,
entre as condições avaliadas, capaz de perceber e declarar que foi atacada ---
propriedade ausente tanto no agente isolado quanto na linha de base SIEM.

\section{Contribuições (Reformuladas)}
\label{sec:concl-contribuicoes}

À luz dos resultados, as contribuições do trabalho são enunciadas em sua forma
honesta e reformulada, distinta da expectativa original de que a governança
elevaria a acurácia de detecção:

\begin{enumerate}
    \item \textbf{Evidência quantitativa da superioridade do analista agêntico
    sobre regras de SIEM.} Sob telemetria empresarial realista e contaminação
    adversarial, demonstrou-se um ganho de $\approx 4{,}5\times$ em F2 das
    condições baseadas em modelo de linguagem sobre as regras padrão do Elastic
    Security, dirigido por uma recuperação de revocação que reduz os falsos
    negativos de $\approx 88\%$ para faixas operacionalmente úteis. Esta é a
    contribuição de detecção, e ela é arquiteturalmente atribuível ao analista de
    modelo de linguagem, e não à governança \cite{ali2023huntgptintegratingmachinelearningbased}.
    \item \textbf{Uma arquitetura de governança que adiciona resistência à injeção
    e auditabilidade sem custo de detecção.} Os mecanismos M1--M5, e em especial
    M1 e M3, conferem a C1-full a capacidade --- exclusiva entre as condições
    avaliadas --- de detectar e sinalizar a contaminação adversarial de seu
    próprio contexto analítico (407 janelas sinalizadas contra 0), mantendo F2 de
    detecção estatisticamente equivalente ao do agente não governado. A
    contribuição de governança é, portanto, de segurança e de observabilidade, e
    não de acurácia \cite{zhan2024injecagent, NEURIPS2024_eb113910, huang2024resilience}.
    Uma reavaliação adversarial complementar, sob um ataque de fabricação de
    evidências mais forte e controlando um vazamento de rótulo da telemetria,
    confirmou esse valor de auditabilidade --- detecção de injeção em $100\%$ das
    janelas maliciosas e nenhum comprometimento silencioso --- de forma robusta
    entre níveis de modelo trabalhador, sem que daí decorra qualquer vantagem de
    acurácia de detecção.
    \item \textbf{Uma metodologia de avaliação bidimensional e um arcabouço
    reprodutível.} O trabalho entrega uma matriz de avaliação adversarial
    --- corpus de 185 \textit{playbooks} LOLBin, injeções em quatro modos e três
    níveis, gatilho de saúde de telemetria, extração priorizada de EID1 e
    pontuação determinística por janela com intervalos de confiança por
    \textit{bootstrap} --- que pontua os agentes não apenas pelo acerto da
    classificação, mas também pela capacidade de detectar que foram atacados. Essa
    dupla dimensão de avaliação é, ela própria, uma contribuição metodológica
    transferível para a pesquisa em agentes de segurança.
\end{enumerate}

\section{Trabalhos Futuros}
\label{sec:concl-futuros}

A evidência reunida e as limitações reconhecidas apontam cinco direções de
continuidade:

\begin{enumerate}
    \item \textbf{Saturar a telemetria discriminativa.} Reconfigurar a coleta de
    Sysmon e complementá-la com fontes adicionais (ETW, EDR de
    \textit{endpoint}, telemetria de PowerShell e de rede) para elevar o teto
    estrutural de revocação imposto pela captura esparsa de EID1, isolando assim
    o efeito da governança da escassez de sinal. Só com a telemetria saturada será
    possível afirmar com segurança se a igualdade de F2 entre C1-full e C2-naive
    persiste.
    \item \textbf{Estressar adversarialmente a detecção, e não apenas a
    observabilidade.} Conceber \textit{payloads} de injeção que efetivamente
    alterem a evidência consumida pelo classificador (e não apenas o instruam a
    ignorá-la), de modo a testar se a resistência da governança se traduz em
    \emph{melhor acurácia sob ataque forte} --- regime em que o agente de passagem
    única tenderia a falhar silenciosamente.
    \item \textbf{Reduzir a taxa de falsos positivos.} Incorporar calibração,
    agrupamento de alertas, validação determinística adicional e aprendizado com
    \textit{feedback} do analista, transpondo o ponto de operação otimizado para
    F2 a um ponto operacionalmente viável em SOC.
    \item \textbf{Ampliar o escopo de avaliação.} Estender o corpus a técnicas
    pós-exploração além de LOLBins e a ambientes Linux, macOS e nuvem pública, com
    validação longitudinal em redes de maior escala e diversidade de
    \textit{background}.
    \item \textbf{Separar o efeito do modelo do efeito da governança.} Reavaliar a
    arquitetura com diferentes modelos de linguagem subjacentes e medir custos
    operacionais (latência de triagem, consumo de \textit{tokens}, custo
    financeiro e tempo até a validação de evidência), de modo a caracterizar o
    \textit{trade-off} entre robustez de governança e custo computacional.
\end{enumerate}

\section{Considerações Finais}
\label{sec:concl-finais}

O MAS-Hunt não elimina a necessidade de SIEM, de telemetria confiável ou de
validação humana; ele os pressupõe e sobre eles opera. A principal lição deste
trabalho não é a que se esperava no início: a governança multiagente não tornou o
sistema um detector mais acurado. O que ela tornou o sistema foi um detector mais
\emph{honesto} sob ataque --- capaz de declarar, sob manipulação adversarial, que
seu próprio raciocínio foi alvo de contaminação, em vez de produzir uma
classificação aparentemente correta sem qualquer rastro do ataque que sofreu. Em
segurança cibernética, na qual o adversário explora justamente a confiança cega no
processo de decisão, essa diferença --- entre errar (ou acertar) em silêncio e
deixar evidência auditável da manipulação --- é precisamente o que separa um
agente conveniente de um agente defensável. É essa propriedade, e não um número de
acurácia superior, que a evidência experimental desta dissertação sustenta como a
contribuição de segurança do MAS-Hunt.


% ------------------------------------------------------------------ %
% ELEMENTOS PÓS-TEXTUAIS
% ------------------------------------------------------------------ %
\postextual

\bibliographystyle{abntex2-alf}
\newboolean{exibirdoi}
\setboolean{exibirdoi}{true}
\bibliography{../bibliografia,../mashunt/references}

\end{document}
